\chapter{El departamento antes del mercado}

\label{cap:siglo}
El Boletín Oficial de Salta empieza a publicarse el 15 de octubre de 1908, y su archivo
descargable arranca con la edición número uno. Las dos primeras décadas de esa colección
contienen unas cuarenta menciones al departamento La Caldera. Son pocas, están dispersas y
ninguna es espectacular. Leídas juntas describen un territorio cuyas relaciones con el
Estado ya estaban planteadas en los términos en que siguen planteadas hoy.

\section{Un municipio que vendía su recaudación}

En febrero de 1909, el Boletín publica el cuadro completo de entradas y salidas del
municipio de La Caldera durante el ejercicio anterior. Es la primera rendición de cuentas
conocida del departamento y cabe en una página.

\begin{ficha}{Inversión de las rentas municipales, año 1908}
Elevado el 15 de enero de 1909 por \textbf{Ángel Fernández\index{Fernández, Ángel}}, Presidente Municipal, al
Ministro de Gobierno, que ordena publicarlo.

\textbf{Entradas:} saldo del año anterior \$283,77; \textbf{producido de la venta por
licitación de los impuestos municipales} \$1.150; derechos de cementerio \$50.
\textbf{Total: \$1.483,77.}

\textbf{Principales salidas:} al Consejo General de Educación, el 20\% adicional, \$200;
secretario y comisario municipal, \$180; farolero municipal, \$54; \textbf{limpieza del
camino nacional desde el Angosto de Arias hasta Castañares y parte de la Calderilla,
\$293,60}; alumbrado, talonarios y escritorio, \$91,90; \textbf{construcción de un lazareto
con motivo de la peste de la viruela, \$63,10}; festejos del 9 de julio, \$80,50; comisión
de señoras para la función del Rosario, \$65; media docena de sillas para la comisión
municipal, \$30.

Existencia en caja al cierre: \$280,12.

\textit{El cuadro cierra: entradas y salidas dan \$1.483,77. Pero la partida educativa no:
el 20\,\% de los \$1.200 que el propio cuadro declara son \$240, y el cuadro imprime \$200.}
\end{ficha}

La partida mayor de ingresos —\$1.150 sobre \$1.483, casi el ochenta por ciento— es «el
producido de la venta por licitación de los impuestos municipales». El municipio no
recaudaba: remataba el derecho a recaudar, y cobraba por adelantado un precio fijo. Quien
ganaba la licitación se quedaba con lo que lograra cobrar.

Ese arreglo tiene una consecuencia que atraviesa el siglo. Un municipio que vende su
recaudación no tiene incentivo ni capacidad para conocer su base imponible: no necesita
saber cuántos inmuebles hay ni cuánto valen, porque su ingreso no depende de eso.

Ciento siete años después, el capítulo \ref{cap:hacienda} documenta lotes urbanos con valuación fiscal de
tres cifras y un municipio cuya superficie crece sin que crezca su recaudación. No es la
misma institución, pero es el mismo problema, y aquí está su forma más antigua.

\textbf{Y hay algo más que este cuadro, leído contra los años que siguen, obliga a decir: no fue el
primero de una serie, fue el único.} Los siete años que van de 1910 a 1916 se leyeron edición por
edición, y cada uno de esos barridos registra el mismo cero: \textbf{el Boletín no publica en esos
siete años un solo balance municipal}, ni de este departamento ni de ningún otro municipio de la
provincia. Sumados, son siete años consecutivos sin una sola rendición de cuentas municipal publicada.

\begin{cautela}
El cero es del \textit{Boletín}, no de los municipios: que no se publicara no prueba que no se
rindiera. Lo que prueba es que \textbf{entre 1909 y 1922 la rendición municipal dejó de ser un acto
público}, y que la de 1908 ---elevada por el presidente municipal al Ministro de Gobierno, que ordena
publicarla--- es un caso aislado y no el primer eslabón de una costumbre. Cuando en 1922 el
departamento vuelva a aparecer con números, ya no será rindiendo cuentas: será un presupuesto que la
comisión municipal sanciona y el Poder Ejecutivo aprueba.
\end{cautela}

Dos datos más de ese cuadro. El gasto mayor del año —un quinto del presupuesto— fue
limpiar el camino nacional entre el Angosto de Arias, Castañares y parte de La Calderilla:
en 2014, la Dirección de Vialidad contrataría al mismo municipio para desmalezar banquinas
y construir badenes en ocho rutas provinciales. Y el único gasto sanitario del ejercicio
fue construir un lazareto durante la peste de viruela.

\section{El año institucional de 1909}

El mismo año en que se publica ese cuadro, el Boletín deja ver cómo se armaba el gobierno del
departamento. \textbf{El 13 de enero de 1909 el Poder Ejecutivo nombró por decreto a los cinco
miembros de la Comisión Municipal de La Caldera} (B.O.\ Nº 26, 1909, h.\ 4) ---Ángel Fernández\index{Fernández, Ángel},
Silvano Murúa (hijo)\index{Murúa, Silvano Ignacio}, el doctor Rafael Usandivaras, Jacinto
Guerrero\index{Guerrero, Jacinto} y Belisario Benítez--- y el decreto dice por qué lo hacía él:
los nombramientos no podían hacerse «\textit{conforme lo prescribe el art.\ 173 de la
constitución}» porque la mayor parte de esas corporaciones no había hecho el sorteo de
renovación por mitades. La excepción funcionaba como regla.

El artículo 2º completa el circuito. Las comisiones nombradas debían recibir de las cesantes el
archivo «\textit{y demás pertenencias}» y elevar al Poder Ejecutivo \textbf{las ternas para el
nombramiento de los jueces de paz} y \textbf{el presupuesto de gastos del año}. El municipio
proponía; la Provincia designaba y aprobaba. Es, con noventa y nueve años de anticipación, la
misma relación que el capítulo \ref{cap:hacienda} describe para el presupuesto municipal de
1932 y para el de hoy.

\textbf{Y después está lo que pasó con esos cinco.} Tres renunciaron en el año, y los tres
reemplazos salieron del mismo grupo de nombres: en febrero, Silvano Murúa\index{Murúa, Silvano Ignacio}
es reemplazado por \textbf{Abelardo Figueroa}\index{Figueroa, Abelardo} (B.O.\ Nº 36, 1909, h.\ 3); en abril, Jacinto
Guerrero\index{Guerrero, Jacinto} por \textbf{Adolfo Vidal} (Nº 51, h.\ 3); y en mayo ---el decreto no trae fecha impresa; se publica en la edición del 22 de ese mes--- el doctor Usandivaras por
\textbf{Augusto Regis}\index{Regis, Augusto} (Nº 58, h.\ 4). Los tres entrantes son los tres que en noviembre
de 1908 habían sido nombrados jurados de reclamos de patentes del departamento (capítulo
\ref{cap:fincas}). \textbf{La comisión municipal que termina 1909 está integrada, en tres de sus
cinco lugares, por el jurado de patentes de 1908.}

El caso de Jacinto Guerrero\index{Guerrero, Jacinto} muestra la escala del pueblo mejor que
cualquier descripción. En enero es miembro de la comisión municipal; en febrero, comisario
suplente de policía del departamento; en marzo, juez de paz suplente, en reemplazo de Augusto
Regis\index{Regis, Augusto}, que había renunciado a ese cargo (B.O.\ Nº 37, 1909, h.\ 3); en abril renuncia a la comisión
municipal. \textbf{El juez de paz propietario era Carlos Matienzo\index{Matienzo, Carlos}, que
firmaba también como secretario del cuadro de rentas municipales.} En un departamento de unos
doscientos inscriptos en el registro cívico ---el decreto de correlación de agosto de 1909
implica ciento setenta y siete inscriptos anteriores---, los cargos de juzgar, mandar la policía,
gobernar el municipio y clasificar el impuesto se repartían entre menos de diez personas.

\textbf{Y 1910 repite el esquema y lo aprieta.} En marzo, Abelardo Figueroa\index{Figueroa, Abelardo}
renuncia a la comisión municipal y entra en su lugar \textbf{Escolástico
Aparicio}\index{Aparicio, Escolástico M.} (B.O.\ Nº 146, 1910, h.\ 3); en abril, el mismo Aparicio es nombrado
\textbf{comisario de policía del departamento} (Nº 148, h.\ 3); en diciembre es \textbf{comisionado de
empadronamiento y clasificación de patentes} de Caldera (Nº 214, h.\ 3) y, en la última edición del año, vuelve a ser nombrado comisario y
figura entre los cinco miembros de la comisión municipal del año siguiente (Nº 219, h.\ 3). \textbf{Cuatro
cargos en un año, y dos de ellos a la vez}: el que mandaba la policía del departamento
integraba el cuerpo que gobernaba el municipio. El decreto de diciembre que designa esas
comisiones repite, palabra por palabra, el argumento de enero de 1909: no se puede aplicar el
artículo 173 de la Constitución. \textbf{Dos años seguidos de excepción es el régimen
ordinario.}

\textbf{Y 1912 lo cierra con la prueba del mecanismo.} El 28 de febrero, el decreto que renueva
los comisarios de todos los departamentos de campaña vuelve a nombrar para Caldera a
\textbf{Escolástico M.\ Aparicio}\index{Aparicio, Escolástico M.} (B.O.\ Nº 320, 1912, h.\ 3):
\textbf{es el quinto año seguido en que el mismo hombre clasifica el impuesto o manda la
policía del departamento}. En mayo, un decreto nombra los jueces de paz «\textit{de acuerdo con
las ternas presentadas por la Comisión Municipal}» (Nº 335, h.\ 3): el circuito que el decreto
de 1909 había ordenado ---el municipio propone, la Provincia designa--- está funcionando y deja
constancia escrita. Y tres semanas después, otro decreto integra esa comisión municipal, que
estaba «\textit{incompleta}» por la renuncia de Augusto Rejis\index{Regis, Augusto}
\textit{[sic]}, con \textbf{Daniel Linares}\index{Linares, Daniel} (Nº 338, h.\ 3), que ese mismo
año figura, en el aviso de remate de San Alejo, como propietario de la finca Angostura que linda
con ella por el este y por el sur.

Dos personas más completan el cuadro de ese año. \textbf{Jacinto Guerrero}\index{Guerrero, Jacinto} es, a la vez, comisario suplente del departamento y juez de paz suplente. Y
\textbf{Eustaquio Murúa}\index{Murúa, Eustaquio}, comisario auxiliar del partido de La Despensa
desde abril, es sorteado en septiembre como \textbf{titular de la segunda mesa receptora de votos}
del departamento, con Adolfo Vidal de primer suplente (Nº 362, h.\ 3). \textbf{El comisario del
partido presidía la mesa donde votaba el partido.}

\textbf{Y en 1913 el mismo hombre cambia de poder sin salir del departamento.} En febrero se
acepta su renuncia a la comisaría y lo reemplaza Julio del Solar (B.O.\ Nº 391, 1913, h.\ 3);
en abril, \textbf{el Poder Ejecutivo lo nombra juez de paz propietario} del departamento, «de
acuerdo con las ternas pasadas por la comisión municipal» (Nº 404, h.\ 3); en agosto y en
octubre firma, como juez de paz, los edictos sucesorios del pueblo; y en noviembre es
designado \textbf{comisionado para revisar las patentes de las casas de comercio} de Caldera
(Nº 448, h.\ 3). \textbf{Seis años seguidos: clasificador de patentes, comisario, juez de paz
y otra vez patentes.} \textbf{Y en mayo de 1914 vuelve al cuerpo municipal}: un decreto del 9
lo nombra miembro de la comisión municipal del departamento en la vacante que dejó
\textbf{Salvador Rosa}\index{Rosa, Salvador} al trasladar su domicilio a la capital (B.O.\
Nº 485, 1914, h.\ 1). \textbf{Son siete años seguidos, y el séptimo lo devuelve al punto de
partida.} La comisión municipal que lo propone para juez tiene ese año tres
miembros ---Daniel Linares\index{Linares, Daniel}, Camilo Cáceres y \textbf{Salvador
Rosa}\index{Rosa, Salvador}--- y el tercero es el mismo acreedor que en 1936 ejecutará a los
herederos Murúa y se quedará con el remate de La Despensa (capítulo \ref{cap:fincas}).

Y el jurado del registro cívico ---el que resuelve los reclamos sobre quién está inscripto para
votar--- lo integran en 1913 y otra vez en 1914 las mismas cuatro personas: Daniel
Linares\index{Linares, Daniel} y Augusto Regis\index{Regis, Augusto} como titulares, Adolfo
Vidal y Silvano Murúa\index{Murúa, Silvano Ignacio} como suplentes (Nº 385, h.\ 2, y Nº 448,
h.\ 3).

\textbf{Y el jurado de patentes de 1915, designado en noviembre de 1914, repite tres de esos
cuatro nombres y cambia uno.} Titulares, Daniel Linares\index{Linares, Daniel} y
\textbf{Bernardo Austerlitz}\index{Austerlitz, Bernardo}; suplentes, Adolfo Videl
\textit{[sic]} y Silvano Murúa\index{Murúa, Silvano Ignacio} (B.O.\ Nº 527, 1914, h.\ 2). Lo
firma el Presidente del Senado en ejercicio del Poder Ejecutivo. \textbf{Es la primera
aparición de Austerlitz en el departamento}, once años antes de figurar como propietario de San
Alejo y Santa Rufina en los pedimentos de 1925; que haya comprado la finca en alguno de los
remates de 1914 es hipótesis y no dato, porque ningún aviso publica el resultado de las
subastas.


\textbf{El resto de los cargos de 1914 confirma el esquema sin novedad.} En abril, el decreto
anual de comisarios de campaña nombra para Caldera a \textbf{Julio del Solar}, el mismo que
había reemplazado a Aparicio en febrero de 1913 (B.O.\ Nº 478, h.\ 1). En febrero se nombra
subcomisario de la \textbf{estación Mojotoro}, «\textit{jurisdicción del Departamento de la
Caldera}», a José Nogales, con antigüedad del 11 de octubre anterior; en julio, otro decreto
reorganiza esa misma subcomisaría y nombra a \textbf{Benjamín López}, dando las gracias al
cesante sin nombrarlo (Nº 469, h.\ 2, y Nº 504, h.\ 2 y 3). Y en noviembre se nombra comisario
auxiliar del \textbf{partido Los Porongos} a Carlos Bongarnier, «\textit{de acuerdo con la
propuesta elevada por el señor comisario de policía del departamento de Caldera}» (Nº 525,
h.\ 2), que es otra vez el circuito de 1909: el municipio o la policía del lugar proponen, la
Provincia designa.

\textbf{Una cosa más, que el archivo de estos años casi no da.} En septiembre de 1914 el
\textbf{juez de paz del departamento} ordena publicar el edicto del juicio sucesorio de don
Salustiano Caro, citando a herederos y acreedores, fechado «La Caldera, septiembre 3 de 1914»
(Nº 510, h.\ 4). El nombre del firmante no se alcanza a leer en la hoja y el año no da otro
acto que lo nombre. \textbf{Es el único acto de 1914 en que el departamento aparece como sede
de un juzgado de paz en funciones}, y el cargo lo ejercía, desde abril de 1913, el mismo
Aparicio.

\begin{cautela}
Los actos dicen los cargos, no las relaciones. Que los mismos apellidos ocupen los cuatro
escalones no prueba acuerdo ni familia: prueba que el conjunto de personas elegibles era chico.
Lo que sí puede afirmarse, porque está en los decretos, es que \textbf{la Provincia designaba,
el municipio proponía en terna y el mismo nombre podía ocupar dos cargos a la vez}.
\end{cautela}

\subsection*{Y 1915 repite el esquema, con una novedad y una retirada}

El año 1915 se leyó edición por edición al cierre de este trabajo: cuarenta y nueve ediciones,
de la 537 a la 586, con la 581 ausente del repositorio y sin ediciones de enero ni de la primera
semana de febrero (apéndice \ref{cap:fuentes}). Deja cinco actos del departamento, y \textbf{ninguno
de ellos es de tierra}.

\begin{ficha}{Los actos del departamento de la Caldera en 1915}
\textbf{26 de enero.} Decreto del Poder Ejecutivo: «\textit{Habiendo terminado el período que fueron
nombrados miembros de la comisión municipal, del departamento de la Caldera los señores Adolfo Vidal
y Angel Fernández}», nombra para ese cuerpo a \textbf{Eustaquio Murúa}\index{Murúa, Eustaquio} y otra
vez a \textbf{Adolfo Vidal}, «\textit{por el término de ley}». Firman Patrón Costas y Julio Cornejo
(B.O.\ Nº 540 del 17 de febrero, h.\ 2, folio impreso 2.077).

\textbf{13 de febrero.} Convocatoria a elecciones: Caldera figura \textbf{en las dos listas del mismo
artículo}, la de los departamentos que eligen representantes al Senado y la de los que eligen
diputados (Nº 539, h.\ 3).

\textbf{Mayo.} Decreto de supresión de puestos «por razones de renta»: entre veintitrés agentes de
comisarías de campaña, \textbf{uno de Caldera} (Nº 552, h.\ 2).

\textbf{2 de septiembre.} Edicto del \textbf{juez de paz del departamento}, \textbf{E.\ M.\
Aparicio}\index{Aparicio, Escolástico M.}, que declara abierto el juicio sucesorio de Asunción Espeche
González de Wayar y cita por treinta días a herederos y acreedores (Nº 568, h.\ 4, y Nº 569, h.\ 3).

\textbf{6 de diciembre.} Decreto que comisiona a una persona por departamento para la clasificación
general de patentes y el empadronamiento de los capitales en giro de 1916: por Caldera,
\textbf{Escolástico M.\ Aparicio}\index{Aparicio, Escolástico M.} (Nº 584, h.\ 2). Un decreto «de la
fecha» nombra los \textbf{jurados} que entenderán en los reclamos de ese empadronamiento: por Caldera,
\textbf{Daniel Linares}\index{Linares, Daniel} y \textbf{Silvano Murúa}\index{Murúa, Silvano Ignacio};
suplentes, \textbf{Adolfo Vidal} y Belisario Nieva (Nº 585, h.\ 2).
\end{ficha}

\textbf{La novedad es de personas, y es la misma forma de siempre.} El hombre que en septiembre firma
como juez de paz del departamento es, en diciembre, el que clasifica sus patentes: \textbf{octavo año
seguido con cargo} para Aparicio\index{Aparicio, Escolástico M.}, que el apéndice \ref{cap:personas}
sigue desde 1908. Y los jurados que van a resolver los reclamos contra esa clasificación son
\textbf{dos propietarios del departamento} ---Linares\index{Linares, Daniel}, dueño de Angostura y de
Los Porongos; Murúa\index{Murúa, Silvano Ignacio}, el que en 1932 hará deslindar cinco fincas
(capítulo \ref{cap:fincas})---, con un suplente, Adolfo Vidal, que once meses antes había sido
nombrado miembro de la comisión municipal por el mismo Poder Ejecutivo que ahora lo nombra jurado.

\begin{cautela}
Conviene decir qué prueba y qué no. \textbf{No prueba ningún acuerdo}, y el decreto que los nombra es
el mismo que nombra jurados en los otros veintidós departamentos. Prueba lo que el capítulo
\ref{cap:prospectiva} sostiene por otra vía: que en un departamento con poco más de quinientos
inscriptos, \textbf{quien clasifica, quien resuelve el reclamo contra esa clasificación y quien
administra el municipio salen de la misma lista, porque no hay otra}. La coincidencia no hizo falta
producirla.
\end{cautela}

\textbf{Y la retirada es de plata.} En mayo la Provincia suprime, por razones de renta, el agente de
la comisaría de campaña de Caldera, uno de los veintitrés que suprime en el mismo artículo. Cuatro
años antes había creado uno para Vaqueros por la obra del puente y por la explotación de canteras. En
las ediciones leídas hasta aquí \textbf{es la primera vez que el archivo registra la supresión de un
cargo del departamento y no su creación}, y ocurre en el año en que ese mismo departamento es
convocado a votar tres veces.

\subsection*{Y 1916 acumula los cargos en dos nombres}

El año 1916 se leyó también edición por edición, y es el tramo más limpio de los seis: las cuarenta
y nueve ediciones de la 587 a la 635, sin un solo salto, con sus 196 hojas (apéndice
\ref{cap:fuentes}). Deja doce actos del departamento, tres de ellos sobre la tierra ---el deslinde de
Loma Bola, el remate de Los Porongos y el segundo deslinde de Santa Mónica y Quitilipe, que el
apéndice \ref{cap:dominio} incorpora--- y nueve sobre los cargos. Los nueve caben en dos nombres.

\textbf{Escolástico Aparicio\index{Aparicio, Escolástico M.} renuncia y vuelve por la puerta de al
lado.} En febrero un decreto acepta su renuncia como miembro de la comisión municipal y nombra en su
lugar a Andrés Marín (B.O.\ Nº 593, h.\ 1). En abril, «\textit{de acuerdo con la terna elevada por la
Comisión municipal de la Caldera}», otro decreto lo nombra \textbf{Juez de Paz propietario} del
departamento «\textit{para el ejercicio del corriente año}» (Nº 601, h.\ 3). \textbf{La terna que lo
propone la eleva el cuerpo que acababa de dejar}, dos meses antes. Es el noveno año seguido con
cargo, y el mecanismo es el que el capítulo describe para 1912: los jueces de paz salen de ternas de
la comisión municipal, y la comisión municipal la designa el gobernador.

\textbf{Eustaquio Murúa\index{Murúa, Eustaquio} junta tres cargos en cinco meses.} El 4 de marzo lo
nombran \textbf{comisario de policía} del departamento, en la vacante que deja Manuel R.\ Cornejo
(Nº 596, h.\ 1). El 16 de marzo, \textbf{expendedor de guías}, «\textit{aceptase la fianza dada en su
favor por don Silvano I.\ Murúa}\index{Murúa, Silvano Ignacio}, \textit{por la suma de 500 pesos
moneda nacional}» (Nº 597, h.\ 2). Y el 5 de agosto, en un decreto que organiza el impuesto al consumo
departamento por departamento, \textbf{receptor del impuesto al consumo} por La Caldera (Nº 616,
h.\ 2). El mismo hombre que en enero de 1915 había sido nombrado miembro de la comisión municipal y
que en septiembre de ese año era candidato a elector de gobernador por el partido oficialista.

\begin{cautela}
\textbf{Tres cargos en un año no son un privilegio: son la aritmética de un departamento chico}, y es
la misma aritmética que produjo los cuatro cargos de Aparicio en 1910. Lo que sí merece registrarse es
\textbf{qué} cargos son. El comisario de policía, el expendedor de guías ---el que extiende el papel
sin el cual no se puede mover ganado--- y el receptor del impuesto al consumo son, los tres, el punto
donde el Estado toca la circulación de bienes del departamento. Y el fiador que responde por el
expendedor es un Murúa, de la familia que quince años después hará deslindar cinco fincas del
departamento (capítulo \ref{cap:fincas}).
\end{cautela}

\textbf{La comisión municipal cambia dos veces en el año} y las dos por renuncia: en febrero sale
Aparicio y entra Andrés Marín; en septiembre sale Adolfo Vidal ---que venía de 1915--- y entra
Salvador Rosa (Nº 624, h.\ 2). Ninguno de los dos decretos acompaña ni cita el escrito de renuncia que
dice haber aceptado, y tampoco se publica la terna que invoca el decreto del juez de paz:
\textbf{tres actos citados y no acompañados en un solo año}, sobre un cuerpo de cinco miembros.

\textbf{Y en diciembre, el departamento aparece por primera vez en la Cámara de Justicia.} La nómina
de juicios sentenciados en el tercer trimestre incluye «\textit{Juan Campilongo versus Amado Juárez,
venido en apelación del Juzgado de Paz «La Caldera», sección contrato}» (Nº 632, h.\ 2). Es la primera
constancia, en las ediciones leídas, de que una sentencia del juzgado de paz del departamento llegó a
segunda instancia.

\subsection*{Dos series que sólo se ven leyendo los nueve años juntos}

Los tramos de 1908 a 1916 se leyeron año por año, y cada barrido registró los cargos de su año. Puestos
en una sola columna, dan dos trayectorias que ninguno de ellos podía ver.

\textbf{La primera es la de Escolástico Aparicio}\index{Aparicio, Escolástico M.}, y el libro ya la
nombra año por año. Completa, es ésta:

{\footnotesize
\begin{longtable}{@{}p{1.1cm}p{11.4cm}@{}}
\toprule
\textbf{Año} & \textbf{Cargo, y acto que lo registra} \\
\midrule
\endhead
1908 & Clasificador de patentes generales de Caldera para 1909 (noviembre) y comisario de policía del departamento para 1909 (diciembre) \\
\addlinespace
1909 & Otra vez clasificador de patentes del departamento, para 1910 \\
\addlinespace
1910 & Miembro de la comisión municipal (marzo, por renuncia de Abelardo Figueroa\index{Figueroa, Abelardo}); comisario del departamento (abril); comisionado de empadronamiento y clasificación de patentes (diciembre) y, el mismo mes, otra vez comisario y otra vez miembro \\
\addlinespace
1911 & \textbf{Comisionado del empadronamiento de los capitales en giro y de la clasificación de patentes para 1912}, por decreto del Ministerio de Hacienda del 9 de noviembre (B.O.\ Nº 295, h.\ 2) \\
\addlinespace
1912 & Comisario de policía del departamento (28 de febrero) \\
\addlinespace
1913 & Renuncia a la comisaría en febrero; \textbf{juez de paz propietario} por terna de la comisión municipal el 18 de abril; firma como tal el sucesorio de Gabriel Ruiloba en octubre; comisionado para revisar las patentes de las casas de comercio en noviembre \\
\addlinespace
1914 & Miembro de la comisión municipal, en la vacante de Salvador Rosa\index{Rosa, Salvador} (9 de mayo) \\
\addlinespace
1915 & Juez de paz del departamento (firma el edicto sucesorio en septiembre); comisionado de la clasificación de patentes de 1916 (6 de diciembre) \\
\addlinespace
1916 & Renuncia a la comisión municipal en febrero; juez de paz propietario en abril, por terna de esa misma comisión \\
\bottomrule
\end{longtable}
}

\textbf{Nueve años seguidos con cargo, bajo tres gobernadores.} Pero lo que la serie completa agrega no
es el número: es \textbf{cuál} es el cargo. En seis de los nueve años ---1908, 1909, 1910, 1911, 1913 y
1915--- lo que ocupa es el puesto que \textbf{clasifica o empadrona a los contribuyentes del
departamento}. La comisaría y el juzgado de paz son la excepción. Este libro lo venía presentando como
el hombre de los cuatro cargos de 1910; el archivo dice que es, antes que nada y durante casi una
década, el que fija cuánto paga cada comercio de La Caldera.

\textbf{La segunda es la de Eustaquio Murúa}\index{Murúa, Eustaquio}, y va en sentido contrario: de un
partido al departamento entero. En 1912 y 1913 es comisario auxiliar del \textbf{partido de La
Despensa}, uno de los trece. En enero de 1915 lo nombran \textbf{miembro de la comisión municipal} y en
septiembre es candidato a elector de gobernador por el partido oficialista. Y en 1916 junta, en cinco
meses, la comisaría de policía del departamento, la expendeduría de guías y la receptoría del impuesto
al consumo. \textbf{Cuatro años para pasar de auxiliar de un partido a los tres cargos que tocan la
circulación de bienes de todo el departamento.}

\textbf{Y los partidos tienen su propia rotación.} Los decretos anuales de comisarios auxiliares
permiten armar, por primera vez, quién mandaba dónde y por cuánto tiempo:

{\footnotesize
\begin{longtable}{@{}p{1.9cm}p{1.95cm}p{1.95cm}p{1.95cm}p{1.95cm}p{1.8cm}@{}}
\toprule
\textbf{Partido} & \textbf{1909} & \textbf{1911} & \textbf{1912} & \textbf{1913} & \textbf{1914} \\
\midrule
\endhead
Calderilla & Genaro Molina & Genaro Molina & Genaro Molina & Genaro Molina & --- \\
\addlinespace
Los Porongos & Teodoro Reyes & Teodoro Reyes & --- & Enrique Leconte & Carlos Bongarnier \\
\addlinespace
La Despensa & Francisco Leeds & Guillermo Wotton & Eustaquio Murúa\index{Murúa, Eustaquio} & Eustaquio Murúa\index{Murúa, Eustaquio} & --- \\
\addlinespace
Vaqueros & Adolfo Agote & --- & Jesús Ramírez & Guillermo Villanueva & --- \\
\addlinespace
San Alejo & --- & Ángel Fernández\index{Fernández, Ángel} & Ángel Fernández\index{Fernández, Ángel} & Manuel Ramos & --- \\
\addlinespace
Los Peñones & Liborio Guerrero & --- & Teodoro Reyes & --- & --- \\
\bottomrule
\end{longtable}
}

Dos cosas saltan. \textbf{Calderilla no rota}: el mismo hombre cuatro años seguidos, y es el único
partido así. Y \textbf{Teodoro Reyes no desaparece de Los Porongos: se muda a Los Peñones}, el mismo
hombre que en 1909 había sido además veedor del Código Rural por Los Porongos. Los cuadros están
armados con los decretos que cada barrido fichó; los guiones no dicen que el partido no tuviera
comisario, sino que el decreto de ese año no lo nombra o no fue fichado.

\subsection*{Y 1917 pone al aserrador en la comisión municipal}

El año 1917 se leyó también edición por edición, y es el más flojo de los once años que este libro leyó enteros: faltan la
edición 640 entera y cuarenta de las 329 hojas, y el entorno no permitió OCR propio, de modo que
doscientas trece hojas se leyeron sólo por la capa de texto del PDF, a unos 150 ppi (apéndice
\ref{cap:fuentes}). Con esa reserva, deja nueve actos del departamento, tres de ellos deslindes que el
apéndice \ref{cap:dominio} incorpora.

\textbf{La comisión municipal se renueva entera el 4 de enero} ---decreto 1056: Daniel
Linares\index{Linares, Daniel}, Andrés Marini y Camilo Cáceres, «\textit{por un nuevo período}»--- y
diecinueve días después, el decreto 1085 nombra jueces de paz del departamento «\textit{de acuerdo con
las ternas elevadas por la H.\ Comisión Municipal}»: propietario, \textbf{Escolástico M.\
Aparicio}\index{Aparicio, Escolástico M.}; suplente, Casimiro Cáceres. \textbf{Décimo año seguido con
cargo}, y la terna la eleva un cuerpo que tenía diecinueve días de renovado. En febrero, el decreto
1101 vuelve a poner a \textbf{Eustaquio Murúa}\index{Murúa, Eustaquio} de comisario de policía del
departamento, el mismo cargo que tenía desde marzo de 1916.

\textbf{Y uno de los tres renovados usa el mes siguiente.} El \textbf{Camilo Cáceres} que el decreto 1056 confirma en la comisión municipal el 4 de enero es el mismo \textbf{Camilo Cácere} ---así, sin la ese final, en el impreso--- que \textbf{quince días después} se presenta «\textit{con títulos bastantes}» a pedir el deslinde, la mensura y la división de condominio de la finca Potrero Gallinato, y obtiene el auto del juez de feria el 19 de enero. \textbf{Quince días entre el nombramiento en el cuerpo que administra el departamento y el escrito que pone en movimiento la justicia sobre una de sus fincas.} Las dos cosas son enteramente legales y no se sigue de ellas ninguna irregularidad: se sigue, otra vez, que en este departamento el que administra y el que pide son la misma lista de apellidos.

\textbf{Y en septiembre entra a ese cuerpo el hombre del aserradero.} El decreto 1377 nombra miembro
de la comisión municipal, en reemplazo de Andrés Marini que renunció, a \textbf{Rafael
Delgado}\index{Delgado, Rafael}. Diecinueve años después, un Rafael Delgado será el dueño del
aserradero del pueblo cuya maquinaria se remata sin base a instancia de una casa comercial de la
capital, y en el mismo año figurará en el nombramiento de jueces de paz del distrito (capítulo
\ref{cap:tierra}). \textbf{Si es la misma persona ---coincidencia de nombre, como siempre, y no
identificación registral---, la única industria documentada del pueblo estuvo sentada en su comisión
municipal desde 1917.}

\textbf{El cuarto nombre del año es el mismo de los tres anteriores.} Daniel
Linares\index{Linares, Daniel}, renovado en la comisión municipal el 4 de enero, es en febrero
\textbf{candidato a Diputado por Caldera} por el partido «Unión Provincial», en la nómina que publica
la Junta de Escrutinio para la elección del 4 de marzo. Y el departamento elige ese año \textbf{un
diputado y ningún senador}, contra el senador y el diputado que le correspondían en 1915: el decreto
de convocatoria lo pone en la lista corta.

\textbf{Y el mismo año, el mismo apellido tiene tierra del otro lado del límite.} Un edicto de deslinde de Campo Santo nombra a \textbf{Daniel Linares}\index{Linares, Daniel} como lindero este de la finca Ántila (B.O.\ Nº 672, h.\ 6), y otro, de la finca «\textit{San Pablo o Puesto Viejo}» del mismo departamento, pone a \textbf{Luis Linares}\index{Linares, Luis} de lindero sud, con el río Mojotoro al oeste (Nº 662, h.\ 3). El miembro de la comisión municipal de La Caldera y candidato a diputado por el departamento era, a la vez, propietario colindante en el departamento vecino.

\begin{cautela}
\textbf{Cuatro actos citados y no acompañados en un año de nueve actos.} Dos son renuncias que el Boletín da por hechas sin publicarlas: la de Andrés Marini a la comisión municipal y la de Emilio Fressart al puesto de subcomisario \textit{ad honorem} del partido de Mojotoro, que el decreto 1419 dice haber aceptado «en la fecha». El tercero son \textbf{las ternas} que el decreto 1085 invoca para nombrar a los jueces de paz y que no se publican. Y el cuarto, los \textbf{títulos y el dictamen del Agente Fiscal} en que se funda el auto de deslinde del Potrero Gallinato. Es el mismo mecanismo que el capítulo documenta para 1916: el cuerpo se rehace, y la justicia se mueve, por escritos que nunca se leen.

\textbf{Y una precisión sobre el perito}, porque el dato podría leerse de más: Herman Pfister\index{Pfister, Hermann} firma los tres deslindes del departamento, pero también es el perito propuesto en otros tres del mismo año que no son de aquí. \textbf{No es el agrimensor de La Caldera: es el agrimensor de la zona}, y que aparezca en todo lo que se mide dice más del tamaño del oficio que de ningún vínculo con el departamento.
\end{cautela}

\subsection*{Y en 1918 el departamento fue municipio electivo durante siete semanas}

El año 1918 se leyó también edición por edición y es el octavo tramo: las ediciones 684 a 736,
534 hojas, sin la 703, sin OCR propio y con un techo de 150 ppi (apéndice \ref{cap:fuentes}).
Dejó \textbf{veintisiete actos del departamento}, ninguno de dominio ---la ficha número
veintiocho del barrido es un falso positivo que este capítulo consigna más abajo. Y dejó, sobre todo, \textbf{el
episodio institucional más importante que este libro encontró antes de 1983}.

\begin{ficha}{Decreto N.\ 1.557 en acuerdo de Ministros --- Salta, 4 de marzo de 1918}
Considerando «\textit{que por la renta recaudada el año próximo pasado, está suficientemente
constatado}» que catorce departamentos y distritos ---\textbf{Rosario de la Frontera, Rosario de
Lerma, Cerrillos, Campo Santo, Metán, La Candelaria, Caldera, Chicoana, Orán, La Viña, Güemes,
Galpón, La Merced y Embarcación}--- «\textit{poseen una renta superior al mínimun de \$ 5.000
establecido}», el Gobernador dispone que desde el \textbf{1.\ de julio} dejen de ser Comisiones
Municipales y pasen a ser \textbf{Municipalidades electivas}.

El artículo 6 pone la inscripción, los actos preparatorios y la elección misma a cargo de las
propias Comisiones Municipales salientes. El artículo 7 manda sortear a los miembros para la
renovación del artículo 173 de la Constitución. Firman Cornejo y Zuviría.
\textit{(B.O.\ Nº 693, 8 de marzo de 1918, h.\ 1 a 3.)}
\end{ficha}

\textbf{Siete semanas después se deshizo, y se deshizo sólo para La Caldera.}

\begin{ficha}{Decreto N.\ 1.660 en acuerdo de Ministros --- Salta, 24 de abril de 1918}
«\textit{Habiéndose incluido entre las Comisiones Municipales, convertidas en Municipalidades
por decreto N 1.557, la del Departamento de la Caldera, cuya renta no alcanza a la suma de \$
5.000, pues en su Cálculo de recursos correspondiente al año próximo pasado, que haciendo a la
suma de \$ 6.471.44, se encuentra agregada la suma de \$ 3.083.94, como saldo del año anterior,
lo que ha motivado un error de apreciación en los datos que han servido de base para dictar el
decreto aludido, siendo en realidad la renta anual de dicha Comisión Municipal de \$ 3.387.50}».

Y considerando «\textit{que por los fundamentos del decreto ya citado y las disposiciones legales
pertinentes, el departamento de la Caldera \textbf{no está en condiciones de ser administrado por
una Municipalidad electiva}}», el artículo 1 deja sin efecto el decreto 1.557 en cuanto la
convierte.

La fecha se imprime «\textit{Salta, Abril 24 de 1913}» [\textit{sic}]: el último dígito se leyó a
600 ppi y la única lectura posible es 1918, porque la edición que lo publica es del 3 de mayo de
1918. Firman Cornejo, Zuviría y Alvarado.
\textit{(B.O.\ Nº 701, 3 de mayo de 1918, h.\ 2, pág.\ 148.)}
\end{ficha}

\begin{cautela}
\textbf{La cuenta cierra y por eso importa.} \$ 6.471,44 menos \$ 3.083,94 da exactamente \$
3.387,50, que es la renta que el decreto declara real. No hay error aritmético: lo que hubo fue
un saldo del ejercicio anterior computado como recurso del año, y esa sola línea contable es la
que decide si un departamento se gobierna a sí mismo o lo gobierna un cuerpo designado.

\textbf{El departamento quedó a mil seiscientos doce pesos con cincuenta centavos de tener
municipalidad electiva.} Es más que todo el presupuesto municipal del departamento diez años
antes, que en 1908 era de mil cuatrocientos ochenta y tres pesos (capítulo \ref{cap:opacidad}).

\textbf{Y la decisión no la tomó nadie contra el departamento.} La tomó una lectura de su propio
cálculo de recursos, hecha en la capital, sobre una planilla que el departamento había elevado.
Es la forma más temprana del mecanismo que el capítulo \ref{cap:opacidad} describe cien años
después: el criterio se aplica sobre un número que el interesado no controla y que nadie está
obligado a explicarle.
\end{cautela}

\textbf{Y un mes después se deshizo para todos.} El \textbf{decreto N.\ 50 del 21 de mayo de
1918}, ya del Interventor Nacional, deja sin efecto el 1.557 \textbf{entero} ---no sólo para La
Caldera--- y suspende la inscripción de electores abierta en los Registros; las Municipalidades y
comisiones existentes continúan hasta que se provean los medios de reemplazarlas
\textit{(B.O.\ Nº 704, h.\ 7 y 8)}. De los catorce departamentos que iban a elegir su gobierno
comunal en julio de 1918, ninguno lo eligió.

\begin{cautela}
Conviene medir qué prueba y qué no.

\textbf{Lo que prueba:} que en 1918 el criterio para tener gobierno municipal propio era un piso
de renta, que La Caldera lo superaba en el papel y no en el cálculo depurado, y que el propio
Estado provincial rehízo la cuenta y la revirtió en siete semanas.

\textbf{Lo que no prueba:} que el departamento hubiera conservado la municipalidad si el decreto
50 no hubiera existido, porque ése la quitó a los catorce. La secuencia de La Caldera es anterior
y es específica: fue el único de los catorce al que se le revisó la renta uno por uno.

\textbf{Y lo que deja establecido para el resto del libro:} que la condición de \textbf{tercera
categoría} con que el departamento aparece en 1983, y que la clasificación estadística M3
todavía arrastra en 2022 (capítulos \ref{cap:poblacion} y \ref{cap:politica}), no empieza con la
Ley 1349 de 1933. En 1918 ya se decidía por renta, y el departamento ya no llegaba.
\end{cautela}

\subsection*{Trece partidos, otra vez, y los mismos apellidos}

El \textbf{decreto 1.564 del 7 de marzo de 1918} nombra los comisarios auxiliares del
departamento «\textit{de acuerdo con las propuestas elevadas por el Comisario del Departamento de
la Caldera}», y al hacerlo enumera en una sola lista \textbf{el Comisario Suplente del Pueblo y
los comisarios auxiliares de trece partidos}: Pueblo, Mojotoro, La
Calderilla, Vaqueros, Potrero de Gallinato, Lesser, \textbf{Yacones}, Potrero de Castillo, Potrero
de Valencia, La Despensa, «\textit{Peñones y Chalchami}» [\textit{sic}], Los Porongos, «\textit{Los
Sauces y Campo Alegre}» y «\textit{San Alejo y Santa Rufina}».

\textbf{Es la misma división que el capítulo \ref{cap:fincas} lee en 1909, nueve años después y
con los mismos nombres.} Y tres de los nombrados vienen de esa lista: \textbf{Liborio
Guerrero}\index{Guerrero, Liborio}, comisario auxiliar de Peñones y Chalchanio en 1909, vuelve al
mismo par de partidos en 1918; \textbf{Zacarías Maidana} y \textbf{Genaro Molina} también son
nombres de 1909. Un cuarto viene del año anterior: \textbf{José Nogales}, que en 1917 figura como
lindero este del Potrero Gallinato en el deslinde de esa finca (apéndice \ref{cap:dominio}), es
en 1918 el comisario auxiliar de ese mismo partido ---y renuncia el 20 de junio, a los tres meses
y medio (decreto 165, B.O.\ Nº 709, h.\ 6).

\begin{cautela}
\textbf{El topónimo Yacones aparece impreso en este decreto y ninguna búsqueda lo devuelve}: la
capa de texto del Boletín lo pierde del todo, y apareció al mirar la imagen. Se consigna porque
es el ejemplo más limpio de lo que el apéndice \ref{cap:fuentes} llama palabra suelta perdida:
una ausencia de la búsqueda que se lee como ausencia del archivo.

\textbf{Y «Chalchami» es una grafía más} del paraje que el Boletín imprime también Chalchanio,
Chalcanio y Chalchamo. Se transcriben todas y no se unifican.
\end{cautela}

\textbf{Dos meses y medio después, la Intervención los suspendió a todos.} El decreto N.\ 27 del
16 de mayo suspende a los subcomisarios auxiliares \textit{ad-honorem} de toda la campaña, que
son precisamente los trece de los partidos de La Caldera, y redistribuye oficiales de línea: al
departamento, el Teniente 1.\ Pedro Greni. Dos días después, el decreto N.\ 39 lo corrige y manda
en su lugar al Teniente \textbf{Horacio Hornstein} ---impreso también Ornstein---, que en junio
pasa además a ocupar interinamente la Receptoría de Rentas del departamento, es relevado el 6 de
diciembre y el 26 recibe la aprobación de su conducta tras un sumario que el Boletín no publica.

\begin{cautela}
\textbf{Es el reemplazo de los vecinos por la tropa, y dura siete meses.} En marzo los trece
partidos del departamento estaban a cargo de trece vecinos propuestos por el comisario
departamental. En mayo los trece quedan suspendidos y el departamento entero queda a cargo de un
teniente destinado desde la capital, que además cobra sus rentas. \textbf{La persona más presente
del departamento en el Boletín de 1918 no es del departamento.}
\end{cautela}

\subsection*{Y el municipio todavía proponía a los jueces de agua}

Dos decretos del año nombran \textbf{Jueces de Agua} para los distritos de Caldera y de La
Calderilla, y los dos lo hacen «\textit{vista la nota que antecede del Sr.\ Comisionado Municipal
de Caldera, proponiendo \ldots\ las personas que han de desempeñar los cargos}». El
\textbf{decreto 113 del 23 de septiembre} nombra a \textbf{Ángel Fernández} para Caldera y a
\textbf{Manuel A.\ Cornejo} para Calderilla; el \textbf{decreto 190 del 7 de noviembre}, seis
semanas después, nombra a \textbf{Manuel A.\ Cornejo Arias} para Calderilla y a \textbf{Rafael
Fernández} para el pueblo de la Caldera. El segundo no cita al primero ni lo deja sin efecto.

\begin{cautela}
\textbf{Esto llena un tramo del hueco que el capítulo \ref{cap:planteo} declara.} Ese capítulo
sitúa el punto de partida de la voz municipal sobre el agua en 1885 ---cuando la Municipalidad de
La Caldera concedió treinta litros por segundo del Wierna--- y su pérdida en 1942, cuando el
derecho se inscribe en los registros provinciales; y declara que los sesenta y dos años del
medio no se leyeron de corrido.

\textbf{1918 dice algo sobre esos años.} La facultad ya no es conceder: es \textbf{proponer}. El
municipio eleva la terna y el Interventor Nacional firma el nombramiento. Y la magistratura del
agua sigue organizada por distrito ---Caldera y Calderilla--- como en el reparto de 1866 de Campo
Santo que el capítulo \ref{cap:planteo} cita como mecanismo.

\textbf{Dos nombramientos en seis semanas para los mismos dos distritos, con nombres que no
coinciden y sin que el segundo mencione al primero}, es además un dato sobre el registro: no se
concilia y se transcribe.
\end{cautela}

\subsection*{Lo demás del año, en orden}

\textbf{La escuela.} El decreto 1.653 del 20 de abril jubila a la \textbf{Directora de la Escuela
de la Caldera, señorita Juana Luque}, a razón de ciento cuatro pesos cincuenta centavos mensuales
desde el 1.\ de mayo. \textbf{No es el primer nombre de la escuela del pueblo que este relevamiento
encuentra}: en 1911 ya se había jubilado otra directora, Jacoba Luquez de Gómez (capítulo \ref{cap:presencia}), y el capítulo \ref{cap:ausencias} registra las dos.

\textbf{Las minas.} Dos avisos, los dos de \textbf{Miguel A.\ Sosa} y \textbf{Dermidio
Arancibia}: el pedimento de la mina de plomo «\textit{20 de Febrero}», ubicada en la finca
«\textit{Cerro Nevado}» de propiedad de los señores \textbf{Garzón y Pintos}, jurisdicción del
departamento (B.O.\ Nº 711, 12 de julio); y la denuncia por despueble de la mina
«\textit{Sara}», también del departamento (B.O.\ Nº 716). Y un tercer acto que los ordena: el
\textbf{decreto 157 del 20 de junio} subdivide la provincia en distritos mineros para el registro
gráfico del Departamento Topográfico, y pone a Caldera en el \textbf{V, Distrito del Toro}, junto
con Rosario de Lerma, Cerrillos, Chicoana, la Capital, Campo Santo y Guachipas. \textbf{Es el
único acto del año que ubica al departamento dentro de una división administrativa nueva}, y el
decreto que organiza el registro gráfico no acompaña el gráfico que manda llevar.

\textbf{Las elecciones.} Dos convocatorias sucesivas ---15 de septiembre, por decreto de julio;
15 de diciembre, por decreto en acuerdo de Ministros de octubre--- declaran las dos que Caldera
elige \textbf{un Diputado, un Senador y dos Electores}. La primera no llegó a realizarse. De la
segunda salen los circuitos del año.

\begin{ficha}{Colegio Electoral N.\ 3 --- Caldera, circuitos y mesas (B.O.\ Nº 730, 22 de noviembre de 1918, h.\ 7, pág.\ 477)}
\textbf{Mesa N.\ 1}, en Caldera, Pueblo, \textbf{Escuela Elemental Mixta}. Circuito: Pueblo,
Potrero de Castillo, Peñones y San Alejo.

\textbf{Mesa N.\ 2}, en Caldera, Pueblo, \textbf{Iglesia Parroquial}. Circuito: Angosto, Angosto
de Arrieta, Angosturo, Campo Alegre, Calderilla, Corral de Barrancas, Chalchanio, Despensa,
Despensilla, El Acheral, El Churcal, El Monte, Las Garrapatas, Los Porongos, Los Sauces, Mesadas,
Potrero de Valencia, Quesería, Santa Rufina, Wierna y Yacone.

\textbf{Mesa N.\ 3}, en Vaqueros, \textbf{Escuela Infantil Provincial}. Circuito: Las Lagunas,
Lesser y Vaqueros.

\textbf{Mesa N,\ 4} [\textit{sic}], en Mojotoro, \textbf{Oficina de Correos}. Circuito: Mojotoro,
Mojotorito, «\textit{Potrero, de Gallinato}» [\textit{sic}] y Túnel.
\end{ficha}

\begin{cautela}
\textbf{Es el inventario de parajes más completo que este libro tiene para la primera mitad del
siglo}, y conviene leerlo junto al padrón de parajes de 1940 que transcribe el capítulo
\ref{cap:fincas}. Trae \textbf{treinta y dos nombres} y siete que este relevamiento no había
registrado nunca: Despensilla, Corral de Barrancas, El Churcal, El Monte, Las Garrapatas, Túnel
y Las Lagunas. «Angosturo» no es un nombre nuevo sino una grafía más de la Angostura que este
libro ya sigue desde 1912, como «Yacone» lo es de Yacones, «Mesadas» de Las Mesadas y
«Quesería» de Quecería.

\textbf{Y trae la geografía institucional del departamento en cuatro puntos}: la escuela del
pueblo, la iglesia parroquial, la escuela de Vaqueros y la oficina de correos de Mojotoro. Son
los cuatro únicos edificios públicos que el Estado provincial reconoce en el departamento en
1918. Este libro no verificó si alguno de los cuatro es el mismo edificio que hoy ocupa ese
lugar.

\textbf{Un falso positivo que conviene dejar dicho}, porque es nuevo: hay un paraje llamado
\textbf{Caldera en el departamento de Rivadavia}, dentro del circuito de la Mesa N.\ 1 de su
Colegio Electoral, y en la misma edición. En la capa de texto el renglón es indistinguible de una
mención de este departamento; lo resuelve el encabezado del bloque, tres renglones más arriba.
\end{cautela}

\textbf{Y las autoridades de las mesas ponen doce nombres.} Preside la Mesa 1 \textbf{Manuel A.\
Cornejo}, el mismo que los decretos 113 y 190 nombran juez de agua de la Calderilla; su segundo
suplente es \textbf{Augusto Regis}\index{Regis, Augusto}, que este libro sigue desde 1908. La
Mesa 2 la preside \textbf{Marcos E.\ Aparicio}, apellido del jurado de patentes de 1908 y del
juez de paz de 1909 a 1917. En la Mesa 3 el primer suplente es \textbf{Enrique J.\ Rauch}. Y en
la Mesa 4, en Mojotoro, el presidente es \textbf{Gregorio Nogales}, el segundo suplente
\textbf{Juan Cruz Nogales} y el primer suplente \textbf{José Yugra González} ---el apellido Yugra,
que los deslindes de la Calderilla de 1917 registran como propietario lindero (apéndice
\ref{cap:dominio}), aparece aquí como autoridad de mesa.

\begin{cautela}
\textbf{La superposición de 1909 vuelve a verse en 1918, y con otros nombres.} El presidente de
una mesa receptora de votos del departamento es, en el mismo año, el juez de agua de uno de sus
dos distritos, nombrado a propuesta del propio municipio. No se afirma nada sobre su conducta:
se registra que en un departamento de esta escala el conjunto de personas que pueden ocupar los
cargos es lo bastante chico como para que la coincidencia no tenga que producirse (capítulo
\ref{cap:prospectiva}).
\end{cautela}

\textbf{Y el único acto del año fechado en el departamento} es un edicto sucesorio: el juicio de
\textbf{Don Delfín Chaile}, declarado abierto por el Juez de Paz Propietario Departamental y
firmado «\textit{Caldera, Septiembre 27 de 1918. C.\ Cáceres, Juez de Paz}». Es el
\textbf{Casimiro Caseres} que el decreto 1697 del 6 de mayo había nombrado juez de paz
propietario del departamento, a propuesta en terna de la Comisión Municipal, junto con el
suplente Gervasio Chalar.

\subsection*{Y en 1919 la Constitución dice por qué, y la razón ya no es la renta}

El año 1919 se leyó también edición por edición y es el noveno tramo: las ediciones 738 a 786,
49 ediciones, 440 hojas de los 478 folios del año (apéndice \ref{cap:fuentes}). Dejó nueve actos
del departamento, ninguno de dominio. Y el primero de ellos, de enero, vuelve sobre el asunto que
ocupó todo 1918 y le cambia el fundamento.

\begin{ficha}{Decreto N.\ 33 --- Salta, 15 de enero de 1919}
«\textit{Siendo de urgente necesidad proceder a la organización del personal de las Comisiones
Municipales de los Departamentos de Campaña}», y considerando

\begin{quote}
«\textit{Que el art.\ 16 de la Ley Orgánica de Municipalidades en que se creaba una Municipalidad
en cada uno de los Departamentos mencionados en el mismo, \textbf{ha sido derogado por el art.\
175 de la Constitución, que exije para la existencia de Municipalidades autónomas un centro
urbano de cinco mil habitantes, de que carecen todos los Departamentos de Campaña}; y hasta tanto
se obtenga la reforma constitucional para conseguir la creación de la Municipalidad electiva en
todos los Departamentos de la Provincia, de acuerdo con el programa fundamental de este
Gobierno}» [\textit{sic} «exije»].
\end{quote}

Y nombra los miembros de las Comisiones Municipales de la campaña. Los de Caldera son
\textbf{Domingo Royo}, \textbf{Francisco Sosaya}, \textbf{José Molina}, \textbf{Julio Mercado} y
\textbf{Camilo Cáceres}, sin distinguir presidente.
\textit{(B.O.\ Nº 738, 17 de enero de 1919, h.\ 5, pág.\ 599.)}
\end{ficha}

\begin{cautela}
\textbf{Lo que esto corrige de 1918.} El decreto 1.557 de marzo de 1918 convirtió catorce
Comisiones Municipales en Municipalidades electivas por tener \textbf{renta} superior a cinco mil
pesos, y el 1.660 sacó a La Caldera de la lista rehaciendo esa cuenta. Diez meses después, el
gobierno siguiente dice que el umbral no era ése: que el artículo de la Ley Orgánica en que se
fundaba aquel decreto \textbf{estaba derogado por la Constitución}, y que lo que la Constitución
exige es un \textbf{centro urbano de cinco mil habitantes}.

\textbf{Son dos umbrales distintos con la misma cifra}, uno de pesos y otro de personas, y el
segundo deja afuera a todos los departamentos de campaña, no sólo a éste. Vistos juntos, los dos
actos dicen que en 1918 y 1919 nadie sabía con certeza cuál era la regla: un gobierno aplicó la
renta, el siguiente declaró que la regla era la población y que el artículo que fundaba al
primero ya no existía.

\textbf{Y los cinco mil habitantes tienen descendencia.} Es la misma cifra que la Ley 1349 de
1933 fijará como techo de la tercera categoría municipal, y que el capítulo \ref{cap:poblacion}
encuentra todavía operando en la clasificación estadística de 2022. \textbf{Viene del artículo
175 de la Constitución provincial y ya estaba en 1919.}

\textbf{Lo que este libro no puede decir} es si esa lectura del artículo 175 era la correcta ni
si la reforma constitucional que el decreto anuncia como programa se intentó. El texto de la
Constitución entonces vigente y el del artículo 16 de la Ley Orgánica quedan pedidos en el
apéndice \ref{cap:pedidos}.
\end{cautela}

\subsection*{De trece partidos a diez}

El \textbf{decreto N.\ 213} nombra los comisarios de partido del departamento para el año, «\textit{en
mérito de las propuestas presentadas por intermedio del señor Jefe de Policía}». Y nombra
\textbf{diez}, no trece:

\begin{quote}
Mojotoro, \textbf{Silverio Caseres}; Vaquero, \textbf{José A.\ Molina}; Gallinato, \textbf{Julio
Zambrano}; Calderilla, \textbf{Delfín Chaile}; Lesser, «\textit{Pdero Vale}» [\textit{sic}];
Yacones, \textbf{Juan Cuchui}; «\textit{Potrero de Castilla}» [\textit{sic}], \textbf{Zacarías
Maidana}; San Alejo, \textbf{Manuel Ramos}; El Sauce, \textbf{Domingo Royo}; el Monte,
\textbf{Liborio Guerrero}.
\end{quote}

\textit{(B.O.\ Nº 745, 7 de marzo de 1919, h.\ 7. El artículo 1.\ imprime «Camisarios»; se
transcribe tal cual.)}

\begin{cautela}
\textbf{La lista de partidos se acorta y se renombra, y las dos cosas importan.} En 1909 y en
1918 eran trece; en 1919 son diez. Desaparecen de la enumeración el Pueblo, Potrero de Valencia,
La Despensa, Los Porongos, Peñones, Santa Rufina y Campo Alegre; aparecen dos nombres que este
libro no había registrado como partido: \textbf{El Sauce} y \textbf{el Monte}.

\textbf{«El Monte» no es nuevo del todo}: es uno de los parajes que el Colegio Electoral de 1918
enumera en el circuito de la Mesa N.\ 2. Lo que cambia es su rango: de paraje de un circuito a
partido con comisario. Y en mayo, cuando Liborio Guerrero\index{Guerrero, Liborio} renuncia, el
decreto lo llama «\textit{Comisario Auxiliar de los partidos de \textbf{Monte y Chalchalnio}}»
[\textit{sic}], de modo que el par Peñones--Chalchanio de 1909 y 1918 se ha vuelto el par
Monte--Chalchalnio.

\textbf{Este libro no puede decir si la división cambió o si cambió la enumeración.} Un decreto
de nombramientos nombra los partidos que tienen candidato propuesto, no necesariamente todos los
que existen. Lo que queda establecido es que \textbf{la nómina publicada pasa de trece a diez en
un año}, y que el paraje que el Boletín escribe de cuatro maneras distintas ---Chalchanio,
Chalcanio, Chalchamo y Chalchalnio--- sigue ahí.
\end{cautela}

\textbf{Y el tercer año seguido con Liborio Guerrero en el mismo lugar.} Comisario auxiliar de
Peñones y Chalchanio en 1909, del mismo par en 1918, del Monte en marzo de 1919 y renunciante en
mayo, cuando lo reemplaza \textbf{Héctor Echanique} ---así, con a, donde 1909 imprime Echenique.
\textbf{Y es el mismo paraje}: el Héctor Echenique que en 1909 figura como veedor del Código
Rural de Chalcanio (capítulo \ref{cap:fincas}) es, diez años después, el comisario auxiliar de
Monte y Chalchalnio. \textbf{El cargo cambia de nombre y de régimen; el hombre y el paraje son
los mismos.}

\subsection*{La Comisión Municipal de 1919, y tres apellidos que duran}

De los cinco miembros que el decreto 33 nombra, cuatro reaparecen en otros actos del propio año o
en otros años del libro.

\textbf{Domingo Royo} y \textbf{José A.\ Molina} son, en marzo, comisarios de partido de El Sauce
y del Vaquero: \textbf{los mismos dos meses en que integran el cuerpo que gobierna el
departamento}. \textbf{Camilo Cáceres} venía de la comisión municipal de 1917, y en 1917 se había
presentado «\textit{con títulos bastantes}» al deslinde del Potrero Gallinato. \textbf{Francisco
Sosaya} es el apellido del peticionante del deslinde de Campo Alegre de 1914. Y \textbf{Julio
Mercado} es el nombre que el apéndice \ref{cap:dominio} registra como comprador en un remate de
1930 y que el capítulo \ref{cap:politica} encuentra en la comisión municipal de 1928.

\begin{cautela}
\textbf{Es la superposición de siempre, ahora con fecha exacta.} El comisario de policía de un
partido y el miembro de la comisión municipal del departamento son, en dos casos de cinco y en el
mismo trimestre, la misma persona. No se afirma nada sobre su conducta ni se infiere acuerdo: se
registra que el cuerpo que eleva las ternas de jueces de paz ---y que las elevó, porque el decreto
117 del 4 de febrero nombra a \textbf{Casimiro Cáceres} juez de paz propietario y a
\textbf{Carmelo Sivila} suplente «\textit{vista las ternas presentadas por la Comisión Municipal}»
[\textit{sic}]--- tiene dos de sus cinco miembros bajo las órdenes del Jefe de Policía.

Y \textbf{Julio Mercado}, si es el mismo, atraviesa \textbf{veintiún años} del archivo: veedor del
Código Rural de Los Sauces y Angostos de Arias y lindero al este del Angosto de Arrieta en 1909,
comisión municipal en 1919 y otra vez en 1928, y en el remate de 1930 el edicto ya lo nombra como
«\textit{herederos de Julio Mercado}». El apellido vuelve en 1972 entre los expropiados para el
embalse (apéndice \ref{cap:dominio}). \textbf{Coincidencia de nombre, como siempre, y no
identificación registral.}
\end{cautela}

\subsection*{Lo demás de 1919}

\textbf{La primera plaza rentada de policía en un partido.} En septiembre, el \textbf{decreto N.\
522} crea la \textbf{Sub-Comisaría de Mojotoro} «\textit{con sueldo}» y nombra para ella a
Silverio Cáceres ---el mismo Silverio Caseres que en marzo era comisario de partido---, con
sueldo retroactivo al 1.\ de agosto. En 1918 Mojotoro aparecía en el Boletín por dos soldados de
tropa con viático de cincuenta centavos; en 1919 tiene plaza con presupuesto.

\begin{cautela}
\textbf{El acto no nombra al departamento en ningún renglón}, y por eso ninguna búsqueda por el
topónimo lo devuelve: apareció buscando por el nombre de la persona. Es el único acto del
departamento en todo 1919 que llegó por esa vía, y conviene decirlo donde se lo usa, porque mide
exactamente lo que un barrido por término de búsqueda no encuentra (apéndice
\ref{cap:fuentes}).
\end{cautela}

\textbf{El registro civil, dos veces.} En enero el decreto 49 nombra encargado de la oficina de
Caldera a \textbf{Leoncio A.\ Mascietti}; y en septiembre se publica la Ley de Presupuesto
para el ejercicio de agosto a diciembre, que asigna en su ítem VII \textbf{un encargado de registro civil para cada una de
cincuenta y cinco localidades de la campaña, Caldera entre ellas, a cuarenta y cuatro pesos cada
uno «\textit{inclusive gastos de oficina y alquiler}»}. El comisario de policía del departamento,
nombrado en la primera edición del año, es otro Mascietti: \textbf{Ambrosio L.} No se infiere
parentesco.

\begin{cautela}
\textbf{Una cuenta del presupuesto que no cierra, y es la que paga a los comisarios de este
departamento.} El ítem V, Policía de la Campaña, imprime una columna mensual cuyos renglones suman
\$16.990 y un subtotal de \$16.930: sesenta pesos de diferencia. La lectura que cerraría la
columna es que el primer renglón fuera 4.800 ---que es exactamente treinta comisarios a ciento
sesenta pesos--- y no 4.860; mirado a 500 ppi, el dígito impreso es un seis. \textbf{Se declaran
las dos lecturas y no se corrige.} La columna «Total» encadena con el subtotal impreso y no con el
sumado, y su multiplicador confirma que el ejercicio es de cinco meses. Los otros catorce
subtotales del cuadro cierran exactos.
\end{cautela}

\textbf{Y el año se gobierna con la firma partida.} La mancheta de las cuarenta y nueve tapas dice
«\textit{Gobierno del Dr.\ don Joaquín Castellanos}», y sin embargo el Ejecutivo lo ejerce el
\textbf{Presidente del Senado} en treinta y dos de ellas, en dos interinatos separados por un
regreso del titular: de enero a junio, de julio a septiembre, y Castellanos en el medio y al
cierre. \textbf{Ningún decreto de apertura ni de cierre de esos interinatos aparece publicado.} La
razón de fondo está en un decreto de mayo, que recuerda que las Cámaras se constituyeron
«\textit{en virtud de la convocatoria hecha por la Intervención Nacional por decreto de fecha 31
de diciembre del año ppdo.}»: \textbf{1918 cerró bajo intervención federal, y 1919 abre con la
Legislatura recién reabierta.}


\begin{figure}[p]
\centering
\includegraphics[width=\textwidth,height=0.86\textheight,keepaspectratio]{fig-parajes-nominas.png}
\caption[La toponimia del departamento en doce nóminas oficiales, 1909--1940]{\textbf{La toponimia
del departamento en doce nóminas oficiales, 1909--1940.} Cada columna es una lista que el Boletín
Oficial publicó ---los comisarios auxiliares de policía de 1909, 1918 y 1919, los veedores del
Código Rural de 1909, y los circuitos electorales del Colegio Nº 3 de 1918 a 1923, los del Nº 5 de 1926 y 1927 ---que ya no cubren Vaqueros, desprendido en el Nº 6--- y los del Nº 2 de 1938 y de 1940, que son la misma lista nombre por nombre---; un cuadrado indica que ese
nombre figura en ella. La columna de la derecha dice si el topónimo tiene hoy lugar en el campo
«finca» del catastro parcelario, en la Base de Asentamientos Humanos, en los dos o en ninguno.
Los once años están leídos edición por edición, y 1922 sólo hasta el 21 de julio (apéndice \ref{cap:fuentes}). Elaboración propia
con el catastro parcelario de la Provincia (IDESA, consultado en septiembre de 2026) y BAHRA;
la tabla y el dibujo los produce \texttt{scripts/tabla\_parajes.py} del repositorio
\texttt{mapas\_caldera} (apéndice \ref{cap:fuentes}).}
\label{fig:parajesnom}
\end{figure}

\begin{figure}[p]
\centering
\includegraphics[width=\textwidth,height=0.86\textheight,keepaspectratio]{fig-parajes-mapa.png}
\caption[Dónde caen hoy los topónimos de 1909--1940]{\textbf{Dónde caen hoy los topónimos de
1909--1940.} A: las parcelas cuyo campo «finca» lleva uno de los cincuenta y dos nombres de la
lámina \ref{fig:parajesnom}, coloreadas según en cuántas de las doce nóminas aparece ese nombre;
los triángulos son los parajes que BAHRA registra en el departamento. B: el corredor La
Caldera--Vaqueros. \textbf{El color no mide superficie ni dominio}: mide permanencia en el
registro del Estado. La geometría es la de la parcela que hoy lleva ese nombre, y la coincidencia
es de nombre y no de deslinde. Veinticuatro de los cincuenta y dos topónimos no se dibujan, y se
enumeran en el recuadro. Elaboración propia con el catastro parcelario de la Provincia (IDESA,
consultado en septiembre de 2026), BAHRA y la cartografía de radios del Censo 2022 (INDEC);
dibujo de \texttt{scripts/mapa\_parajes.py} del repositorio \texttt{mapas\_caldera}
(apéndice \ref{cap:fuentes}).}
\label{fig:parajesmapa}
\end{figure}

\subsection*{Doce nóminas del mismo territorio, y lo que queda de ellas}

Entre 1909 y 1940 el Estado provincial publicó \textbf{doce listas} de las unidades en que se
divide este departamento, y \textbf{sólo dos de las doce coinciden entre sí}. Tres son de policía ---los
comisarios auxiliares de 1909, de 1918 y de 1919---, una es del Código Rural ---los veedores de
1909---, y ocho son electorales ---los circuitos de 1918, 1921, 1922, 1923, 1926, 1927, 1938 y 1940---. Las doce
salen de tramos leídos edición por edición, de modo que son la lista entera y no una muestra.

Puestas una al lado de la otra dan \textbf{cincuenta y dos nombres distintos}, y el resultado se
lee de un golpe en la lámina \ref{fig:parajesnom}.

\begin{cautela}
\textbf{Tres cosas se leen ahí, y las tres importan para lo que este libro sostiene.}

\textbf{La primera: hay un núcleo que no se mueve.} \textbf{Diez} nombres aparecen en once o en
las doce listas, con treinta y un años y cuatro autoridades distintas de por medio:
Los Yacones, Potrero de Castillo, San Alejo,
Santa Rufina, Los Porongos, La Despensa, Potrero de Valencia, Peñones, Chalchanio y Los Sauces.
\textbf{Vaqueros, Mojotoro, Lesser y Potrero de Gallinato salieron de esa lista por una sola
razón}: desde 1926 están en el Colegio Electoral Nº 6 y no en el Nº 5, de modo que faltan de las
dos últimas nóminas del Colegio 5 sin haber dejado de existir. \textbf{Son los mismos con los que el capítulo \ref{cap:cot} muestra
que el municipio se describe a sí mismo en el artículo 6.º de su plan de 2015}, cien años después.
\textbf{Y hay un decimoquinto que se cae de esa lista por una sola ausencia}: \textbf{La
Calderilla} está en diez de las doce y falta justamente en la de 1922, donde el mismo circuito
imprime «Candelaria» y donde 1921 y 1923 imprimen Calderilla. Es la discordancia que este capítulo registra más arriba y no resuelve, y
es la única razón por la que el paraje que da nombre a medio departamento no figura entre los
permanentes.

\textbf{La segunda: alrededor de ese núcleo, la división cambia todo el tiempo.} Los veedores del
Código Rural de 1909 usan una partición distinta de la de la policía del mismo año ---aparecen
Wierna, Bandera Angosta, Mojotorillo y Chaguadero, que la lista policial no nombra---; 1918
enumera trece partidos y 1919, diez; y el circuito electoral de 1918 nombra treinta y dos parajes,
el de 1921 treinta y cuatro, el de 1922 otros treinta y cuatro que no son los mismos, el de 1923 los de 1921 otra vez, los de 1926 y 1927 veintiséis cada uno ---sin Vaqueros--- y los de 1938 y 1940 cuarenta y seis, \textbf{que son la misma lista}. \textbf{Ninguna de las siete listas
deslinda nada}: todas
nombran, y el libro no puede saber si lo que cambia es el territorio o la enumeración.

\textbf{Y la tercera, que es la que no se esperaba: veinticuatro de los cincuenta y dos nombres no
tienen hoy lugar en ningún nomenclador.} No están en el campo «finca» del catastro parcelario ni
en la Base de Asentamientos Humanos. Entre ellos hay siete que el circuito electoral de 1918
nombra como lugares donde vive gente que vota ---Despensilla, Corral de Barrancas, El Churcal, El
Monte, Las Garrapatas, El Túnel y La Despensa---, y nombres que los circuitos de 1921 y el padrón de 1940 todavía
registraban, como El Bordo, la Estación del ferrocarril, Severino, Los Chusos, Higueritas, Los
Matos y Nogalito. \textbf{Una parte de la toponimia con la que el Estado organizaba el voto y la
policía de este departamento se borró del registro oficial sin que ningún acto lo dispusiera}, y
el capítulo \ref{cap:ausencias} la retoma como lo que es: una ausencia medida.
\end{cautela}

\begin{cautela}
\textbf{Qué prueba la lámina \ref{fig:parajesmapa} y qué no.} Prueba que veintiocho de los
cincuenta y dos topónimos siguen teniendo una geometría asociada, y dónde caen. \textbf{No prueba
que esa geometría sea la de entonces}: la coincidencia es de nombre y no de deslinde, y el
catastro no dice si los límites actuales son los de 1909. Por eso el mapa no dibuja polígonos de
partido: \textbf{ninguno de los siete instrumentos deslinda las unidades que nombra}, y trazar sus
límites sería inventar una frontera que el archivo no da. Esa correspondencia entre los nombres y
las matrículas actuales es, exactamente, lo que el apéndice \ref{cap:pedidos} reclama a la
Dirección General de Inmuebles.
\end{cautela}

\subsection*{Y 1920 es un año de una sola plaza y de una elección}

El año 1920 se leyó también edición por edición y es el décimo tramo: las ediciones 788 a 839,
48 ediciones, 526 hojas distintas sobre seiscientos folios (apéndice \ref{cap:fuentes}). Dejó
diez actos del departamento, ninguno de dominio, y se reparte en dos materias y nada más: una
plaza de policía y una elección.

\subsubsection*{La plaza que se creó en 1919 rota tres veces en 1920}

En septiembre de 1919 el decreto 522 creó la Sub-Comisaría de Mojotoro «\textit{con sueldo}».
\textbf{En 1920 el puesto cambia de manos tres veces}, y ninguno de los tres decretos cita al
anterior:

\begin{ficha}{El puesto de Sub-Comisario de Mojotoro en 1920}
\textbf{Decreto N.\ 677, 4 de febrero.} Acepta la renuncia de \textbf{Silverio Cáceres},
«\textit{del puesto de Sub Comisario de Policía de Mojotoro, jurisdicción del Departamento de la
Caldera}». No dice quién lo reemplaza. \textit{(B.O.\ Nº 795.)}

\textbf{Decreto N.\ 724, 26 de febrero.} Nombra a \textbf{José Agustín Martínez} «\textit{para
ocupar el puesto vacante de Sub-Comisario de Mojotoro, \textbf{partido} del departamento de la
Caldera}», a propuesta de la Jefatura de Policía. \textbf{No menciona la renuncia que el propio
Boletín había publicado tres semanas antes.} \textit{(B.O.\ Nº 798.)}

\textbf{Decreto N.\ 901, 16 de junio.} Nombra a \textbf{Ciro Cáceres} Sub-Comisario de Policía de
Mojotoro, «\textit{Departamento de La Caldera}», por la nota N.\ 1.622 M.-19 de la Jefatura de
Policía, que no se publica. \textbf{Tampoco dice que Martínez haya renunciado ni lo nombra}: la
vacante se declara sin nombrar a quien la dejó. \textit{(B.O.\ Nº 816.)}
\end{ficha}

\begin{cautela}
\textbf{Tres cosas se leen ahí.}

\textbf{La primera es la rotación.} La única plaza rentada de policía \textbf{de partido} que el
departamento consiguió en toda esta serie ---la comisaría departamental sí estaba presupuestada,
y el ítem V de 1919 paga una treintena de comisarios de campaña, según cuál de las dos lecturas del renglón se tome--- duró cuatro meses y medio en manos de su primer titular y tres meses y medio en las del segundo
---si es que Martínez llegó a tomar posesión, cosa que el Boletín no publica---. Entre la renuncia
del primero y el nombramiento del segundo pasaron tres semanas.

\textbf{La segunda es de forma, y es la que interesa a este libro.} De los tres decretos, dos
\textbf{declaran una vacante sin nombrar a quien la dejó}. El acto dice que el puesto está vacante
y no dice por qué ni desde cuándo. Es el mismo mecanismo que el capítulo \ref{cap:opacidad}
describe cien años después con otro contenido: lo publicado permite saber que algo se decidió y
no qué se decidió.

\textbf{Y la tercera es de nombre.} La misma unidad se llama «\textit{jurisdicción del
Departamento}» en febrero y «\textit{partido del departamento}» tres semanas después, en dos
decretos del mismo gobierno. La nómina de partidos que 1909, 1918 y 1919 publican cada año deja
de publicarse: \textbf{1920 no trae lista de partidos}, y Mojotoro aparece solo.
\end{cautela}

\textbf{Y en octubre se cubre el registro civil, también sobre una vacante sin dueño.} El decreto
N.\ 1140 del 25 de octubre nombra encargado de la Oficina del Registro Civil de La Caldera a
\textbf{Casimiro Cáceres} ---el mismo que fue juez de paz propietario del departamento en 1918 y
en 1919--- «\textit{con antigüedad al 5 de Agosto ppdo.\ fecha desde la cual presta servicios en
carácter de interino}». \textbf{El interinato que el decreto reconoce nunca se publicó}, y la
vacante se declara otra vez sin nombrar a quien la dejó.

\subsubsection*{La elección de 1920, y un departamento que figura dos veces}

Nueve de las dieciséis apariciones del año son electorales. La convocatoria es del 6 de febrero,
para el comicio del 7 de marzo, y trae un dato que conviene poner entero porque no se concilia.

\begin{ficha}{Decreto N.\ 685 --- Salta, 6 de febrero de 1920}
«\textit{Art.\ 1.º---Convócase al pueblo de la Provincia a elecciones de Senadores y Diputados
provinciales, que se llevará a cabo el 7 de Marzo próximo en la siguiente forma: Por la Capital
cinco Diputados; Molinos un Diputado; \textbf{Caldera un Diputado}; Cafayate un Senador;
Guachipas un Diputado; Chicoana dos Diputados; Metán un Diputado; Cerrillos un Diputado;
\textbf{Caldera un Diputado}; Rivadavia dos Diputados y un Senador; Iruya un Diputado; Orán dos
Diputados; Rosario de Lerma, dos Diputados}».
\textit{(B.O.\ Nº 795, h.\ 5 y 6, pág.\ 1.136.)}
\end{ficha}

\begin{cautela}
\textbf{Caldera figura dos veces en la misma enumeración, con la misma grafía y con un diputado
cada vez.} Está confirmado en la imagen y se transcribe como está. Este libro \textbf{no concilia}
y no infiere que la segunda mención deba decir otro departamento.

Y la lista es además \textbf{parcial}: no nombra a La Viña, Anta, Campo Santo, Rosario de la
Frontera, La Candelaria, San Carlos, Cachi ni Santa Victoria. \textbf{Un acto que convoca a toda
la provincia a votar enumera doce departamentos distintos en trece renglones}, repite uno y deja
sin nombrar al menos ocho.

No se saca de ahí ninguna conclusión sobre la representación del departamento: se registra que el
instrumento que la fija, ese año, no cierra consigo mismo.
\end{cautela}

\textbf{Las cuatro mesas son las mismas de 1918.} El mismo decreto ubica el Colegio Electoral
N.\ 3 de Caldera en cuatro puntos: la \textbf{Escuela Elemental Mixta} y la \textbf{Iglesia
Parroquial} del pueblo, la \textbf{Escuela Infantil Provincial} de Vaqueros y la \textbf{Oficina
de Correos} de Mojotoro. Son los cuatro edificios de 1918, dos años después y sin variación. Lo
que 1920 \textbf{no} publica son los circuitos: los treinta y dos parajes que 1918 enumeraba uno
por uno no vuelven a aparecer (lámina \ref{fig:parajesnom}).

\textbf{Tres partidos proclamaron candidato por el departamento}, y el tercero es un nombre
conocido: por la Unión Cívica Radical Intransigente, \textbf{José María Ocaranza}; por la Unión
Provincial, \textbf{José Antonio Aráoz}; y por la Unión Cívica Radical, \textbf{Dermidio
Arancibia}.

\begin{cautela}
\textbf{Dermidio Arancibia es, dos años antes, uno de los dos descubridores de las minas del
departamento.} El aviso de julio de 1918 lo nombra, junto con Miguel A.\ Sosa, como peticionante
de la mina de plomo «20 de Febrero» en la finca Cerro Nevado, y en agosto los dos denuncian el
despueble de la mina «Sara». En 1920 es candidato a diputado por el mismo departamento.

\textbf{Coincidencia de nombre, como siempre, y no identificación registral.} Pero si es la misma
persona, es el primer caso documentado en esta serie de alguien que pide una concesión sobre el
suelo del departamento y se presenta a representarlo, y el capítulo \ref{cap:politica} lo retoma
donde corresponde.
\end{cautela}

\textbf{Y a Augusto Regis lo sacan de la mesa antes de votar.} La nómina de la Junta Electoral de
febrero registra, para el Colegio Electoral N.\ 3 de Caldera, mesa N.\ 1: «\textit{Pte.,
Fernández Alberto Tomás, en reemplazo de Regis Augusto}»\index{Regis, Augusto}. El mismo Augusto
Regis que este libro sigue desde 1908 y que en 1918 era segundo suplente de esa misma mesa. El
acto no dice por qué.

\textbf{Y la urna del departamento se anuló.} El decreto N.\ 772 del 22 de marzo convoca a
elecciones complementarias para el 4 de abril «\textit{de las urnas que han sido anuladas}», y la
primera de la lista es «\textit{La Caldera --- C.\ Electoral N.º 3 Mesa N.º 4}».

\begin{cautela}
\textbf{La mesa 4 es la de Mojotoro}, según el decreto de febrero que ubica las cuatro. El decreto
772 no lo dice: nombra el colegio y el número de mesa, no el lugar, \textbf{y el cruce lo hace
este libro}. Tampoco dice por qué se anuló.

Puesto junto al resto del año, el cuadro es éste: la única plaza rentada de policía de partido del
departamento está en Mojotoro y cambia de manos tres veces en cinco meses; la única urna anulada
del departamento es la de Mojotoro; y ninguno de los dos actos explica nada. \textbf{Este libro no
afirma que haya relación entre las dos cosas}, y lo deja anotado porque el que siga puede
buscarla.
\end{cautela}

\textbf{Cierra el año un nombramiento militar.} Un decreto en acuerdo de Ministros de noviembre
nombra Coroneles de la Guardia Nacional por departamento; por Caldera, \textbf{Mariano Coll}. El
artículo 2.º nombra tenientes coroneles sólo por cinco departamentos, y Caldera no está entre
ellos.

\begin{cautela}
\textbf{Un año sin una sola cifra del departamento.} 1920 publica \textbf{diez resúmenes del
movimiento de la Tesorería General}, más que ningún otro de los años cuyos resúmenes este libro
contó uno por uno ---1911 publicó uno, 1912 cuatro, 1918 tres y 1919 ninguno; del tramo 1923--1946
no hay recuento---, y el propio Boletín los encadena mes a mes por la existencia de caja. \textbf{Y
no publica un solo balance municipal}, ni de este departamento ni de ningún otro. El Estado
provincial rinde su caja diez veces en un año; los municipios, ninguna (apéndice
\ref{cap:fuentes}).
\end{cautela}

\subsection*{Y 1921 devuelve los deslindes, y devuelve a los mismos linderos}

El año 1921 se leyó también edición por edición y es el undécimo tramo: las ediciones 840 a 888,
de las que seis son marcadores de ausencia del repositorio, de modo que son 43 ediciones y 742
hojas (apéndice \ref{cap:fuentes}). Dejó dieciocho actos del departamento, y \textbf{dos de ellos
son de dominio: los primeros desde 1917}, después de tres años seguidos en cero.

\begin{ficha}{Deslinde de la finca «El Angosto» --- B.O.\ Nº 869, 19 de agosto de 1921, aviso Nº 286}
Se presenta el procurador Elías Gallardo «\textit{con poder y títulos bastantes}» de la señora
\textbf{Antonia P.\ de Peralta} pidiendo el deslinde, mensura y amojonamiento de la finca «El
Angosto», departamento de La Caldera,

\begin{quote}
«\textit{limitada al Norte, propiedad de los herederos de \textbf{Luisa A.\ Costas}, hoy
\textbf{Daniel Linares} y de \textbf{Julio Mercado}; al Sud, el monte del \textbf{Quitilipe}
separado por el río de La Caldera y propiedad de \textbf{Susana Aguilera}; al Este, propiedad de
\textbf{Augusto Reyes}; al Oeste, herederos de \textbf{Luis A.\ Costas} \textit{[sic]} hoy Daniel
Linares}».
\end{quote}

Juez Alberto Mendioroz, agrimensor Rodolfo Chávez, secretario Tomás Izarrualde.
\end{ficha}

\begin{cautela}
\textbf{Los tres linderos del este, del sud y del norte son los mismos de 1909.} El remate de una
décima parte indivisa del \textbf{Angosto de Arrieta}, publicado en 1909, declara como linderos
«\textit{al este Julio Mercado, Augusto Rejis\index{Regis, Augusto} y Susana Aguilera}» (apéndice
\ref{cap:dominio}). \textbf{Doce años después, sobre el paraje contiguo, vuelven los tres
nombres}: Julio Mercado al norte, Susana Aguilera al sud y un Augusto ---impreso Reyes en 1921 y
Rejis en 1909--- al este.

\textbf{Este libro no unifica Reyes con Rejis ni afirma que las dos fincas sean la misma.} Lo que
queda establecido es que \textbf{la vecindad de propietarios del sector del Angosto no cambió en
doce años}, y que uno de ellos ---Julio Mercado--- es, en 1919 y en 1928, miembro de la Comisión
Municipal del departamento. Es la tercera tesis en su forma más literal: las unidades y los
nombres duran.

\textbf{Y aparece Quitilipe}, que el archivo registra en 1912 y otra vez en 1916, las dos en el
deslinde de «Santa Mónica y Quitilipe», y que aquí es el monte del sud, separado por el río. \textbf{Y aparece Daniel
Linares}\index{Linares, Daniel}, que en 1917 era miembro de la comisión municipal y candidato a
diputado, ahora como propietario en dos de los cuatro rumbos.

\textbf{Una errata del propio edicto que no se concilia}: el mismo texto escribe «Luisa A.
Costas» en el lindero del norte y «Luis A.\ Costas» en el del oeste. Dos personas o una sola mal
impresa.
\end{cautela}

\textbf{Y el segundo deslinde del año es el que este libro ya tenía, ahora con su edición.} El
edicto de «\textbf{Cerro de Buena Vista}» y «\textbf{Sausal o Curuzú}», partido de Vaqueros, de
Francisco S.\ Urquiza\index{Urquiza, Francisco S.}, sale en la edición 843 del 28 de enero. El
barrido agrega tres cosas: el juez es Daniel Etcheverry y el secretario Tomás Izarrualde; la
resolución se imprime «\textit{Salta, Octubre 18 de 1020}» \textit{[sic]}, con un cero donde la
fecha pide un nueve; y el auto ordena publicar el edicto \textbf{treinta días en dos diarios y
«\textit{por una sola vez}» en el Boletín Oficial}.

\begin{cautela}
\textbf{Esa instrucción vale más que el edicto.} En 1921 un juez de primera instancia distingue,
en el mismo auto, la publicación que da noticia ---treinta días en dos diarios de la capital---
de la que deja constancia ---una sola vez en el Boletín---. \textbf{El Boletín no es, en esta
época, el lugar donde se entera el interesado}: es el lugar donde queda el acto.

Cien años después, el capítulo \ref{cap:ribera} documenta líneas de ribera que salen dos días en
el Boletín y no se asientan en la matrícula, y el capítulo \ref{cap:opacidad} lo llama el canal:
la información existe, se publica y no viaja por donde mira el interesado. \textbf{En 1921 el
juez lo sabía y por eso mandaba las dos publicaciones.}
\end{cautela}

\subsection*{Un distrito cuya elección se anula entera}

De los dieciocho actos del año, nueve son electorales, y la serie termina mal.

\begin{ficha}{Decreto N.\ 1828 --- Salta, 20 de septiembre de 1921}
«\textit{Habiendo comunicado el señor Presidente de la Junta Electoral de la Provincia, en nota
N.o 30 de fecha 6 del corriente, que dicha Junta ha resuelto \textbf{declarar nula la elección
practicada el día 7 de Agosto ppdo., en los distritos eleectorales} \textit{[sic]} \textbf{de la
Caldera}, Rosario de Lerma, Rosario de la Frontera, La Candelaria, Guachipas, Iruya, Anta, Metán,
Santa Victoria, San Carlos, Orán y Rivadavia, \textbf{por no haberse verificado elección legal en
dos tercios de las mesas}}» ---artículos 75 y 84 de la Ley de Elecciones---, se convoca de nuevo
para el segundo domingo de octubre. A Caldera le corresponden \textbf{dos Convencionales para la
reforma de la Constitución}.
\textit{(B.O.\ Nº 879, 28 de octubre de 1921, h.\ 15 y 16.)}
\end{ficha}

\begin{cautela}
\textbf{Dos años seguidos con la elección del departamento anulada.} En 1920 se anuló una urna
---la de la mesa 4, que es la de Mojotoro--- y hubo complementarias. En 1921 se anula
\textbf{toda la elección del distrito}, y no por un incidente sino porque no hubo elección legal
en dos tercios de las mesas. El departamento no está solo: son doce distritos de los veintidós.

\textbf{Y la convocatoria a la nueva elección se publica el 28 de octubre, diecinueve días después
del segundo domingo de octubre, que cayó el 9.} El Boletín publica la citación cuando la fecha
señalada ya pasó. El decreto es del 20 de septiembre y la nota que lo funda, del 6: cincuenta y
dos días entre la resolución y su publicidad.

Este libro \textbf{no puede decir si la elección de octubre se hizo}, ni con qué convocatoria, ni
qué pasó con los dos convencionales del departamento en la reforma constitucional. Ni la nota N.\
30 de la Junta ni la resolución que declara la nulidad se publican, y los dos quedan pedidos en
el apéndice \ref{cap:pedidos}.
\end{cautela}

\textbf{Y el año trae cuatro nóminas de autoridades y de candidatos.} En febrero la Unión Cívica
Radical proclama por La Caldera al senador \textbf{Emeterio Royo} y al diputado \textbf{Lucio
Ortiz}; en diciembre, para el comicio de enero de 1922, la Unión Provincial proclama a
\textbf{José Antonio Aráoz} ---el mismo candidato que presentó en 1920--- y la Unión Cívica
Radical a \textbf{Natal López Cross}. Las autoridades de mesa de febrero y las de diciembre
traen veinticuatro nombres, y entre ellos vuelven \textbf{Domingo Royo}, \textbf{Daniel
Linares}\index{Linares, Daniel}, \textbf{Enrique J.\ Rauch}, \textbf{Gregorio Nogales},
\textbf{Nolasco Echenique} y \textbf{Eustaquio Murúa}\index{Murúa, Eustaquio}, que preside la
mesa 1 en diciembre.

\subsection*{Y Escolástico Aparicio vuelve, y dura dos meses}

El capítulo daba por cerrada en 1918 la serie de \textbf{Escolástico M.\
Aparicio}\index{Aparicio, Escolástico M.}, cuando la Intervención lo trasladó de la Receptoría de
Caldera a la de Campo Santo. \textbf{Vuelve en 1921}: el decreto 1669 del 7 de julio lo nombra
Juez de Paz Suplente del departamento, «\textit{vistas las ternas elevadas por la Comisión
Municipal del Departamento de la Caldera}», con \textbf{Laureano Peralta} de propietario. Y el
decreto 1816 del 12 de septiembre acepta su renuncia, «\textit{en atención a los motivos en que
la funda}», que el decreto no transcribe.

\begin{cautela}
\textbf{Dos meses y cinco días}, y el acto no nombra reemplazante: en lo que queda del año no se
publica otro que cubra esa vacante. \textbf{Y es la primera vacante sin destino de la serie}: las
tres que 1920 dejó eran del signo contrario ---declaraban el hueco sin nombrar a quien lo
dejaba---, y ésta nombra al que se va y no al que entra. El mismo año trae las dos formas: el
decreto 1669 que lo nombra tampoco dice quiénes eran los jueces salientes.

\textbf{Y el decreto que lo nombra tampoco nombra a los salientes} ni cita el acto por el que
cesaron: el visto habla de ternas para cubrir los cargos y no dice quién los dejaba. \textbf{Doce
años de cargos del departamento, y el archivo no permite armar la sucesión completa de ninguno de
ellos.}
\end{cautela}

\textbf{Y en noviembre un vecino pierde dos cargos en cuatro días.} El decreto 2021 del 7 declara
cesante al Comisario de Policía de La Caldera, \textbf{don Julio Royo}, y nombra «\textit{al
Ernesto Tedin}» \textit{[sic]}; el decreto 2062 del 11 lo separa del cargo de \textbf{Receptor de
Rentas y Expendedor de Guías} del departamento «\textit{por razones de mejor servicio}» y nombra
a \textbf{Ernesto J.\ Tedin}, con fianza de doña Encarnación Tedin de López por cinco mil pesos.

\begin{cautela}
\textbf{Los dos cargos del departamento que manejan dinero y fuerza pública pasan, en cuatro
días, del mismo hombre al mismo hombre}, y la fianza la da alguien del mismo apellido que el
entrante. El segundo decreto no cita al primero y ninguno de los dos da la causa: «razones de
mejor servicio» es la fórmula. Este libro no infiere nada sobre la conducta de nadie y registra
la operación.

\textbf{Y hay tres Royo en el departamento ese año}: Julio, el cesante; \textbf{Domingo}, miembro
de la Comisión Municipal en 1919 y presidente de la mesa 2 en febrero de 1921; y
\textbf{Emeterio}, candidato a senador por el departamento en la lista radical de febrero. No se
afirma parentesco.
\end{cautela}

\subsection*{Y por fin las escuelas tienen dotación}

El presupuesto de sueldos de las escuelas de la campaña, publicado en noviembre, trae el
departamento entero en dos renglones.

\begin{ficha}{Presupuesto de escuelas de la campaña --- B.O.\ Nº 883, 25 de noviembre de 1921, h.\ 10}
\textbf{Departamento Caldera.} Escuela \textbf{Elemental de 2.\textsuperscript{a}}, Parroquia:
director \$125 y dos maestros de grado a \$110 cada uno, \$345 mensuales y \$1.035 por trimestre.
Escuela \textbf{Infantil de 2.\textsuperscript{a}}, Vaqueros: director \$110 y un maestro \$100,
\$210 mensuales y \$630 por trimestre.
\end{ficha}

\begin{cautela}
\textbf{Es el primer dato de dotación escolar del departamento que este relevamiento encuentra}, y
dice tres cosas. Que el departamento tenía \textbf{dos escuelas} y cinco cargos docentes en total.
Que la del pueblo funcionaba en la \textbf{Parroquia}, que es también una de las cuatro mesas
electorales del departamento desde 1918. Y que \textbf{Vaqueros tenía la suya}, la Escuela
Infantil Provincial, que es la tercera mesa.

Los dos subtotales cierran: 125 más 220 son 345, y 110 más 100 son 210; y la segunda columna es
el trimestre, porque 345 por tres da 1.035 y 210 por tres da 630, igual que en los cuatro ítems
vecinos de la misma hoja.

El capítulo \ref{cap:ausencias} registra la escuela como una de las dimensiones que este libro no
cubre. \textbf{Sigue sin cubrirla}: esto es un renglón de presupuesto, no una historia de la
escuela. Pero con la jubilación de Juana Luque en 1918 y estos dos renglones de 1921, el
departamento deja de ser, en materia escolar, una página en blanco.
\end{cautela}

\section{La primera concesión de agua}

En enero de 1924 la Provincia otorgó una concesión de agua en el departamento. Los
considerandos explican, sin proponérselo, casi todo lo que este libro discute noventa años
después. \textbf{Y el barrido de 1923 aportó, al cierre, la pieza que falta antes del decreto:
el escrito con que el interesado la pidió}, publicado por edicto el 28 de diciembre de 1923.

\begin{ficha}{Pedido de concesión de agua --- B.O.\ Nº 990, 28 de diciembre de 1923, h.\ 13 y 14 (aviso 437)}
«\textit{Salta, Diciembre 10 de 1922.\ ---Francisco S.\ Urquiza}\index{Urquiza, Francisco S.},
con domicilio en Balcarce 415, expone que «\textit{según resulta del título que acompaño poseo
un bien raíz situado en el departamento de la Caldera, \textbf{partido de Vaqueros}, denominado
«Cerro de Buena Vista»}», y pide «\textit{hacer uso de los sobrantes de agua del \textbf{Río de
Wierna}}» en cantidad de \textbf{cien litros por segundo}, «\textit{siendo la extensión que me
propongo regar de \textbf{ciento noventa hectáreas}}».

Acompaña «\textit{un croquis explicativo del recorrido de la acequia}», que arranca del río y
sigue «\textit{por el curso actual de la acequia de la \textbf{finca Vaqueros}, al Oeste del
camino nacional de ésta a Jujuy}»; la única obra de arte necesaria es «\textit{un sifón que
atraviesa el cauce de un río seco que a veces recoge las corrientes veraniegas}». Y declara que
la acequia «\textit{atraviesa únicamente la finca nombrada de propiedad de \textbf{los herederos
Toranzos y doctor Carlos Serrey}}».

Pide, por último, «\textit{se recabe informes de la Municipalidad de la Caldera y se publique
edictos a fin de que los que se crean lesionados en sus derechos hagan valer su oposición en
forma}».

\textbf{Providencia del Ministerio de Gobierno, 12 de enero de 1923:} «\textit{de acuerdo a lo
que prescribe el Código Rural en su artículo 112, inciso 5º, publíquese durante treinta días en
el diario \textbf{La Provincia}}». Firma Ortelli.
\end{ficha}

\begin{cautela}
\textbf{Cinco cosas que el decreto de 1924 no dejaba ver.}

\textbf{Primera: la superficie.} Cien litros por segundo para \textbf{ciento noventa hectáreas}.
El decreto concede el caudal y no dice para cuánta tierra; el pedido sí. Es la primera
equivalencia entre agua y superficie que este libro tiene para el departamento.

\textbf{Segunda: el croquis existe y no se publicó.} El escrito dice acompañarlo y el Boletín
no lo trae ---igual que el plano de Hessling de 1911 y que los anexos que el capítulo
\ref{cap:opacidad} inventaría cien años después---. Citado y no acompañado.

\textbf{Tercera: la acequia atraviesa tierra ajena, y el dueño es un senador.} El propio escrito
declara que el canal cruza la finca Vaqueros «de propiedad de los herederos Toranzos y doctor
\textbf{Carlos Serrey}»\index{Serrey, Carlos}. Cinco meses después, en mayo de 1923, ese mismo
Carlos Serrey se incorpora al Senado provincial \textbf{como senador por Cerrillos}, en la misma
nómina que trae al senador por Caldera. \textbf{Coincidencia de nombre y no identificación
registral}; si es la misma persona, el propietario cuyo campo debe atravesar la acequia de la
primera concesión de agua del departamento legisla la provincia que la otorga. Y el apellido
\textbf{vuelve sobre la misma finca sesenta y dos años después}: la Ley 6334 de 1985 expropia
dos hectáreas de la Finca Vaqueros «\textit{de dominio de Manuel Serrey y otros}», y el capítulo
\ref{cap:expropiacion} registra tres expropiaciones a titulares de ese apellido en dos años.

\textbf{Cuarta: dónde mandó publicar la Provincia.} La providencia ordena publicar treinta días
\textbf{en el diario \textit{La Provincia}}, no en el Boletín Oficial ---donde sin embargo
apareció---. Es, en 1923, la misma distinción que un juez hacía en 1921 al mandar un edicto
treinta días a dos diarios y «por una sola vez» al Boletín: \textbf{la publicidad que da noticia
se hace en el diario; el Boletín deja constancia}.

\textbf{Y quinta, que es la más incómoda: el municipio al que el escrito pide informes estaba
intervenido.} El pedido reclama «\textit{se recabe informes de la Municipalidad de la Caldera}»,
y el decreto de 1924 consigna que «\textbf{la Comisión Municipal aconseja el despacho
favorable}». Entre una cosa y la otra, el 24 de septiembre de 1923, el Poder Ejecutivo
\textbf{intervino esa Comisión} y nombró un Interventor. \textbf{Este libro no puede establecer
si el dictamen favorable lo firmó la comisión o el Interventor}, porque el decreto no lo dice y
porque ningún acto del Interventor se publicó; y el apéndice \ref{cap:pedidos} lo pregunta. Lo
que queda establecido es que \textbf{el informe municipal que fundó la primera concesión de agua
del departamento se emitió mientras su gobierno estaba intervenido por la misma Provincia que
concedía}.

\textbf{Y una distancia de fechas que no se explica.} El escrito es del 10 de diciembre de 1922,
la providencia del 12 de enero de 1923 y el edicto del 28 de diciembre de 1923: \textbf{once
meses y medio entre la orden de publicar y la publicación}, y trece meses y medio entre el
pedido y el edicto. Ningún acto del año da la razón.

\textbf{Y un mes después del edicto, el decreto.} El 1.429 del \textbf{28 de enero de 1924}
---expediente \textbf{3898 letra C}--- concede el uso de los sobrantes, con el fundamento del
Departamento de Obras Públicas de que el río, «\textit{después de proveer a las fincas de
«Wierna» y «Vaqueros», deja correr un sobrante abundante}» (B.O.\ Nº 996, h.\ 1 y 2).
\textbf{Treinta y un días entre el último día del edicto y el decreto que concede}, para un
trámite que había tardado trece meses en publicarse.

\textbf{Y tres meses después, el concesionario entra al cuerpo municipal que informó su propia
concesión.} El decreto 1576 del 15 de abril de 1924 designa a Francisco Urquiza\index{Urquiza, Francisco S.}
miembro de la Comisión Municipal de La Caldera. La secuencia está entera en actos publicados y
este libro la registra sin afirmar relación entre sus términos; la desarrolla más abajo, en el
apartado de 1924.
\end{cautela}

\begin{ficha}{Decreto Nº 1429 del 28 de enero de 1924}
Expediente 3898. \textbf{Francisco S.\ Urquiza\index{Urquiza, Francisco S.}} solicita el uso de los sobrantes del
\textbf{río Wierna} para regar su propiedad \textbf{Cerro de Buena Vista}.

\textbf{El Departamento de Obras Públicas informa} «no poder afirmar si dicho río tiene o
no un caudal permanente como el que se pide, \textbf{en razón de no haberse efectuado
aforos de los ríos de la Provincia}»; y opina, en consecuencia, que la concesión sólo
podría otorgarse con carácter precario.

\textbf{La Comisión Municipal aconseja el despacho favorable}: el río, «después de proveer
de agua a las fincas de Wierna y Vaqueros, deja correr un sobrante abundante por su curso
natural»; las pérdidas se evitarían encauzándolo, y como desagua en el río Vaqueros y éste
en el Mojotoro, «siempre serían aprovechadas por otros campos cultivables». Agrega que la
finca «es apta para cualquier clase de cultivos» y que la concesión \textbf{«permitirá
fácilmente colonizar esas tierras»}.

Publicada la solicitud treinta días conforme el Código Rural, no hubo oposición. Se concede
el uso de los sobrantes, \textbf{cien litros por segundo}, «con carácter \textbf{precario y
eventual}, en previsión de la existencia posible de derechos legítimos adquiridos
anteriormente».
\end{ficha}

Cuatro cosas de este decreto siguen vigentes.

No había aforos. En 1924 el Estado provincial reconoce por escrito que no puede decir
cuánta agua lleva el río porque nunca lo midió, y por eso —no por la condición del
solicitante— otorga la concesión como precaria. La precariedad nace de la ignorancia del
Estado sobre su propio recurso. Noventa años después, el capítulo \ref{cap:ribera} documenta que el
procedimiento de línea de ribera pone en primer lugar de prelación los aforos «con
estadística no menor de cincuenta años», y que en la práctica se recurre al método
geológico-geomorfológico. La falta de mediciones sigue siendo el problema de fondo, y ahora
tiene fecha de inicio.

El municipio informa y aconseja, y su dictamen está publicado. En 2021, cuando se aprobó el
plano del loteo El Durazno, el municipio también intervino —a fojas 58, 84 y 97/101— y este
libro todavía no sabe qué dijo.

El propósito declarado es colonizar. No regar un cultivo existente: poblar. El agua se
concede para que la tierra pueda ocuparse. Es la formulación más temprana y más franca de
lo que este libro llama conversión de suelo, y viene del propio Estado.

Y la finca es Cerro de Buena Vista, la misma que Urquiza\index{Urquiza, Francisco S.} había hecho deslindar en 1921 y la
misma cuyo nombre reaparece en 2014 en una audiencia pública ambiental.

Esa concesión duró cuatro años. Vecinos y autoridades del departamento de Campo Santo, río
abajo, reclamaron contra ella; en abril de 1928 el gobierno rechazó el reclamo, y tres
semanas más tarde, con otro gobernador, la dejó sin efecto con una sola línea de
fundamento. Urquiza\index{Urquiza, Francisco S.} demandó a la Provincia y en 1931 la Corte de Justicia le dio la razón: ordenó
restablecer la concesión y dejó dicho que los motivos de interés público «deben aparecer
establecidos por formalidad y no meramente invocados». El capítulo \ref{cap:politica} reconstruye la
secuencia completa, porque lo que en ella se ve no es una cuestión de agua sino de quién
informa sobre qué, y de qué tiene que decir un acto para revocar un derecho.

\begin{cautela}
Guárdese el fundamento de esa precariedad, porque el capítulo \ref{cap:ribera} lo discute en su versión
contemporánea. En 1924 la concesión es precaria porque no se sabe cuánta agua hay. En 2014
y 2015, con el artículo 143 del Código de Aguas declarando eventuales todas las concesiones
subterráneas, dos gestiones para loteos —cuatrocientos lotes en uno, quinientos ochenta
habitantes en otro— se tramitaron como permanentes, y este libro no pudo establecer sobre
qué base.
\end{cautela}

\section{Minas, cateos y una fundición}

La actividad extractiva del departamento no empieza con las canteras de áridos de 2014. Es
tan vieja como el Boletín que la publica.

En diciembre de 1908, Odorico Díaz —en su nombre y en el de Emilio Fressart— se presenta
ante el Ministro de Hacienda exponiendo que «en el departamento de La Caldera, partido del
Potrero de Castillo» ha encontrado señas evidentes de minerales de cobre, y solicita
permiso de cateo. El trámite exigía presentar croquis del lugar y denunciar el domicilio de
los propietarios del terreno, y el pedido los nombra: «Los dueños del terreno son los señores
Garzon y Pinto, vecinos de Buenos Aires» (B.O.\ Nº 22, 30/12/1908, p.\ 106; edicto del 14 de
diciembre firmado por el escribano de minas Waldino Riarte). Pide cuatro unidades de quinientas
hectáreas, sobre la senda de las Ánimas y al sur del paralelo del cerro Nevado, porque los
terrenos «no son labrados ni cercados». La superficie no cierra con las medidas que el propio
pedido da ---quinientos metros por cuatro mil son doscientas hectáreas, no las «2ooo hectáreas»
que declara---, y el croquis no se publicó.

\textbf{Y la piedra empezó antes que el cobre.} En diciembre de 1911 un decreto crea una plaza
de agente para la comisaría de Vaqueros (B.O.\ Nº 305, 1911, h.\ 4) por dos motivos: la obra del puente sobre el río de ese
nombre y «\textit{la explotación de las canteras en las proximidades del puente sobre el río de
Mojotoro}». \textbf{Es la primera explotación de áridos que el archivo registra en el
departamento, ciento tres años antes de las concesiones de 2014} (capítulo \ref{cap:defensas}),
y aparece por el mismo motivo que hoy: hacía falta un policía donde se juntaban los
trabajadores.

Ese último requisito merece subrayarse, porque contrasta con todo lo que este libro
documenta un siglo después: para pedir cateo en 1908, la identificación del dueño del suelo
era condición de admisibilidad, no un dato que hubiera que deducir de actos ajenos.

\textbf{Pero la fiebre de 1925 tiene un antecedente de nueve años antes, y estaba en el mismo
expediente.} En noviembre de 1916 José María Romero Escobar pide al Ministro de Hacienda cuatro
unidades de cateo ---dos mil hectáreas--- por plomos argentíferos en el «\textit{Puesto de las
Pirguas}», «\textit{departamento de la Caldera}», en terrenos de los señores \textbf{Garzón y Pintos},
en las faldas del cerro «\textit{Cerro Nevado}» (B.O.\ Nº 633, h.\ 2). El informe del Escribano de
Minas, publicado a continuación, advierte que sobre esa misma finca pesan \textbf{dos pedimentos
anteriores}: uno de Casiano Hoyos, «\textit{por sí y por don Juan Larrán}»\index{Larrán, Juan}, en el
lugar «\textit{El Potrero}»; y otro de Octavio Espeche y Odorico Díaz, en la falda este del cerro
(h.\ 3).

\textbf{Ese nombre importa.} Juan Larrán\index{Larrán, Juan} es el único solicitante de 1925 cuyo
cateo el archivo registra concedido, y el que en 1930 vuelve con dos minas más. Aparece aquí
\textbf{nueve años antes}, asociado a otro solicitante y sobre los mismos terrenos de Garzón y Pintos
que el primer pedimento del departamento, el de 1908, ya nombraba. La continuidad de la propiedad del
suelo minero ---los mismos Garzón y Pintos en 1908, en 1916 y en 1925--- y la del solicitante quedan
documentadas, y la fiebre de 1925 deja de ser un comienzo.

\begin{cautela}
\textbf{Y el cerro no aparece en 1916 por primera vez.} El deslinde general de San Alejo y Santa
Rufina, de 1909, cierra su límite poniente «\textit{donde concluyen las propiedades en el cerro nevado,
su punta más alta}». Si es el mismo accidente, \textbf{el mojón poniente del deslinde más importante del
departamento es el cerro sobre el que siete años después se pide cateo}, y los dos hechos que este libro
trata en capítulos distintos recaen sobre el mismo punto. \textit{Si} es el mismo: en 1909 la expresión
puede ser descriptiva y no nombre propio, el edicto no lo decide y este libro tampoco.

El documento dice «\textit{distrito de las Pirguas}» y lo ubica en el departamento de la Caldera. El
nombre no figura en el vocabulario de topónimos del departamento y se asocia habitualmente a otra
jurisdicción. \textbf{Se transcribe lo que el escrito dice y no se lo concilia}: o el paraje existió
con ese nombre en el norte del departamento, o el escrito erró la jurisdicción. El expediente de
cateo lo diría; está pedido en el apéndice \ref{cap:pedidos}.
\end{cautela}

En el verano de 1925 hubo algo parecido a una fiebre. Sólo en enero se presentaron al menos seis pedimentos de cateo por minerales de primera categoría; tres de ellos, todos de dos mil hectáreas, publicaron el nombre del propietario del suelo:

\begin{ficha}{Pedimentos de cateo de enero de 1925}
\begin{center}\small
\begin{tabular}{@{}p{1.7cm}p{3.6cm}p{3.9cm}p{2.5cm}@{}}
\toprule
\textbf{Expte.} & \textbf{Presentantes} & \textbf{Finca} & \textbf{Propietario del suelo} \\
\midrule
1165 C & Roldán y Prémoli & San Alejo y Santa Rufina & Bernardo Austerlitz\index{Austerlitz, Bernardo} \\
1164 C & Povoli y Prémoli (h) & San Alejo y Santa Rufina & Bernardo Austerlitz\index{Austerlitz, Bernardo} \\
1168 C & Juan Larrán & \textbf{Las Nieves} & Garzón y Pintos; E.\ Quipildor \\
\bottomrule
\end{tabular}
\end{center}
Escribano de Minas: Zenón Arias. El pedimento de Larrán fue concedido el 1º de agosto de
1925.
\end{ficha}

Esos tres pedimentos, de dos mil hectáreas cada uno, suman seis mil, sobre un departamento de unas ciento nueve mil (capítulo \ref{cap:fincas}): sólo con ellos se solicitó explorar
cerca del cinco y medio por ciento del territorio. Y como el edicto exigía identificar a los propietarios, sabemos que San
Alejo y Santa Rufina —partido de 1909, finca del artículo 6º— pertenecía en 1925 a Bernardo
Austerlitz\index{Austerlitz, Bernardo}, domiciliado en Alberdi 374 de la capital, y que Las Nieves era de Garzón y
Pintos, con terrenos linderos de E. Quipildor. Para San Alejo es el primer nombre de propietario
que este libro obtiene \textbf{en propiedad plena}, porque el archivo de 1909 da el estado
anterior: las fincas San Alejo y Santa Rufina eran entonces un \textbf{condominio de la familia
Fernández}\index{Fernández, Ángel}, en juicio de división promovido por Benítez y Aguilar, con
deslinde general ordenado en mayo de 1909 a cargo del perito Walter Hesling (B.O.\ Nº 59, 1909, h.\ 4). \textbf{El
administrador de las fincas era Ángel Fernández}\index{Fernández, Ángel} ---el mismo que firmó
como Presidente Municipal el cuadro de rentas de 1908 y que ese año fue nombrado comisario
auxiliar del partido de San Alejo y Santa Rufina---, según lo reconoce él mismo en la rendición
de cuentas que el juzgado falló en noviembre de 1909 (B.O.\ Nº 114, 1909, h.\ 2). Dos de sus linderos al naciente eran
Daniel Linares\index{Linares, Daniel} y Abelardo Figueroa\index{Figueroa, Abelardo}, que ese año
integraban el jurado de patentes y la comisión municipal.

\textbf{El pleito terminó en remate, y el aviso da la primera superficie medida de una finca de
este departamento.} El 30 de diciembre de 1911, por orden del juez Vicente Arias y «de común
acuerdo de partes», se remató San Alejo y Santa Rufina (B.O.\ Nº 306, 1911, h.\ 4), «\textit{perteneciente a los herederos
Fernández}»\index{Fernández, Ángel}, con \textbf{17.874 hectáreas 27 áreas y 79 cuadrados según
mensura practicada por el agrimensor señor Walter Hessling y aprobado por el Departamento
Topográfico}, y base de \textbf{ciento ochenta mil pesos al contado}. Son unas dieciséis
hectáreas de cada cien del departamento en un solo remate. Los límites vuelven a nombrar a los
mismos: al este, la finca Angostura de Daniel Linares\index{Linares, Daniel}, con Eusebio
Herrera, Francisco Sajama, Arturo Escudero y Abelardo Figueroa\index{Figueroa, Abelardo}; al
sud, otra vez Linares, los herederos del doctor Desiderio Ruiz, \textbf{Garzón y Pintos} y
Silverio Bejarano; al norte, la provincia de Jujuy, la cabaña de Marcos Zorrilla, Uracatao o
Pueblo Viejo de Domingo Royo y Gavino Sánchez; al oeste, José Bruzzo y otra vez Zorrilla.

Dos comparaciones ordenan la escala. La base del remate ---\$180.000--- es \textbf{ciento veintiún
veces la renta municipal de todo el año 1908}. Y el honorario que el mismo juzgado aprobó en
junio de 1911 al agrimensor por ese deslinde, \textbf{\$5.000} (B.O.\ Nº 258, 1911, h.\ 1), es más del triple de esa renta
anual. \textbf{Medir una finca del departamento costaba tres años y medio de municipio.}

\textbf{No se vendió.} En abril de 1912 el mismo martillero la vuelve a ofrecer (B.O.\ Nº 328, 1912, h.\ 4), ahora bajo el
título «\textit{La finca más importante del Depto.\ de la Caldera}», y \textbf{con base de ciento
treinta y cinco mil pesos}: un cuarto menos que cuatro meses antes. El aviso agrega lo que el
remate anterior no decía y es lo más parecido a una descripción del territorio que este libro
encuentra para la época: «\textit{bosques variados con riquísimas maderas}», «\textit{pasto
suficiente para sostener 8.000 animales vacunos}», «\textit{agua en abundancia y clara que nace
en la misma propiedad}» y \textbf{una Escuela Nacional} dentro de la finca. Esa escuela es
veinticinco años anterior a los trámites de la Ley Láinez que el capítulo \ref{cap:presencia}
documenta para los parajes del departamento.

\textbf{Y tampoco se vendió en 1912.} El barrido edición por edición de 1914, incorporado al
cierre, muestra que la misma finca salió a remate \textbf{tres veces más en ocho meses}, con el
mismo martillero, el mismo juez y el mismo juicio, y con la base cayendo en cada llamado.

\begin{ficha}{Los tres remates de 1914}
Los tres avisos repiten el título ---«\textit{¡¡Extraordinario remate!! --- La finca más
importante del Departamento de La Caldera}»--- y casi palabra por palabra el mismo cuerpo, por
orden del juez de primera instancia \textbf{Vicente Arias}, en el juicio de división de
condominio de los herederos Fernández\index{Fernández, Ángel} y «\textit{de común acuerdo de
todos sus propietarios actuales}». Martillero: \textbf{M.\ R.\ Alvarado}.
\begin{itemize}[leftmargin=*,nosep]
\item \textbf{21 de febrero}, en el Jockey Club: base \textbf{\$5,50 la hectárea} (B.O.\ Nº 463, h.\ 3).
\item \textbf{Martes 12 de mayo}, en el escritorio del martillero, España 530: base \textbf{\$4,25} (B.O.\ Nº 476, h.\ 4).
\item \textbf{Lunes 5 de octubre}: base \textbf{\$3,20} (B.O.\ Nº 509, h.\ 3).
\end{itemize}
El de octubre agrega lo que los anteriores no traen: \textbf{canteras en explotación},
yacimientos de cal y de arcilla «\textit{cuyas muestras obtuvieron premio en la exposición de
Buenos Aires de 1911}», y otra vez la escuela nacional. Y publica las condiciones de pago:
diez por ciento en el acto, veinticinco a la firma de la escritura dentro de los treinta días,
y el saldo del sesenta y cinco en dos cuotas a seis y doce meses, al siete por ciento anual y
con garantía hipotecaria sobre la misma propiedad; comisión del dos por ciento a cargo del
comprador.
\end{ficha}

\textbf{Puestas en serie, las cinco bases son el primer precio de la tierra que este libro
puede seguir en el tiempo.}

\begin{center}\small
\begin{tabular}{@{}lll@{}}
\toprule
\textbf{Remate} & \textbf{Base publicada} & \textbf{Equivalente} \\
\midrule
30 dic.\ 1911 & \$180.000 el total & \$10,07 la hectárea \\
abr.\ 1912 & \$135.000 el total & \$7,55 la hectárea \\
21 feb.\ 1914 & \$5,50 la hectárea & \$98.308 el total \\
12 may.\ 1914 & \$4,25 la hectárea & \$75.966 el total \\
5 oct.\ 1914 & \$3,20 la hectárea & \$57.198 el total \\
\bottomrule
\end{tabular}
\end{center}

\textit{La columna de equivalencias es de este libro y no de los avisos: se calcula sobre las
17.874,2719 hectáreas que los cinco declaran. Los de 1911 y 1912 publican sólo el total; los
tres de 1914, sólo el precio por hectárea.}

\textbf{De diciembre de 1911 a octubre de 1914 la base cayó un sesenta y ocho por ciento}, y
dentro de 1914 un cuarenta y dos por ciento en ocho meses. Ningún aviso dice por qué y este
libro no lo sabe. Lo que queda documentado es que \textbf{la finca más grande del departamento
se ofreció cinco veces en tres años y el precio de salida se redujo a un tercio}.

\begin{cautela}
El aviso de 1912 no coincide consigo mismo ni con el de 1911. La superficie se imprime
«17,874 hectáreas 27 áreas 79 mts.\ cdados.» en el título y «17874 hectáreas, 27 áreas, 19 metros
cuadrados» doce renglones más abajo. \textbf{Los tres avisos de 1914 imprimen las tres veces «27 áreas y 19 metros cuadrados»}, y el de febrero lo imprime además en cuerpo grande: la forma que se repite es ésa. Y \textbf{Garzón y Pintos}, que en diciembre de 1911
lindaban al sud, en abril de 1912 aparecen al oeste. Este libro transcribe las dos versiones y no
elige: para la ubicación de esa propiedad hace falta el plano de Hessling, que ninguno de los dos
avisos publica y que el apéndice \ref{cap:pedidos} pide.

\textbf{Con el apellido del lindero del naciente pasa lo mismo, y 1914 lo resuelve.} En esa
misma posición de la enumeración, los avisos imprimen «Sajama» en 1911, «\textit{Sosayo}» en
mayo de 1914 y «\textit{Sosavo}» en octubre, y un edicto de 1932 imprimirá «\textit{Sosauya}»
(capítulo \ref{cap:fincas}). Pero en septiembre de 1914 el interesado se presenta él mismo a
pedir el deslinde de su propia finca, y el edicto lo nombra \textbf{Francisco Sosaya}. Ésa es
la forma que este libro usa; las otras quedan anotadas y no se corrigen en las citas.
\end{cautela} Para la zona de Las Nieves no: \textbf{los mismos apellidos ---«Garzon y
Pinto, vecinos de Buenos Aires»--- son los dueños que declaró el cateo de 1908} en el partido del
Potrero de Castillo. Que se trate del mismo inmueble es probable y no está probado: el partido era
más grande que la finca. Pero si lo es, la tierra que hoy es el catastro 102, fiscal, y que la
comunidad Kondorwaira reclama (capítulo \ref{cap:tierra}) tuvo durante al menos diecisiete años
dueños privados que vivían en Buenos Aires.

La finca Las Nieves da nombre, además, al río sobre el cual en 2015 se concesionó una
cantera de áridos. El edicto de 1925 menciona como punto de referencia el «Abra de las
Minas»: el topónimo ya registraba actividad anterior.

\begin{cautela}
Un detalle de esos expedientes merece anotarse, y cinco años después el archivo lo amplía.

En marzo de 1930, la Dirección General de Minas declaró caducas por abandono de trámite
\textbf{ocho solicitudes de cateo} del departamento, todas presentadas en 1925. Sus
solicitantes fueron:

\begin{center}\small
\begin{tabular}{@{}llll@{}}
\toprule
\textbf{Expte.} & \textbf{Solicitantes} & \textbf{Presentada} & \textbf{Mineral} \\
\midrule
1163 C & Dorindo F.\ \textbf{Prémoli} y Ramón E.\ Juárez & 5 ene.\ 1925 & 1ª categoría \\
1164 C & Rodolfo Povoli y Enrique \textbf{Prémoli} (h) & 5 ene.\ 1925 & 1ª categoría \\
1165 C & José Manuel Roldán y Apolo A.\ \textbf{Prémoli} & 5 ene.\ 1925 & 1ª categoría \\
1167 C & Ignacio \textbf{Prémoli} y José María Oñativia & 15 ene.\ 1925 & 1ª categoría \\
1169 C & Emilio D.\ Sylvester & 21 ene.\ 1925 & 1ª categoría \\
1173 C & Federico Gauffín & 16 feb.\ 1925 & carbón de piedra \\
1195 C & Julio Suárez & 21 ago.\ 1925 & auríferos y cobre \\
1198 C & Francisco G.\ de la Vega & 18 set.\ 1925 & oro, plata y cobre \\
\bottomrule
\end{tabular}
\end{center}

\textbf{Cuatro de los ocho expedientes llevan el apellido Prémoli}, y son cuatro personas
distintas —Dorindo, Enrique hijo, Apolo e Ignacio—, cada una asociada a un socio diferente,
presentando cuatro solicitudes separadas en el término de diez días. Dos de ellas ---la 1164 C y la 1165 C, ya consignadas en la ficha de enero--- recaen sobre la misma finca, el mismo propietario y el mismo punto de arranque, con
el rectángulo girado noventa grados.

Este libro no sabe si eso respondía a una estrategia para cubrir más superficie que la
permitida por solicitante, a una división de tareas dentro de una familia o a una
concurrencia simple. El Código de Minería de la época fijaba límites por peticionante, y ése
es el punto a verificar: si el límite existía y era de dos mil hectáreas, cuatro parientes
presentando por separado cubrían ocho mil.

Y el desenlace da la medida de lo que fue aquella fiebre. Las ocho solicitudes se declararon
abandonadas por sus propios titulares ---según el decreto, entre abril y septiembre de \textbf{1925}, es decir a las pocas semanas o meses de presentarse; una de las solicitudes es del 18 de septiembre, de modo que la fecha debe cotejarse--- y la
caducidad se declaró formalmente cinco años más tarde. Nadie exploró: se pidió, se dejó
caer, y el registro tardó un lustro en anotarlo.

Se consigna porque es el tipo de operación que este libro busca en el presente y casi nunca
puede ver: varios actos que por separado son rutinarios y juntos dibujan otra cosa. En 1925
se pueden leer juntos porque todos se publicaron con el nombre del solicitante, el del
propietario del suelo, la superficie y el punto de arranque.
\end{cautela}

\subsection*{Y en 1930 aparecieron el plomo y la plata}

La fiebre de 1925 no dejó nada, pero cinco años después alguien sí llegó a manifestar un
descubrimiento.

En 1930 ---el padrón minero registra la manifestación el 30 de abril---, \textbf{Juan Larrán}, el único solicitante de 1925 cuyo cateo concedido registra el archivo ---ajeno a
los ocho caducados, y sobre la finca Las Nieves---, presenta el descubrimiento de una
mina de \textbf{plomo y plata} de tres pertenencias, denominadas Caldera Nº 1, Nº 2 y Nº 3, sobre la finca \textbf{San José}.

\textbf{Ese mismo año se le concedió otro cateo.} El 15 de octubre de 1930 la Dirección General de Minas otorgó a \textbf{Alberto Jaime Harrison y Juan Larrán} un permiso de exploración de dos mil hectáreas, para minerales de primera categoría, sobre la finca \textbf{Las Nieves}, «Departamento Caldera», en terrenos de \textbf{Julio Suárez y Francisco de la Vega}. El pedido era del 1º de junio de 1929 (expediente 38-H; Boletín Oficial Nº 1.346). Pero el padrón minero de 1933 fecha el registro de la manifestación de las minas (expediente 44-L) el \textbf{30 de abril de 1930}: el descubrimiento se registró mientras este segundo pedido estaba en trámite y seis meses antes de que se concediera. Qué relación hubo entre ambos, y con el cateo de 1925, el archivo relevado no lo dice; el apéndice \ref{cap:pedidos} pide los expedientes. El edicto muestra, además, que Las Nieves, de Garzón y Pintos en 1925, figuraba en 1930 a nombre de Suárez y de la Vega.

El punto de partida que fija el edicto merece transcribirse, porque describe cómo se median
las cosas: «el vértice Noroeste de la casa de \textbf{los hermanos Quipildor}, ubicado a
quinientos metros al Oeste verdadero del \textbf{Cerro de la Inga}». Una casa de familia
como origen de coordenadas de una concesión minera.

Los propietarios del suelo son los mismos \textbf{Quipildor Hermanos}, junto a \textbf{Julio
Suárez} y \textbf{Francisco de la Vega}. Los dos últimos figuran también como solicitantes de
cateo en 1925, de modo que en 1930 son a la vez propietarios superficiarios y peticionantes
mineros en el mismo departamento.

El padrón minero provincial cerrado el 30 de marzo de 1933, publicado en el Boletín Oficial del 30 de junio, registra esas minas bajo el número de orden 93 ---expediente 44-L, inscripto el 30 de abril de 1930---, a nombre de Juan Larrán: veintisiete hectáreas, tres pertenencias,
distrito San José, plomo y plata, canon al día. Y en la columna de estado del expediente
consigna dos palabras: \textbf{«a mensurarse»}.

Tres años después de manifestado el descubrimiento, la única mina registrada del departamento
todavía no había sido medida. Es el mismo desenlace que tuvieron las ocho solicitudes de
1925, sólo que más lento: se pide, se registra, se paga el canon y no se mide.

Vale detenerse en el mineral. Es \textbf{plomo}, el mismo que fundía el establecimiento de
Mojotoro entre 1917 y 1925 con mineral traído desde Jujuy. Cuando esa fundición salió a remate,
a fines de 1925, el departamento no tenía plomo declarado propio. Cinco años más tarde
lo tuvo, y la fundición ya no existía.

Este libro no sabe si ambas cosas están conectadas —si el cateo buscaba abastecer una
industria que se había liquidado, o si el hallazgo fue posterior e independiente—, y el
archivo no lo dice. Registra la secuencia: una fundición de plomo sin mineral local, y cinco
años después mineral local sin fundición.

\subsection*{Y en 1941 se remató}

La única mina registrada del departamento tuvo un final, y el decreto que lo ordena cierra el
hilo entero.

\begin{ficha}{Decreto Nº 4769 del 14 de mayo de 1941, y edicto de remate (Boletín Oficial Nº 1.899, 30 de mayo de 1941)}
El Ministerio de Hacienda, Obras Públicas y Fomento dispone el remate administrativo de
varias minas cuyas concesiones caducaron \textbf{por falta de pago del canon minero},
conforme el artículo 7º de la Ley 10.273. Martillero designado: Francisco Peñalba Herrera.
Firman Cornejo y Ricardo E.\ Usandivaras. El decreto comprende siete minas; entre ellas, \textbf{«La Colorada»}, de cobre, en la quebrada de Cobres (La Poma), inscripta a nombre de Alejandro Bonari \textbf{y de Juan Larrán}, con base de \$1.800.

Entre las rematadas, expediente 44-L: \textbf{tres pertenencias, veintisiete hectáreas, minas
de plomo y plata denominadas «Caldera Nº 1, 2 y 3», en la finca San José, departamento La
Caldera, inscriptas a nombre de Juan Larrán}. Base de venta: \$2.500.

Remate: sábado 17 de julio de 1941 ---así lo dice dos veces el edicto publicado en el Boletín Oficial Nº 1.899 del 30 de mayo, aunque el 17 de julio de 1941 fue jueves: el error es del original---, en las oficinas de la Dirección General de Minas, calle
Córdoba 56. Seña del veinte por ciento en el acto.
\end{ficha}

La secuencia completa de la minería de este departamento cabe en cinco fechas y no produjo
nada.

En \textbf{1908} se pide cateo de cobre en el partido de Potrero de Castillo. En \textbf{1925}
se presentan al menos nueve solicitudes ---ocho que caducarán, y la de Larrán; al menos tres de ellas de dos mil hectáreas cada una; cuatro de ellas a nombre de
cuatro Prémoli distintos--- y \textbf{las ocho se abandonan a los pocos meses}; la caducidad se
declara cinco años después. En \textbf{1930}, el único solicitante de 1925 cuyo cateo
prosperó según el archivo ---Larrán, ajeno a los ocho caducados--- manifiesta el descubrimiento de plomo y plata en la finca San José. En \textbf{1933}
el padrón minero provincial registra esa mina con el canon al día y el expediente «a
mensurarse». Y en \textbf{1941} se la remata porque dejó de pagarse el canon.

\textbf{Y el propio 1941 registra el paso intermedio}, que el barrido de ese año incorporó al
cierre: un acto posterior al remate, en aplicación del artículo 7º de la Ley 10.273,
\textbf{«\textit{dispone se inscriba la presente mina “Caldera nº 1, 2 y 3”, expediente
nº 44-letra L, en calidad de \textbf{vacante}}»} en el libro respectivo.

Y en agosto de \textbf{1943} la Dirección General de Rentas pidió la anulación de las patentes
de esas minas, «en mérito de encontrarse caducas por haber sido rematadas sin haberse
obtenido su venta». Es decir: el remate de 1941 \textbf{quedó desierto}. Nadie las compró, y
el Estado borró la patente de sus registros.

\textbf{De modo que el final tiene tres actos y no dos}: se remata en mayo, \textbf{se inscribe
como vacante} en el mismo 1941 y se le anula la patente en 1943. \textbf{Treinta y tres años
desde el primer cateo de 1908 hasta la palabra «vacante»}, y ni una tonelada de mineral en el
medio.

\textbf{Treinta y cinco años, once expedientes ---doce, con el cateo de 1929 documentado después--- y ni una tonelada documentada.} La única mina
que llegó a inscribirse nunca fue mensurada, y su titular la perdió por deuda con el fisco.

Conviene retener eso al leer el capítulo \ref{cap:resistencias}, que documenta seis concesiones de áridos entre 2014 y 2015 --- tres de ellas sobre tierra fiscal --- y una manifestación de descubrimiento de
fosfatos de mil doscientas cincuenta y seis hectáreas en la Sierra de Mojotoro. La actividad
extractiva de este departamento tiene un siglo de intentos y muy poca producción verificable;
lo que sí produjo, de manera sostenida, fueron \textbf{expedientes sobre superficie}.

\subsection*{La fundición de Mojotoro}

El departamento tuvo, además, industria pesada.

En octubre de 1925 el juez Carlos Gómez Rincón declara disuelta la sociedad «Levilly y
Compañía», constituida por escritura pública del 6 de noviembre de 1917 en la ciudad de
Buenos Aires entre Francisco Sabatie y Emilio Oliver Levilly, «como socios para la
explotación de la mina de plomo denominada Pulpera, ubicada en la Provincia de Jujuy, y el
establecimiento de fundición situado en el Departamento de la Caldera de esta provincia
denominado Mojotoro». Tres semanas después, el martillero José María Leguizamón anuncia el
remate de ambos bienes juntos con base de \$222.144,5.

No era un obraje ni una cantera: era un establecimiento metalúrgico, alimentado con mineral
traído de otra provincia. El capital era porteño —sociedad constituida en Buenos Aires,
liquidada por un pleito entre firmas de afuera—. Y la escala está fuera de todo lo demás
que el departamento registra en el período: el presupuesto municipal de 1908 era de \$1.483
y la finca Potrero del Gallinato salió a remate en 1922 con base de \$800. El activo industrial
valía más de cien veces una finca entera.

Esto corrige una impresión que este libro venía sosteniendo. Veníamos describiendo una
economía que pasa de la agricultura al suelo urbano, con la extracción de áridos como
actividad intermedia. La fundición muestra otra capa, anterior y de naturaleza distinta,
que llegó de afuera, duró unos ocho años y terminó liquidada en un remate judicial.

La historia económica del corredor no es una línea recta. Es una sucesión de actividades
que llegan de afuera, duran una década y se liquidan.

Conviene además dejar asentado cómo apareció ese remate, porque mide el alcance de este
relevamiento. El aviso no dice «Caldera» en ninguna parte: se encontró buscando «Mojotoro».
Y el escaneo imprimió «Eevilly» en lugar de «Levilly», de modo que buscar por el nombre de
la sociedad tampoco lo habría hallado. Un activo de doscientos veintidós mil pesos del
departamento se publicó sin nombrar al departamento y con el nombre del titular mal
impreso. Cuántas piezas equivalentes quedaron fuera de este trabajo es algo que no se puede
estimar.

\section{Una comisión municipal designada, un presupuesto aprobado desde afuera}

El primer semestre de 1922 se leyó edición por edición y es el duodécimo tramo, el único que
cubre medio año y no un año: las ediciones \textbf{890 a 916}, veintisiete ediciones y 446 hojas,
que van del 13 de enero al 21 de julio. \textbf{De la 917 en adelante no hay archivo}, de modo que
el segundo semestre del año no se leyó (apéndice \ref{cap:fuentes}). Con esa reserva ---que es
grande y vale para todo lo que sigue---, el tramo deja \textbf{veintitrés actos del departamento},
ninguno de dominio salvo un aviso de remate, y \textbf{el año en que el departamento recupera
órgano municipal}.

El decreto 510, del 18 de abril, considera «\textit{necesario designar las Comisiones Municipales
en los departamentos y distritos de la Provincia que en la actualidad carecen de ella}», y
el Interventor Nacional nombra miembros de la de La Caldera a Daniel Linares\index{Linares, Daniel}, Augusto
Regis\index{Regis, Augusto}, Domingo Royo, Jorge Rauch y Manuel Condori (B.O.\ Nº 906, h.\ 3). El decreto 172, del 15 de julio, aprueba
el presupuesto que esa comisión sancionó en sesión del 28 de mayo, previo dictamen del
Fiscal General (Nº 916, h.\ 4).

\begin{cautela}
\textbf{Cuatro palabras del decreto 510 cambian lo que este libro venía diciendo, y conviene
leerlas despacio: «que en la actualidad carecen de ella».}

Este capítulo describió, año por año, un cuerpo municipal que se renueva, que eleva ternas de
jueces de paz y que pierde y repone miembros: en 1909, en 1910, en 1915, en 1916, en 1917 y en
1919. \textbf{En abril de 1922 el departamento no tiene Comisión Municipal}, y el acto que la crea
lo dice en su propio fundamento. El archivo leído no registra cuándo dejó de tenerla ---el último
acto sobre ese cuerpo es el decreto 1669 del 7 de julio de 1921, que invoca «\textit{las ternas
elevadas por la Comisión Municipal del Departamento de la Caldera}», de modo que en julio de 1921
todavía existía---, ni si en el intervalo hubo comisionado municipal designado, como el que el
mismo decreto 510 da por terminado para Molinos en su artículo 1º.

\textbf{No se infiere que no lo haya habido}: el tramo empieza el 13 de enero de 1922 y el
segundo semestre de 1921 se leyó entero sin hallar el acto. Lo que queda establecido es más
acotado y basta: \textbf{entre algún momento posterior a julio de 1921 y el 18 de abril de 1922 el
departamento no tuvo el cuerpo que lo administraba}, y quien lo repuso fue un Interventor
Nacional, no una elección.
\end{cautela}

\subsection*{Y en junio un decreto rectifica ese acto y no cambia nada que se vea}

Dos meses después, ya con la Intervención terminada y el gobierno del doctor Adolfo Güemes en
funciones, un decreto provincial vuelve sobre el mismo artículo.

\begin{ficha}{Decreto del Poder Ejecutivo de la Provincia --- B.O.\ Nº 910, h.\ 9, junio de 1922}
«\textit{[\ldots] segundo del decreto N 510, de fecha 18 de Abril ppdo.\ en cuanto nombra miembros
de la Comisión Municipal de La Caldera a los señores \textbf{Augusto Regis, Daniel Linares,
Domingo Royo, Manuel Condori y Jorge A.\ Rauch}}».

\textit{El número propio de este decreto no está impreso en la parte del acto que el corpus trae.}
\end{ficha}

\begin{cautela}
\textbf{El decreto dice que rectifica y repite los mismos cinco nombres.} Las únicas diferencias
visibles entre el artículo 2º del decreto 510 y su rectificación son \textbf{el orden de la
enumeración} ---Regis\index{Regis, Augusto} pasa del segundo lugar al primero, Linares\index{Linares, Daniel} del primero al segundo, Rauch del
cuarto al quinto--- y \textbf{una inicial}: «Jorge A.\ Rauch» donde el 510 dice «Jorge Rauch».

\textbf{Qué se rectifica no está en la parte legible del acto}, y este libro no lo infiere. Lo
registra porque es, en su forma más pura, el mecanismo que el capítulo \ref{cap:opacidad}
describe cien años después: \textbf{un acto publicado que declara corregir otro sin decir qué
corrige}. El lector del Boletín de 1922 tenía delante los dos textos y no podía saber cuál era la
diferencia; el de 2026 tampoco.
\end{cautela}

\subsection*{Cuatro cargos, cuatro cadenas, y ningún acto que diga por qué}

El tramo se ordena en cuatro hilos, y los cuatro tienen la misma forma: el acto nombra o
declara cesante, no cita al anterior y no da la causa.

\begin{ficha}{Los cuatro cargos del departamento entre noviembre de 1921 y julio de 1922}
\textbf{Comisaría de policía del departamento.} \textbf{Decreto 32} del 24 de noviembre de 1921:
acepta la renuncia de \textbf{Ernesto J.\ Tedin} (Nº 890, h.\ 17 y 18). \textbf{Decreto 40} del
mismo día: nombra a \textbf{Julio Royo} «\textit{habiéndose aceptado en la fecha la renuncia
presentada por el señor Comisario del Departamento}», sin nombrarlo (h.\ 19). \textbf{Decreto 349}
del 8 de marzo de 1922: declara cesante a Royo y nombra a \textbf{Jorge A.\ Rauch} (Nº 899,
h.\ 10). Tres titulares en tres meses y medio.

\textbf{Subcomisarías.} \textbf{Decreto 176} del 2 de enero: declara cesante al subcomisario de
Mojotoro \textbf{Medardo Sarmiento} «\textit{en vista del continuo abandono que hace de su
puesto}» y nombra a \textbf{Rodolfo Pérez} (Nº 895, h.\ 3). \textbf{Decreto 399} del 21 de marzo:
nombra a \textbf{Juan de la Cruz Nogales} subcomisario de Mojotoro y a \textbf{Eustaquio Guerra}
subcomisario \textit{ad honorem} de \textbf{Vaqueros} (Nº 902, h.\ 1). Junio: \textbf{Víctor
Masclef}, subcomisario \textit{ad honorem} «\textit{de Mojotoro, La Caldera}» (Nº 912, h.\ 1).
Cuatro titulares para Mojotoro en seis meses.

\textbf{Juzgado de paz.} \textbf{Decreto 123} del 17 de diciembre de 1921: acepta la renuncia de
\textbf{Laureano Peralta} al juzgado de paz propietario (Nº 894, h.\ 1). \textbf{Decreto 105} del
12 de junio de 1922: «\textit{vista la terna elevada por la Comisión Municipal de La Caldera}»,
nombra propietario a \textbf{Carlos F.\ Matienzo}\index{Matienzo, Carlos} y suplente a \textbf{Vicente Vivas} (Nº 912,
h.\ 3). \textbf{Seis meses sin juez de paz propietario publicado.}

\textbf{Registro Civil.} \textbf{Decreto 402} del 22 de marzo: declara cesante a \textbf{Casimiro
Cáceres} y nombra a \textbf{Julio Royo Ortiz} (Nº 902, h.\ 2). \textbf{Decreto 524} del 20 de
abril: deja sin efecto el nombramiento de «\textit{D.\ Julio Royo}» \textbf{por no haber tomado
posesión del cargo} y nombra a \textbf{Santiago Guzmán}, expediente 4054 letra E (Nº 906, h.\ 8).
\textbf{Decreto 135} del 27 de junio: nombra a \textbf{Carlos F.\ Matienzo}\index{Matienzo, Carlos} (Nº 914, h.\ 4).
\textbf{Cuatro encargados en cuatro meses}, y ningún acto dice qué pasó con Guzmán.
\end{ficha}

\begin{cautela}
\textbf{Tres cosas se leen ahí, y la tercera es la que le importa a este libro.}

\textbf{La primera: un solo relevo del año da su motivo.} De los once actos de la ficha, el único
que dice por qué es el decreto 176, y lo que dice es «\textit{el continuo abandono que hace de su
puesto}». Los otros diez usan la fórmula o no usan ninguna. \textbf{Y ese acto no contiene la
palabra Caldera}: se atribuye al departamento porque el decreto de junio escribe «\textit{de
Mojotoro, La Caldera}». Una búsqueda por el nombre del departamento no lo encuentra.

\textbf{La segunda: los cargos se acumulan.} \textbf{Carlos F.\ Matienzo}\index{Matienzo, Carlos} es, con quince días de
diferencia, juez de paz propietario del departamento y encargado de su Registro Civil. Es el
mismo apellido del \textbf{Carlos Matienzo}\index{Matienzo, Carlos} que en 1909 era juez de paz propietario y firmaba
además como secretario del cuadro de rentas municipales. \textbf{Coincidencia de nombre, como
siempre, y no identificación registral}; si es la misma persona, son trece años y la misma
combinación de cargos. Y \textbf{Jorge Rauch} es comisario de policía desde marzo y miembro de la
Comisión Municipal desde abril: el que manda la policía integra otra vez el cuerpo que gobierna
el municipio, como Aparicio\index{Aparicio, Escolástico M.} en 1910.

\textbf{Y la tercera: el apellido Royo aparece cuatro veces en cinco meses, en cuatro cargos
distintos y con cuatro formas del nombre.} \textbf{Julio Royo}, comisario; \textbf{Julio Royo
flor} ---así impreso, con efe minúscula---, Receptor de Rentas y Expendedor de Guías con fianza
de cinco mil pesos (Nº 900, h.\ 10); \textbf{Julio Royo Ortiz}, encargado del Registro Civil; y
\textbf{Domingo Royo}, miembro de la Comisión Municipal. A eso se suma que el decreto 524 deja
sin efecto el nombramiento de «D.\ Julio Royo» para el cargo que el decreto 402 había dado a
«Julio Royo Ortiz»: \textbf{o el segundo abrevia el nombre del primero, o son dos personas y uno
de los dos actos está mal}. \textbf{Este libro no concilia y no elige}: transcribe las cuatro
formas. Y anota que un Julio Royo había perdido esos mismos dos cargos ---comisaría y
receptoría--- en cuatro días de noviembre de 1921, y que los recupera, uno cada uno, en los cinco
meses siguientes.
\end{cautela}

\subsection*{La división electoral vuelve a publicarse, y una unidad cambia de nombre}

El 20 de enero de 1922 la Junta Electoral publica otra vez los circuitos del \textbf{Colegio
Electoral Nº 3} del departamento (Nº 891, h.\ 8), con las mismas cuatro mesas y los mismos cuatro
edificios de 1918 y 1921 ---la Escuela Elemental Mixta y la Iglesia Parroquial del pueblo, la
Escuela Infantil Provincial de Vaqueros y la Oficina de Correos de Mojotoro--- y con
\textbf{treinta y cuatro parajes}, los mismos que 1921.

\begin{cautela}
\textbf{Es la cuarta nómina electoral del departamento y la octava lista de sus unidades}, y la
lámina \ref{fig:parajesnom} la incluye como su séptima columna.

\textbf{Y trae una discordancia que el año siguiente resuelve.} En el circuito de la mesa
2, donde 1918 y 1921 imprimen \textbf{Calderilla}, 1922 imprime \textbf{Candelaria}. Las formas se
leyeron en la imagen, cada una en su año. \textbf{Y en 1923 el mismo circuito vuelve a imprimir
Calderilla}, con la misma lista y en el mismo lugar de la enumeración (lámina
\ref{fig:parajesnom}). De las cuatro nóminas electorales, tres dicen Calderilla y una dice
Candelaria. \textbf{Este libro sigue sin conciliar} ---transcribe lo que cada año imprime--- pero
deja registrado que \textbf{1922 es el único año que lo escribe así}, y que el paraje que da
nombre a medio departamento falta de una sola de las nueve listas que la lámina compara. Lo
que no cambia es lo demás: los treinta y cuatro nombres, las cuatro mesas y los cuatro edificios
son los de 1918, cuatro años y tres gobiernos después.

Dos observaciones de impresión de la misma hoja, transcriptas y no corregidas: el original
escribe «\textit{Colejio}» con jota para La Caldera y «Colegio» con ge para Güemes y Metán en la
columna vecina; y \textbf{el rótulo de la mesa Nº 1 no está impreso} ---el bloque pasa del título
del colegio electoral a «\textit{En Caldera Pueblo, Escuela Elemental Mixta}», y se sabe que es
la primera porque la siguiente se rotula «MESA N.\ 2»---.
\end{cautela}

\subsection*{Y las escuelas, por primera vez, con el sueldo del director}

El presupuesto de instrucción primaria de 1922 (Nº 893, h.\ 5) trae el departamento en dos
ítems, y esta vez con la columna de importes completa.

\begin{ficha}{Presupuesto de instrucción primaria --- Departamento Caldera, 1922}
\textbf{Ítem 3, Escuela Elemental de 2.\textsuperscript{a} de Parroquia}: director \$120 y dos
maestros a \$95, que el cuadro totaliza en \$190; \textbf{\$310 mensuales} y \textbf{\$3.720}
al año.

\textbf{Ítem 4, Escuela Infantil de 2.\textsuperscript{a} de Vaqueros}: director \$105 y un
maestro \$95; \textbf{\$200 mensuales} y \textbf{\$2.400} al año.

\textit{Los cuatro subtotales cierran contra lo impreso: 120 más 190 son 310, por doce dan 3.720;
105 más 95 son 200, por doce dan 2.400.}
\end{ficha}

\begin{cautela}
\textbf{Seis mil ciento veinte pesos al año para las dos escuelas del departamento.} La suma es de
este libro y no del cuadro, que no imprime subtotal por departamento.

Puesta al lado del presupuesto de 1921 ---\$345 y \$210 mensuales, que dan \$6.660 al año---,
\textbf{la dotación escolar del departamento bajó un ocho por ciento entre los dos ejercicios}, y
lo hizo perdiendo un cargo: la Elemental de la Parroquia pasa de un director y dos maestros a
\$110 cada uno, a un director y dos maestros a \$95. \textbf{El paraje se llama Parroquia}: es
donde la mesa 2 del colegio electoral monta su urna, en la Iglesia Parroquial, y no un adjetivo
de la escuela.
\end{cautela}

\begin{cautela}
\textbf{Y el presupuesto del municipio, que es el acto más importante del tramo, no se publica.}
El decreto 172 lo aprueba y no lo transcribe. \textbf{En las veintisiete ediciones leídas no hay
una sola cifra de la hacienda del departamento}: las únicas cifras del año que lo nombran son las
dos fianzas de cinco mil pesos del receptor de rentas, la base de ochocientos del remate del
Potrero del Gallinato y los dos ítems escolares. \textbf{Ningún balance municipal}, ni de éste ni
de ningún otro municipio, en todo el tramo ---igual que en 1915, 1916, 1918, 1919 y 1920---.

Es el mismo mecanismo que el apartado inicial de este capítulo describe para 1908, invertido: en
1908 el Boletín publicó el cuadro entero de un municipio que vendía su recaudación; en 1922
publica el decreto que aprueba el presupuesto y no el presupuesto. \textbf{Catorce años después,
lo que se publica es la autorización y no el número.}
\end{cautela}

El circuito queda completo. El Interventor Nacional designa a los cinco miembros del cuerpo. Ese
cuerpo sanciona su presupuesto. El presupuesto sube al Ministro de Gobierno, dictamina el
Fiscal General y lo aprueba un decreto del Poder Ejecutivo. El municipio no decide cuánto
gasta: propone, y la Provincia autoriza.

Sesenta y un años después, en 1983, una ordenanza de La Caldera que rebajaba tasas a
jubilados seguía el mismo recorrido: sancionada por un presidente municipal designado,
elevada a la Secretaría de Estado de Municipalidades «para su aprobación», y aprobada por
decreto ministerial. Ése es el punto de partida institucional del municipio que hoy
administra el corredor: cuando el capítulo \ref{cap:hacienda} discute su capacidad de control sobre el
fraccionamiento, discute sobre una institución que durante buena parte de su historia no
tuvo autonomía ni siquiera para aprobar su propio presupuesto.

Dos de los cinco designados en 1922 son Daniel Linares\index{Linares, Daniel} y Augusto Regis\index{Regis, Augusto}: los mismos que en
1910 integraban la comisión de siete vecinos encargada de las defensas contra el río, y
Linares además como presidente de esa comisión, el que en 1911 rindió las cuentas del
déficit.

El hombre que dirigió la primera obra hidráulica del departamento pasó a integrar el cuerpo
que administraba su presupuesto. En un departamento de pocos miles de habitantes eso
no tiene nada de extraordinario, y no se presenta aquí como hallazgo escandaloso. Se
presenta como dato de estructura: la administración del río y la del municipio fueron,
desde el principio, las mismas personas.

\section{Y en 1923 la Provincia intervino el municipio}

El año 1923 se leyó entero, edición por edición, y es el decimotercer tramo: las
\textbf{cincuenta y dos ediciones} de la \textbf{939 a la 990}, \textbf{673 hojas}, una por cada
viernes del año, \textbf{sin un solo hueco en la serie numérica y sin ninguna edición doble}
(apéndice \ref{cap:fuentes}). Dejó \textbf{diecisiete actos} del departamento, ninguno de
dominio salvo dos avisos de remate, y \textbf{el acto de mayor alcance que el archivo temprano
registra sobre el gobierno de este territorio}.

\begin{ficha}{Decreto Nº 1161 --- Salta, 24 de septiembre de 1923 (B.O.\ Nº 978, h.\ 1 y 2)}
«\textit{Vistas las presentaciones de que instruye el expediente N.\ 4541 F, pidiendo la
intervención del P.\ E.\ a la Comisión Municipal del departamento de La Caldera}», y
considerando que conforme al artículo 173 de la Constitución y al 22 de la Ley Orgánica de
Municipalidades las Comisiones Municipales las nombra el Poder Ejecutivo, «\textit{el que, de
acuerdo al artículo 24 de dicha Ley, puede removerlas de su sola y directa autoridad}», de lo
que «\textit{surge evidente que es también de su exclusiva jurisdicción intervenirlas}»; y que
los hechos de esas presentaciones hacen procedente la intervención «\textit{para regularizar su
funcionamiento y con él obtener la observancia en ella de la Constitución y las leyes}»:

\textbf{Art.\ 1º} --- «\textit{Decláranse intervenida la Municipalidad del departamento de La
Caldera y nómbrase Interventor de ella al señor \textbf{Eduardo F.\ Harbin}}».
\textbf{Art.\ 2º} --- «\textit{Désele las instrucciones respectivas por el Ministerio de
Gobierno}». Firman Güemes y Luis López.
\end{ficha}

\begin{cautela}
\textbf{Diecisiete meses.} El 18 de abril de 1922 un Interventor Nacional creó la Comisión
Municipal del departamento porque «en la actualidad carece de ella»; el 24 de septiembre de 1923
el gobierno provincial la interviene. \textbf{Entre las dos fechas, ese cuerpo perdió dos de sus
cinco miembros}: Augusto Regis\index{Regis, Augusto} renunció en junio de 1922 y otra vez, ya
cubierta la vacante, en enero de 1923 ---el decreto 620 acepta su renuncia y designa a
\textbf{Genaro Molina}, el mismo hombre que fue comisario auxiliar de la Calderilla cuatro años
seguidos entre 1909 y 1913---, y \textbf{Domingo Royo} renunció en mayo de 1923, reemplazado por
\textbf{Julio Royo Ortiz} (decreto 886). \textbf{El cuerpo que se intervino en septiembre ya
había cambiado dos de sus cinco integrantes en cinco meses.}

\textbf{Y lo que decide no se publica.} El expediente 4541 F y las presentaciones que piden la
intervención no se publican; tampoco las instrucciones que el artículo 2º manda darle al
Interventor, ni un solo acto del Interventor dentro del año, ni un solo acto de la Comisión
intervenida. \textbf{Lo publicado permite saber que el municipio fue intervenido y no por qué}:
el decreto dice «los hechos anotados en las presentaciones» y no dice cuáles son. Es la misma
estructura que el capítulo \ref{cap:opacidad} describe cien años después, aplicada esta vez al
acto que suspende el gobierno de un departamento.

\textbf{Y el acto no se llama igual dos veces.} El sumario de la edición dice \textbf{Comisión
Municipal}; el epígrafe del cuerpo, \textbf{Municipalidad}; el vistos, Comisión Municipal; el
considerando tercero, \textbf{Comuna}; y el artículo 1º, Municipalidad. Cuatro nombres para el
mismo cuerpo en dos hojas. Se transcriben todos y no se normalizan.
\end{cautela}

\begin{cautela}
\textbf{Puesto en la serie, 1923 cierra un ciclo de seis años y lo vuelve legible.} En 1918 el
Estado decidió por un umbral de renta que el departamento no estaba «\textit{en condiciones de
ser administrado por una Municipalidad electiva}»; en 1919 declaró que el umbral no era la renta
sino la población; en algún momento entre julio de 1921 y abril de 1922 el departamento se quedó
sin Comisión Municipal, sin que ningún acto publicado lo dijera; en abril de 1922 un Interventor
Nacional se la devolvió; y en septiembre de 1923 el gobierno provincial la intervino.

\textbf{Cinco actos en seis años, ninguno de los cuales cita al anterior}, y los cinco resuelven
lo mismo: quién administra este departamento. Ninguno lo decide el departamento.
\end{cautela}

\subsection*{Tres comisarios en nueve días}

\begin{ficha}{La comisaría del departamento, febrero y marzo de 1923}
\textbf{Decreto 701} del \textbf{21 de febrero}: acepta la renuncia de \textbf{Eustaquio
Murúa}\index{Murúa, Eustaquio} al cargo de Comisario de Policía del departamento, «\textit{atento
sus fundamentos}», que no transcribe (Nº 948, h.\ 8). Es el mismo hombre que en 1912 y 1913 era
comisario auxiliar de La Despensa y que en 1916 juntó la comisaría del departamento, la
expendeduría de guías y la receptoría del impuesto al consumo.

\textbf{Decreto 712} del \textbf{26 de febrero}: nombra interinamente comisario a \textbf{Juan de
la Cruz Nogales}, «\textit{vista la propuesta hecha por el Señor Jefe de Policía}» (Nº 949,
h.\ 3).

\textbf{Decreto 725} del \textbf{1º de marzo}: deja sin efecto ese nombramiento ---«\textit{conferido
por decreto de 24 de Febrero ppdo.}», fecha que no coincide con la del propio decreto 712--- y
nombra comisario del departamento a \textbf{Julio Royo}, hasta entonces comisario de La Silleta,
«\textit{vista la nota 568 de la Jefatura de Policía}», que no se publica (Nº 949, h.\ 7).
\end{ficha}

\begin{cautela}
\textbf{Tres actos sobre el mismo cargo en nueve días}, y los tres se apoyan en piezas que no se
publican: los fundamentos de la renuncia, la propuesta del Jefe de Policía y la nota 568.
\textbf{Y el acto que revoca cita mal la fecha del que revoca}: dice 24 de febrero donde el
decreto revocado imprime 26. Este libro transcribe las dos y no concilia.

\textbf{Y los nombres vienen todos de 1922.} Juan de la Cruz Nogales era, en marzo de 1922,
subcomisario de Mojotoro; en junio de 1923 aparece \textbf{como depositario judicial} de los
bienes rematados en el departamento (Nº 961, h.\ 12). Julio Royo era comisario del departamento
en 1921 y perdió el cargo en cuatro días de noviembre de ese año: \textbf{vuelve al mismo puesto
dieciséis meses después}, y lo hace desde la comisaría de La Silleta, en otro departamento. El
Boletín no dice que el Julio Royo comisario, el Julio Royo Ortiz de la Comisión Municipal y el
Domingo Royo al que éste reemplaza sean parientes, y este libro no lo infiere.
\end{cautela}

\subsection*{Y el año electoral no deja rastro del diputado}

El decreto 629 del 22 de enero convoca a la renovación de la Legislatura para el \textbf{4 de
marzo}, y pone a La Caldera entre los seis departamentos que eligen \textbf{un senador y un
diputado} (Nº 943, h.\ 2 y 3). Dos partidos proclaman: la \textbf{Unión Cívica Radical Nacional},
al senador Ignacio Ortiz y al diputado \textbf{Augusto P.\ Matienzo}; la \textbf{Unión
Provincial}, al senador \textbf{José Antonio Aráoz} ---el mismo candidato que ese partido
presentó por el departamento en 1920 y en 1921--- y al diputado Pedro J.\ Frías (Nº 946,
h.\ 8 y 9). En mayo, la nómina de senadores incorporados registra por Caldera a \textbf{José
A.\ Aráoz} (Nº 956, h.\ 9).

\begin{cautela}
\textbf{Del diputado no hay noticia.} La nómina de mayo es de senadores y el año no trae ninguna
de diputados: quién resultó electo por el departamento, y si llegó a incorporarse, no consta en
las 673 hojas. \textbf{Aráoz es el tercer intento del mismo partido con el mismo candidato}, y
el único que llega.

\textbf{Y las autoridades de mesa se rehacen tres veces en una sola edición.} La edición 947
publica tres nóminas de insaculados de la Junta de Escrutinio, fechadas el 26 de febrero, el 27
y el 1º de marzo. En dos de ellas el mismo hombre ---\textbf{Abelardo López}, suplente segundo
de la mesa 4 del departamento--- es reemplazado dos veces con tres días de diferencia, primero
por Roberto Varela y después por José Isidoro Roldán. Los dos actos se transcriben y no se
concilian.

\textbf{Y la mesa 2 cambia de presidente}: entra Pedro Arévalo en reemplazo de \textbf{Pablo
Choque}. Es, en toda la serie de autoridades electorales del departamento que este capítulo
sigue desde 1912, \textbf{el primer apellido de origen andino que aparece presidiendo una mesa}
---y aparece en el renglón que lo saca---.
\end{cautela}

\begin{cautela}
\textbf{Una nota sobre cómo se encontraron estas tres fichas}, porque mide lo que un barrido por
término no ve. Las nóminas de la Junta de Escrutinio y el aviso de remate de junio imprimen el
nombre del departamento \textbf{con las letras separadas} ---«C A L D E R A», «C ald era»---, de
modo que \textbf{ninguna búsqueda literal las devuelve}. Aparecieron con un patrón que admite
espacios entre todas las letras. \textbf{Tres de los diecisiete actos del año, incluido el único
de dominio, estaban fuera del alcance de la búsqueda por la que este libro encontró casi todo lo
demás} (apéndice \ref{cap:fuentes}).
\end{cautela}

\subsection*{Y la Calderilla sale a remate}

En noviembre, un aviso judicial ofrece «\textit{con base de \$\,7.000 la \textbf{Finca
Calderilla} ubicada en el departamento de La Caldera}», por disposición del juez Cánepa y como
correspondiente al \textbf{juicio sucesorio de don Silvano Murúa}\index{Murúa, Silvano Ignacio};
en el mismo acto, la novena parte indivisa de una casa de la capital con base de \$5.000 y
doscientas diecisiete cabezas de ganado vacuno sin base. Martillero, José M.\ Leguizamón; remate
fijado para el 11 de diciembre de 1923 (Nº 982, h.\ 11).

\begin{cautela}
\textbf{Es el primer acto de dominio del departamento desde 1921}, y recae sobre el nombre que
este capítulo sigue desde el padrón de partidos de 1909. \textbf{La Calderilla aparece aquí como
finca y no como paraje}, y este libro no afirma que sea la misma unidad que el circuito
electoral del mismo año nombra.

\textbf{Y el causante es un Murúa.} Silvano Murúa es el nombre que el archivo registra como
jurado de patentes en 1913, 1914 y 1915, y el apellido de los que en 1924 y 1932 hacen deslindar
fincas del departamento y que en 1936 pierden La Despensa en un remate. \textbf{Coincidencia de
nombre y no identificación registral.} El aviso no dice la superficie: convertir los siete mil
pesos en hectáreas exige el folio, que el apéndice \ref{cap:pedidos} reclama. \textbf{Y el
archivo no dice si se vendió}: el remate cae a trece días del final del año leído.
\end{cautela}

\section{Y en 1924 el cuerpo municipal se desarma entero}

El año 1924 se leyó entero y es el decimocuarto tramo: las \textbf{cincuenta y dos ediciones} de
la \textbf{991 a la 1042}, \textbf{792 hojas}, una por cada viernes del año ---\textbf{el primer
año de toda la serie en que el calendario cierra solo}, sin huecos, sin ediciones dobles y sin
una sola fecha de tapa que no caiga donde debe---. Dejó \textbf{veinticuatro actos} del
departamento en veinte ediciones. \textbf{No es el año más denso del archivo temprano}, como decía una versión anterior de este párrafo: 1918 dejó veintisiete y 1925 dejará veinticinco.

\begin{ficha}{La Comisión Municipal del departamento en 1924}
\textbf{7 de enero, decreto 1381}: acepta la renuncia de \textbf{Julio Royo Ortiz} (Nº 992, h.\ 3).

\textbf{3 de abril, decreto 1560}: acepta las renuncias de \textbf{Daniel Linares}\index{Linares, Daniel}
y \textbf{Genaro Molina}, «\textit{atento el carácter de indeclinable de la misma}» (Nº 1006,
h.\ 3).

\textbf{15 de abril, decreto 1576}, expediente 4969 F: «\textit{vistas las comunicaciones}» que
obran en él, designa para integrar la Comisión a \textbf{Julio Royo Ortiz, Camilo Cáceres,
Francisco Urquiza\index{Urquiza, Francisco S.} y Serafín Yáñez} (Nº 1007, h.\ 3).

\textbf{4 de agosto, decreto 1796}, expediente 5180 F: acepta la renuncia de \textbf{Serafín
Yáñez} (Nº 1023, h.\ 4). \textbf{No se publica reemplazante}: la vacante queda abierta en el
corpus.
\end{ficha}

\begin{cautela}
\textbf{De los cinco que el Interventor Nacional nombró en abril de 1922 no queda ninguno.}
Augusto Regis\index{Regis, Augusto} renunció dos veces y salió en enero de 1923; Domingo Royo, en
mayo de 1923; Daniel Linares y Genaro Molina, en abril de 1924 y con renuncia «indeclinable»;
Jorge Rauch desaparece del archivo sin acto que lo diga. \textbf{Dos años, cinco miembros, seis
renuncias publicadas y una intervención de por medio.}

\textbf{Y el que vuelve es el que se había ido.} Julio Royo Ortiz renuncia el 7 de enero y es
designado otra vez el 15 de abril, al mismo cuerpo, por el mismo Poder Ejecutivo, sin que ningún
acto explique ni la salida ni el regreso. Seis días después de su renuncia, el decreto 1397 lo
nombra ---como «\textbf{Julio Royo}», sin el segundo apellido--- \textbf{expendedor de guías y
marcas del departamento}, con la fianza del artículo 77 de la Ley de Contabilidad. \textbf{Este
libro no unifica las dos formas del nombre}, que el Boletín usa con seis días de diferencia.

\textbf{Y los nombres nuevos no son nuevos.} Camilo Cáceres es el comisario auxiliar de Los
Sauces de 1909 y de Vaqueros de 1910 y el juez de paz suplente de 1917 y 1919: entra a la
Comisión Municipal \textbf{quince años después} de su primer cargo en este mismo archivo.
\end{cautela}

\subsection*{Y el concesionario del agua entra al cuerpo que informó su concesión}

\begin{cautela}
\textbf{Francisco S.\ Urquiza}\index{Urquiza, Francisco S.} aparece \textbf{cuatro veces en
1924, en cuatro calidades distintas}, y conviene ponerlas en fila.

\textbf{En enero} el decreto 1429 le concede los sobrantes del río Wierna para regar «Cerro de
Buena Vista», con el informe favorable de la Comisión Municipal del departamento.
\textbf{En marzo} el edicto de deslinde de la finca «Vaqueros» o «Entre Ríos» lo nombra como
\textbf{propietario del lindero este}. \textbf{En abril} el decreto 1576 lo \textbf{designa
miembro de esa misma Comisión Municipal}. \textbf{Y en octubre} preside la mesa de Vaqueros en
la insaculación electoral.

\textbf{Tres meses separan al beneficiario de la primera concesión de agua del departamento del
cuerpo cuyo informe la fundó.} El Boletín no dice que haya relación entre una cosa y la otra, y
este libro no la afirma: registra la secuencia, que está entera en actos publicados. \textbf{Es
la superposición más completa que el archivo temprano ofrece} ---concesionario de agua,
propietario lindero, miembro del cuerpo municipal y autoridad electoral, en un solo año--- y es
el mismo hombre que veinte años después figura como el mayor de los cuarenta y siete titulares
del revalúo de 1944 (capítulo \ref{cap:fincas}).
\end{cautela}

\subsection*{Una escuela nacional en San Alejo}

El decreto 1580 del 16 de abril, expediente 3247 D, «\textit{presta la aquiescencia del Gobierno
para la instalación de una escuela de la \textbf{Ley 4874}, en la localidad denominada
«\textbf{San Alejo}», Departamento de La Caldera}», a pedido de la Inspección Nacional de
Escuelas y con informe del Consejo General de Educación (Nº 1007, h.\ 3).

\begin{cautela}
\textbf{La Ley 4874 es la ley Láinez}: escuelas nacionales instaladas en territorio de provincia,
con el consentimiento de ésta. \textbf{Es la primera escuela nacional del departamento que este
relevamiento encuentra}, y cae en San Alejo ---el paraje que el circuito de la mesa 1 nombra
desde 1918 y cuya finca este capítulo sigue desde el remate de 1911---. El acto no dice cuántos
alumnos, ni quién la atiende, ni dónde funcionará: \textbf{presta la aquiescencia y nada más}. En
1922 el departamento tenía dos escuelas provinciales, la Elemental de la Parroquia y la Infantil
de Vaqueros; en 1924 se le suma una tercera, y es de la Nación.
\end{cautela}

\subsection*{Y hay que ir a buscar los libros del Registro Civil}

\begin{ficha}{Resolución del Ministro de Gobierno --- Salta, 1º de julio de 1924 (B.O.\ Nº 1017, h.\ 9 y 10)}
«\textit{Vista la comunicación precedente de la Dirección del Registro Civil, de la que resulta,
que en el departamento de la Caldera se impone la necesidad de \textbf{investigar el destino que
hayan tenido los libros} de esa dependencia}». Considerando «\textit{que la medida que se demanda
es de naturaleza urgente, dado el hecho que la motiva e intereses que el mismo puede
comprometer}», autoriza al Jefe del Registro Civil a practicar la diligencia con un empleado,
«\textit{debiendo la comisaría del departamento de la Caldera prestarle toda la cooperación que
fuere necesaria}». Firma Luis López.
\end{ficha}

\begin{cautela}
\textbf{Es el único acto del año que no es decreto ni edicto}, y el único de todo el archivo
temprano en que el Estado provincial manda averiguar qué pasó con sus propios registros en este
departamento. \textbf{No dice que los libros hayan desaparecido}: dice que hay que investigar qué
destino tuvieron. Y no dice cuál es «el hecho que la motiva».

\textbf{Puesto en serie, cobra otro tamaño.} En 1922 el departamento tuvo \textbf{cuatro
encargados del Registro Civil en cuatro meses}, y de uno de ellos ---Santiago Guzmán--- ningún
acto vuelve a decir nada. \textbf{Dos años después hay que ir a buscar los libros.} El archivo no
une las dos cosas y este libro tampoco: las pone una al lado de la otra. \textbf{Y no publica el
resultado}: ningún acto del año dice qué encontró la diligencia.

\textbf{Once años más tarde aparece lo único parecido a una respuesta, y es un renglón de
cuadro.} El Boletín del 14 de junio de 1935 publica un \textbf{inventario de libros del Registro
Civil por oficina}, con el libro, el año de comienzo, el de terminación, los años que encierra,
los folios y el total de actas. El renglón del departamento dice, entero:
«\textit{C[a]ld[e]ra.. matrimonios 1913 1926 13 400 116}».

\textbf{Un solo libro, y de matrimonios.} Las otras oficinas del cuadro traen dos o tres libros
---nacimientos, matrimonios, defunciones---; \textbf{la de La Caldera trae uno}. El cuadro no dice
por qué, no declara que falte ninguno, y este libro no lo infiere. Lo que queda establecido es lo
que el renglón dice: que en 1935, del Registro Civil de este departamento, el inventario oficial
enumeraba \textbf{un libro de matrimonios de 1913 a 1926, cuatrocientos folios y ciento dieciséis
actas}. La resta cierra: 1926 menos 1913 son trece, y la columna dice trece.
\end{cautela}

\begin{cautela}
\textbf{Y conviene anotar quiénes ocupan las oficinas del departamento en los años treinta,
porque cambia la composición.} A la receptoría de rentas, que \textbf{Clara Iradi} ejerce en 1929
y en 1931, la sucede en 1933 \textbf{Julia Regis de Regis}; y el Registro Civil lo toma en 1934
\textbf{Laura Aráoz de Murúa}\index{Murúa, Carlos}. \textbf{Tres mujeres en dos oficinas del
departamento entre 1929 y 1934}, cuando el archivo temprano no registra ninguna entre 1908 y
1928. Los dos apellidos de las dos últimas ---Regis y Murúa--- son de los que este capítulo sigue
desde 1912 y 1922.
\end{cautela}

\subsection*{La Calderilla se remata dos veces y las dos versiones no coinciden}

\begin{ficha}{Los dos avisos de 1924 sobre la misma finca y la misma ejecución}
\textbf{14 de marzo} (Nº 1001, h.\ 18): el martillero José M.\ Leguizamón, por disposición del
juez Figueroa y en la ejecución de Marcos Alsina contra la sucesión de los esposos Mayo,
rematará el 16 de abril «\textit{los derechos y acciones de \textbf{Rafaela de Mamani, Pedro Ríos
y Manuel Lozano} en la Finca «Calderilla»}», \textbf{con base de \$533,33}.

\textbf{27 de junio} (Nº 1016, h.\ 15): el mismo martillero, el mismo juez y la misma ejecución;
rematará el 11 de agosto «\textit{la mitad de la finca Calderilla \ldots\ correspondiente a
\textbf{Pedro Ruiz y Manuel Lozano}}», \textbf{con base de \$266,66}.
\end{ficha}

\begin{cautela}
\textbf{Los dos avisos no dicen lo mismo y este libro no los concilia.} Cambian los titulares
---aparece y desaparece Rafaela de Mamani, y «Pedro Ríos» pasa a ser «Pedro Ruiz»---, cambia el
objeto ---derechos y acciones en la finca, contra la mitad de la finca--- y \textbf{la base de
junio es exactamente la mitad de la de marzo}. Puede ser que el primer remate fracasara y el
segundo saliera a base reducida sobre una porción distinta, que es lo habitual; puede ser que uno
de los dos avisos esté mal compuesto. \textbf{El archivo no lo dice, y el resultado de ninguno de
los dos consta.}

\textbf{Lo que sí queda establecido es quiénes eran los titulares}: tres nombres, dos de ellos de
origen andino, sobre una finca que se remata por quinientos treinta y tres pesos. \textbf{Es la
primera vez que el archivo temprano nombra a los dueños de una fracción de La Calderilla}, y los
nombra en el acto que se la quita.
\end{cautela}

\subsection*{Y un sucesorio firmado en el pueblo}

\begin{ficha}{Edicto sucesorio --- Caldera, 15 de marzo de 1924 (B.O.\ Nº 1012, h.\ 7, aviso 587)}
«\textit{Por disposición del suscrito Juez de Paz de Caldera, se cita, llama y emplaza por el
término de treinta días \ldots\ a todos los que se consideren con derecho a los bienes dejados
por fallecimiento de don \textbf{Simón Yufra}}». Firma \textbf{Carlos F.\ Matienzo}\index{Matienzo, Carlos},
Juez de Paz.
\end{ficha}

\begin{cautela}
\textbf{Es el único acto de todo el año dictado dentro del departamento y fechado allí.} Los
otros veintitrés se dictan en la capital. Un juez de paz de pueblo abre una sucesión, la manda
publicar y firma: \textbf{es la primera vez, en dieciséis años de archivo leído, en que el
departamento aparece en el Boletín Oficial por un acto propio y no por un acto que se le hace}.
\textbf{Y es el comienzo de una serie}: el juez de paz del pueblo vuelve a firmar sucesorios
allí en 1927, 1942, 1943, 1945 y 1946, hasta volverse rutina. \textbf{Lo que en 1924 era único
en dieciséis años, en los años cuarenta es lo normal}, y ese cambio ---de la excepción a la
costumbre--- es él mismo un dato sobre el departamento.

\textbf{Y el apellido importa.} Yufra es uno de los seis que el revalúo de 1944 registra como
titulares de partidas del departamento y que el capítulo \ref{cap:fincas} describe como de origen
andino. \textbf{Veinte años antes, un Simón Yufra muere en este departamento con bienes bastantes
para abrir una sucesión.} El libro no afirma parentesco ni continuidad de dominio: registra que
el apellido está en el departamento en 1924 y sigue estándolo en 1944.

\textbf{Y Matienzo firma como juez de paz en marzo}, cuatro meses antes de que el decreto 1786
del 31 de julio lo nombre juez de paz propietario «\textit{vista la terna elevada por la
H.\ Comisión Municipal de La Caldera}» junto con Vicente Vivas. \textbf{Es el mismo par que el
decreto 105 de junio de 1922 había nombrado}, y el decreto de 1924 no cita al de 1922 ni explica
qué pasó en el medio ---en el medio está la intervención de septiembre de 1923---. La terna que
invoca no se publica.
\end{cautela}

\subsection*{Y el presupuesto municipal se aprueba en noviembre, sin una cifra}

El decreto 1975 del 14 de noviembre, expediente 5162 F, aprueba «\textit{el referido Presupuesto
de Gastos y Cálculo de Recursos de la Municipalidad de La Caldera, para el ejercicio del
corriente año}», con dictamen del Fiscal General e informe del Consejo de Educación (Nº 1037,
h.\ 4).

\begin{cautela}
\textbf{Tercera vez que ocurre lo mismo y la tercera es la más nítida.} El decreto aprueba el
presupuesto por remisión al expediente y \textbf{no trae una sola cifra}; es el único acto
económico del departamento en las 792 hojas del año; y \textbf{se aprueba el 14 de noviembre para
un ejercicio que termina en seis semanas}. En 1908 el Boletín publicó el cuadro entero de
entradas y salidas de este municipio; en 1922 y en 1924 publica el decreto que aprueba el
presupuesto y no el presupuesto.

\textbf{Y lo único que el año deja ver de la hacienda del departamento está en el presupuesto
provincial}: el encargado del Registro Civil, con el mismo sueldo común a más de cuarenta
oficinas de campaña, y \textbf{mil pesos para la Iglesia de La Caldera}. Un subsidio al culto es,
en todo 1924, la única cifra propia que el Estado provincial le asigna por su nombre a este
departamento.
\end{cautela}

\section{Y en 1925 el departamento se parte en dos colegios electorales}

El año 1925 se leyó entero y es el decimoquinto tramo: las \textbf{cincuenta y dos ediciones} de
la \textbf{1043 a la 1094}, \textbf{798 hojas}, una por cada viernes del año y sin huecos en la
serie, como en 1924. Dejó \textbf{veinticinco actos} del departamento.

\begin{ficha}{Las elecciones de 1925 --- decretos de enero y febrero (B.O.\ Nº 1046 h.\ 9, 1047 h.\ 22 y 1048 h.\ 6)}
El decreto de convocatoria del 28 de enero llama a elegir para el domingo 1º de marzo y asigna a
Caldera \textbf{un diputado} ---no un senador y un diputado, como en 1923---.

La Junta de Escrutinio publica la composición de las mesas del \textbf{Colegio Electoral
Núm.\ 5 ---Caldera}: dos mesas, presidentes Fernando Robles y Cirilo E.\ Bejarano, con
Ángel Fernández, Lucas Molina, \textbf{Anastacio Valencia} y Pedro Arévalo de suplentes.

Y el decreto que fija dónde se vota separa por primera vez dos colegios dentro del mismo
departamento:

\textit{Caldera, Colegio Electoral Núm.\ 5} --- Mesa 1, Escuela Elemental Mixta de Caldera;
Mesa 2, Iglesia Parroquial.
\textit{Vaqueros, Colegio Electoral Núm.\ 6} --- Mesa 1, Escuela Infantil Provincial; Mesa 2,
Oficina de Correos de Mojotoro.
\end{ficha}

\begin{cautela}
\textbf{Dos cosas cambian de una vez, y las dos importan.}

\textbf{La primera: el departamento deja de ser un colegio electoral y pasa a ser dos.} De 1918 a
1924 fue siempre el \textbf{Colegio Electoral Nº 3}, con cuatro mesas y un solo rótulo. En 1925
es el \textbf{5} para Caldera y el \textbf{6} para Vaqueros, con dos mesas cada uno ---los mismos
cuatro edificios de siempre, repartidos en dos jurisdicciones electorales---. \textbf{Es la
primera vez que el Estado provincial traza una línea administrativa dentro de este departamento
que separa a Vaqueros del pueblo}, y lo hace sesenta y nueve años antes de que Vaqueros sea
municipio (capítulo \ref{cap:vaqueros}).

\textbf{Y la segunda: en 1925 la nómina de parajes no se publica.} Los actos de ese año fijan los
edificios y las autoridades \textbf{y no enumeran los circuitos}, después de cinco años seguidos
---1918, 1921, 1922, 1923 y 1924--- que publicaron entre treinta y dos y treinta y cuatro nombres
cada uno. \textbf{Vuelve en 1926}, y vuelve recortada: ver más abajo.
\end{cautela}

\subsection*{Un hombre, la comisaría y la caja, en dos días}

\begin{ficha}{Carlos Murúa\index{Murúa, Carlos} en 1925}
\textbf{Decreto 2453, 11 de mayo} (expte.\ 7017 E): lo nombra \textbf{Comisario de Policía de La
Caldera}, en reemplazo de Julio Royo Ortiz (Nº 1063, h.\ 3).

\textbf{Decreto 2465, 12 de mayo}: «\textit{debiendo procederse a la reorganización del personal
dependiente del P.\ Ejecutivo}», lo nombra \textbf{Receptor de Rentas y Recaudador de Impuestos
al Consumo} del departamento, con \textbf{fianza de \$5.000} (Nº 1063, h.\ 8).

\textbf{Decreto 2962, 29 de octubre} (expte.\ 7469 E): acepta su renuncia a la comisaría,
«\textit{dándosele las gracias por los servicios prestados}», y nombra a Francisco Arias
(Nº 1088, h.\ 2). \textbf{Del cargo de receptor no dice nada.}
\end{ficha}

\begin{cautela}
\textbf{Veinticuatro horas entre el que manda la policía del departamento y el que cobra sus
impuestos, y es la misma persona.} Es el mismo par de cargos que Eustaquio Murúa\index{Murúa, Eustaquio}
acumuló en 1916 y que Julio Royo tuvo y perdió en cuatro días de noviembre de 1921. \textbf{Tres
veces en nueve años, con tres personas distintas y dos de ellas del mismo apellido.}

\textbf{Y la salida es asimétrica.} En octubre renuncia a la comisaría con agradecimiento
expreso; el decreto no menciona la receptoría, y ningún otro acto del año la vuelve a nombrar.
\textbf{Este libro no infiere que la haya conservado}: registra que el acto que le quita un cargo
calla sobre el otro.
\end{cautela}

\subsection*{Y la Comisión Municipal, «desintegrada», se rehace con los que se habían ido}

El decreto 2575 del 2 de junio acepta la renuncia de \textbf{Julio Royo Ortiz} ---el mismo a quien
tres semanas antes habían reemplazado en la comisaría--- y el \textbf{decreto 2671 del 25 de
junio}, «\textit{encontrándose desintegrada la H.\ Comisión Municipal de Caldera}», nombra para
integrarla a \textbf{Daniel Linares\index{Linares, Daniel}, Rafael Delgado y Genaro Molina}
(Nº 1067 h.\ 4 y Nº 1070 h.\ 6).

\begin{cautela}
\textbf{Linares y Molina son exactamente los dos que en abril de 1924 habían renunciado con
carácter «indeclinable»}, y vuelven al mismo cuerpo catorce meses después. El decreto no cita
aquella renuncia ni explica la vuelta. \textbf{Es la tercera composición de la Comisión Municipal
en tres años}, después de la del Interventor Nacional de 1922 y la del decreto 1.576 de 1924, y
ninguna de las tres sale de una elección.

\textbf{Y el que entra nuevo también acumula.} \textbf{Rafael Delgado} es nombrado en junio a la
Comisión Municipal, en agosto encargado interino del Registro Civil del departamento por licencia
de José Delgado, y en octubre titular de ese Registro por renuncia del mismo (decretos 2765 y
2947). \textbf{El Boletín no dice que haya parentesco entre José y Rafael Delgado y este libro no
lo infiere}; registra que en el mismo semestre uno reemplaza al otro en la oficina que anota los
nacimientos y las muertes del departamento, y que el que lo reemplaza integra además el cuerpo
que lo gobierna.
\end{cautela}

\subsection*{Y Urquiza era, además, subcomisario de Vaqueros}

\begin{ficha}{Decreto 2674 --- Salta, 25 de junio de 1925 (expte.\ 7163 E; B.O.\ Nº 1070, h.\ 7)}
«\textit{Nómbrase con carácter ad-honorem Subcomisario de «Vaqueros» Caldera, al señor Antonio
Agüero (hijo) \textbf{en reemplazo de D.\ Francisco S.\ Urquiza}}»\index{Urquiza, Francisco S.}.
\end{ficha}

\begin{cautela}
\textbf{Es la quinta calidad del mismo hombre y la que este libro no esperaba.} Urquiza fue, en
poco más de un año: \textbf{concesionario} de los sobrantes del Wierna (enero de 1924),
\textbf{propietario lindero} en el deslinde de la finca Vaqueros (marzo de 1924), \textbf{miembro
de la Comisión Municipal} del departamento (abril de 1924), \textbf{presidente de mesa} en
Vaqueros (octubre de 1924) y \textbf{subcomisario \textit{ad honorem} de Vaqueros} hasta junio de
1925.

\textbf{El acto que lo saca no dice desde cuándo estaba}, y ningún acto del archivo leído lo
nombra: \textbf{el nombramiento no se publicó, o se publicó fuera de los tramos leídos}. Es el
mismo mecanismo que este capítulo describe desde 1909 ---la vacante que se declara sin nombrar a
quien la dejó---, invertido: aquí lo que falta es el principio y no el final.

\textbf{Y conviene decir qué es una subcomisaría \textit{ad honorem}}: el policía del paraje, sin
sueldo. De modo que el mismo hombre que riega ciento noventa hectáreas con agua del Wierna, que
linda con la finca que se deslinda, que integra el cuerpo que informó su propia concesión y que
preside la mesa donde vota Vaqueros es también, sin cobrar, \textbf{la autoridad policial del
paraje}. Este libro no afirma que ninguna de esas calidades haya servido a otra: afirma que las
cinco están publicadas y que son de la misma persona.
\end{cautela}

\subsection*{Los cateos de 1925 tienen ahora nombre, finca y superficie}

\begin{ficha}{Los pedimentos mineros de enero de 1925, publicados a lo largo del año}
\textbf{1166-C} --- Povoli y Premoli, fincas \textbf{San Alejo y Santa Rufina} (Nº 1058, h.\ 9).
\textbf{1165-C} --- Roldán y Premoli, las mismas dos fincas (Nº 1072, h.\ 15).
\textbf{1195-C} --- Julio Juárez, terrenos de los herederos de Juan A.\ Pinto (Nº 1080, h.\ 9).
\textbf{1169-G} --- Emilio D.\ Sylvester, finca \textbf{Las Nieves} de Garzón y Pintos,
\textbf{cuatro unidades (2.000 hectáreas)} (Nº 1081, h.\ 12).
\textbf{1168-C} --- \textbf{Juan Larrán}, finca Las Nieves de Garzón y Pintos «\textit{y en
terrenos también del señor E.\ Quipildor}», cuatro unidades (Nº 1069, h.\ 15), \textbf{con su auto
de concesión del 1º de agosto de 1925} (Nº 1077, h.\ 7), anotado en la Sección Minas del
Departamento Topográfico bajo el Nº 87.
\end{ficha}

\begin{cautela}
\textbf{Esto confirma con documento lo que este capítulo afirmaba por la tabla de caducidades}:
de los nueve pedimentos de enero de 1925, \textbf{el de Larrán es el que se concede}, y ahora
consta el auto, su fecha y su número de registro.

\textbf{Y agrega dos cosas que la tabla no traía.} La primera: \textbf{las fincas}. Los cateos no
caen en cualquier parte: caen sobre \textbf{San Alejo, Santa Rufina y Las Nieves}, que son tres de
las fincas que este libro sigue desde 1909, y \textbf{Las Nieves es de «Garzón y Pintos»}. La
segunda: \textbf{un tercer propietario, E.\ Quipildor}, cuyos terrenos el pedimento de Larrán
nombra junto a los de Garzón y Pintos. \textbf{Es el primer apellido de origen andino que el
archivo temprano registra como propietario de terrenos en el departamento}, y aparece en el acto
que autoriza a explorar debajo de ellos.

\textbf{Y los dos edictos de Las Nieves se apoyan en un deslinde de 1889}, del agrimensor
P.\ Arias, del que citan el vértice 49. \textbf{Es el documento más antiguo que este libro ve
citado sobre el departamento}, treinta y seis años anterior al corpus, y queda como citado y no
acompañado: el apéndice \ref{cap:pedidos} lo reclama.
\end{cautela}

\begin{cautela}
\textbf{Y un deslinde de agosto nombra, en un solo renglón, a tres de los que este capítulo
persigue.} El edicto de las fincas «Sauce» o «Paraíso» y la fracción Ceival, del \textbf{doctor
José María Solá} ---parte en Campo Santo y parte en La Caldera---, describe por linderos
«\textit{la fracción denominada \textbf{«Mosquera», del señor Daniel Linares}}»\index{Linares, Daniel},
«\textit{los «Noques», de Reymundo Echenique}» y, al poniente, «\textit{\textbf{«La Despensa», de
los herederos Murúa}}» (Nº 1075, h.\ 11).

\textbf{Daniel Linares es, desde hace dos meses, miembro de la Comisión Municipal del
departamento}, y aparece aquí como propietario de una fracción que se deslinda. \textbf{La
Despensa figura en manos de los herederos Murúa} once años antes del remate de 1936 que este
libro registra, y \textbf{Murúa es el apellido del comisario y del receptor de rentas del mismo
año}. Coincidencias de nombre, todas, y no identificaciones registrales. Lo que queda dicho es
que en 1925 el archivo nombra, en el mismo departamento y en el mismo semestre, a los mismos
apellidos como autoridad y como propietario.
\end{cautela}

\section{1926 y 1927: la lista vuelve sin Vaqueros, y la hacienda municipal reaparece}

Los dos años se leyeron enteros y son el decimosexto y el decimoséptimo tramo: \textbf{1926}, las
ediciones \textbf{1096 a 1147} y \textbf{769 hojas}; \textbf{1927}, las \textbf{1148 a 1199} y
\textbf{745 hojas}. Los dos son semanales, los dos tienen una edición por cada uno de los
cincuenta y dos viernes del año y ninguno tiene huecos en la serie. Dejaron \textbf{dieciocho y
veinticuatro actos} del departamento.

\subsection*{La nómina vuelve, con veintiséis nombres en lugar de treinta y cuatro}

El decreto 3156 del 30 de enero de 1926 fija otra vez las mesas y esta vez \textbf{sí enumera los
circuitos}: «\textit{La Caldera Colegio Electoral Nº 5}», mesa 1 en la Escuela Elemental Mixta con
Pueblo, Potrero de Castillo, Peñones y San Alejo; mesa 2 en la Iglesia Parroquial con Angosto,
Angosto de Arrieta, Angostura, Campo Alegre, Calderilla, Corral de Barranca, Cachanio, Despensa,
Despensilla, El Acheral, Churcal, \textbf{El Monte de las Garrapatas}, Los Porongos, Los Sauces,
Mesada, Potrero de Valencia, Quesería, Santa Rufina, Wierna, Yacones, El Bordo y Las Lagunas.
\textbf{Veintiséis en total}, y a continuación, en el mismo acto, «\textit{Vaqueros Colegio
Electoral Nº 6}». La lista de 1927 repite exactamente la misma (Nº 1153).

\begin{cautela}
\textbf{La lista no se acortó: se partió.} Los ocho nombres que faltan respecto de 1924 no
desaparecieron del departamento ---Lesser, Vaqueros, Mojotoro, Mojotorito, Potrero de Gallinato,
el Túnel y la Estación del Ferrocarril--- sino que \textbf{pasaron al Colegio Electoral Nº 6}, el
de Vaqueros, que el mismo decreto enumera aparte. \textbf{Desde 1926 el departamento tiene dos
listas de parajes y no una}, y la lámina \ref{fig:parajesnom} sólo puede seguir la del Colegio 5
porque es la que el acto atribuye a «La Caldera» por su nombre.

\textbf{Y el octavo nombre que falta es de otra especie}: 1926 imprime «\textbf{El Monte de las
Garrapatas}» como un solo topónimo donde 1918 a 1924 imprimían «El Monte» y «Las Garrapatas»
separados ---1918 sin coma entre los dos, 1922 y 1923 con ella---. Es la misma forma que el padrón
de 1940 volverá a usar. \textbf{Este libro no decide si son uno o dos}: cuenta lo que cada año
imprime y deja ver la oscilación, que es de la fuente y no de la lectura.
\end{cautela}

\subsection*{Y la hacienda del municipio vuelve al Boletín, sin una cifra}

\begin{ficha}{Los dos actos económicos de 1927 (B.O.\ Nº 1168 h.\ 5 y Nº 1170 h.\ 2)}
\textbf{Decreto 5024, 17 de mayo de 1927}, expediente 1099 M. «\textit{Vista la memoria y
rendición de cuentas de la inversión de las rentas recaudadas en 1926 que ha presentado la
Comisión Municipal de La Caldera y atento el informe del Consejo General de Educación que
manifiesta su conformidad en lo que respecta a la \textbf{renta escolar} que por dicho ejercicio
le corresponde}», aprueba esa rendición de cuentas.

\textbf{Decreto 5036, 20 de mayo de 1927}, expediente 1725 M: aprueba el \textbf{Presupuesto de
Gastos y Cálculo de Recursos} de la misma Comisión para el ejercicio en curso.
\end{ficha}

\begin{cautela}
\textbf{Es la primera rendición de cuentas municipal del departamento que el archivo temprano
registra}, y llega dieciocho años después del cuadro de 1908. \textbf{Y no trae una sola cifra}:
los dos decretos aprueban por remisión al expediente. Con el decreto 172 de 1922, el 1.975 de
1924 y el 3420 de 1926, son \textbf{cinco actos económicos del municipio publicados sin su
número}, en seis años.

\textbf{Lo que sí agregan es el mecanismo}: los dos invocan el informe del Consejo General de
Educación, y el de la rendición dice por qué ---la \textbf{renta escolar} que le corresponde a ese
Consejo sobre la recaudación municipal---. \textbf{De modo que una parte de lo que el municipio
recaudaba iba a la educación provincial}, y el Consejo tenía que conformarse antes de que la
Provincia aprobara las cuentas. Es la primera vez que el archivo muestra a un tercero con derecho
sobre la caja de este municipio.

\textbf{Y el decreto 5036 no está solo.} El acto inmediato siguiente hace lo mismo para General
Güemes y el sumario anuncia otro para La Poma: \textbf{es una tanda de aprobaciones municipales
despachadas el mismo día}. El trámite que a La Caldera le costó, en 1922, tres meses entre la
sanción y la aprobación, en 1927 sale en lote.
\end{cautela}

\subsection*{Y Los Sauces se remata tres veces en cuatro meses}

\begin{ficha}{La finca «Los Sauces» en 1927}
\textbf{3 de junio} (Nº 1169, h.\ 14, aviso 2188): «\textit{Por disposición del señor Juez de
1ª Instancia y 1ª Nominación en lo Civil y Comercial, doctor Ángel María Figueroa y como
perteneciente al \textbf{juicio sucesorio de Domingo Royo}}», el martillero Roberto del Carril
venderá el 14 de junio «\textit{con base de Cincuenta mil pesos moneda nacional, o sea la de su
avaluación pericial, una hermosa finca denominada «LOS SAUCES»}».

\textbf{8 de julio} (Nº 1174, h.\ 11, aviso 2233): el mismo remate, \textbf{con base de
\$14.000}.

\textbf{26 de agosto} (Nº 1182, h.\ 14, aviso 2349): el mismo texto, la misma base y los mismos
linderos, con otro día de remate y otro número de aviso.
\end{ficha}

\begin{cautela}
\textbf{El causante es Domingo Royo}, el miembro de la Comisión Municipal del departamento que el
Interventor Nacional nombró en abril de 1922 y que renunció en mayo de 1923. \textbf{Cuatro años
después, la finca Los Sauces se remata en su sucesión}, y la base cae de cincuenta mil pesos a
catorce mil ---a menos de un tercio--- \textbf{en cinco semanas}. El archivo no dice si se vendió
ni a quién.

\textbf{Y Los Sauces es uno de los catorce nombres permanentes} de la lámina
\ref{fig:parajesnom}: está en once o en las doce nóminas de 1909 a 1940, era partido policial en 1909 y es
circuito de la mesa 2 en 1926 y 1927. \textbf{Es la primera vez que el archivo temprano lo nombra
como finca con dueño y con precio.}

\textbf{Y un detalle de método que el año deja ver.} Los tres avisos son el mismo remate y llevan
\textbf{tres números de aviso distintos}: cotejar los números no los habría reunido; cotejar el
texto, sí. Es la razón por la que este libro cuenta actos y no avisos.
\end{cautela}

\begin{cautela}
\textbf{Y en 1927 vuelve a firmar un juez de paz dentro del departamento.} El edicto sucesorio de
\textbf{Milagro Catacata}, fechado «\textit{Caldera, Junio 13 de 1927}» y firmado por
\textbf{Ángel Fernández, Juez de Paz} (Nº 1171, h.\ 11, aviso 2200). Es el tercero de la serie,
después del de Simón Yufra de 1924 y los dos de 1925, y el apellido es otra vez de origen andino.

\textbf{Los cuatro sucesorios que el archivo temprano registra dictados dentro del departamento
llevan apellidos que el revalúo de 1944 volverá a encontrar en sus partidas.} Este libro no
afirma continuidad de dominio ni parentesco: registra que las mismas familias que abren sucesión
ante el juez de paz del pueblo en los años veinte figuran como titulares veinte años después.
\end{cautela}

\section{1928: el año en que todo el aparato del departamento se renueva}

El año 1928 se leyó entero y es el decimoctavo tramo: las ediciones \textbf{1200 a 1251},
\textbf{52 ediciones} y una por cada viernes. Dejó \textbf{veintitrés actos} del departamento.
\textbf{Es el año que el capítulo \ref{cap:politica} cuenta}, y el barrido lo confirma entero y
le agrega el principio y el final, que ese capítulo no tenía.

\begin{ficha}{Lo que pasa en el departamento entre marzo y junio de 1928}
\textbf{26 de marzo, decretos 6776 y 6777}, uno detrás del otro: aceptan las renuncias de
\textbf{Victorio Offredi} al juzgado de paz propietario y de \textbf{Ramón Balvoa} ---así
impreso--- a la comisaría de policía. Ninguno da la causa (Nº 1214, h.\ 14).

\textbf{21 de abril, decreto 6877}, expediente \textbf{3898 C} ---el mismo de la concesión de
1924---: rechaza los reclamos de vecinos y de la Comisión Municipal de Campo Santo contra la
concesión de agua de Urquiza\index{Urquiza, Francisco S.}, sobre el informe de Obras Públicas que
dice que la acequia «\textit{va a morir en el Río Vaqueros}» y que, lejos de disminuir el caudal
del Mojotoro, tiende a aumentar el del Vaqueros (Nº 1217, h.\ 11).

\textbf{30 de abril, decreto 6958}, expediente 674 U: acepta la renuncia de Urquiza a la Comisión
Municipal (Nº 1218, h.\ 12).

\textbf{9 de mayo, decreto 7095}: nombra comisario de policía a \textbf{Julio Royo Ortiz}.
\textbf{Decreto 7096}: «\textit{debiendo procederse a la reorganización de la Municipalidad de La
Caldera, \textbf{para asegurar la existencia del régimen municipal}}», declara en acefalía la
Comisión y nombra comisionado \textit{ad honorem} a Luis E.\ Guardo. \textbf{Decreto 8002}, mismo
día: \textbf{Teófilo Reyes} al Registro Civil, en reemplazo de Ramón Balboa (Nº 1222, h.\ 6 y 7).

\textbf{12 de mayo, decreto 8023}: deja sin efecto la concesión de agua, «\textit{considerando que
el mantenimiento de dicha concesión es inconveniente}». \textbf{Decreto 8024}: acepta la renuncia
de Eustaquio Guerra a la subcomisaría de Mojotoro ---a la que el decreto 8001 lo había nombrado
\textbf{tres días antes}--- y nombra a Juan de la Cruz Nogales.

\textbf{18 de mayo, decreto 8099}: Royo Ortiz, además, receptor de rentas y expendedor de guías.
\textbf{26 de junio, decreto 9306}: nueva Comisión Municipal.
\end{ficha}

\begin{cautela}
\textbf{El capítulo \ref{cap:politica} lee esta secuencia desde el conflicto del agua. El barrido
del año la ensancha, y lo que muestra es más grande y menos dirigido.}

\textbf{No cambió sólo la comisión municipal: cambió todo.} Entre marzo y junio de 1928 el
departamento renovó \textbf{el juez de paz propietario, el comisario de policía, el encargado del
Registro Civil, el receptor de rentas, el subcomisario de Mojotoro, la comisión municipal entera
y, en octubre, el juez de paz suplente}. \textbf{Siete cargos en siete meses, y no hay uno solo
que quede en pie.} Dos de esos cambios ---las renuncias de Offredi y de Balvoa--- son del 26 de
marzo, \textbf{seis semanas antes de la acefalía} y casi un mes antes del decreto que rechaza el
reclamo de Campo Santo.

\textbf{Eso no desmiente la lectura del capítulo \ref{cap:politica}: la encuadra.} La renuncia de
Urquiza y la acefalía ocurren dentro de un recambio general que empieza antes del episodio del
agua y sigue después, con dos gobernadores ---Corbalán hasta el 9 de mayo, Cornejo desde
entonces--- y bajo la fórmula «\textit{para asegurar la existencia del régimen municipal}», que el
decreto 7094 del mismo día usa igual para Embarcación. \textbf{Lo que el archivo prueba es la
secuencia, y lo que la secuencia prueba es que en este departamento el aparato entero cuelga de
la firma del gobernador de turno.}

\textbf{Y el recambio no fue prolijo.} El decreto 8001 nombra a Eustaquio Guerra subcomisario de
Mojotoro el 9 de mayo y el 8024 le acepta la renuncia el 12: \textbf{tres días}, con el
nombramiento y la renuncia publicados en la misma edición, en hojas contiguas y sin una remisión
entre sí. \textbf{La lámina \ref{fig:cargos} muestra ese año como lo que es}: un corte vertical
en el que ningún segmento cruza entero. Y el decreto 9464 del 1º de agosto, que reconoce servicios a los ex-encargados del
Registro Civil, le asigna a Ramón Balboa un período que \textbf{empieza el 9 de junio y termina el
13 de mayo}: se transcribe y no se concilia.
\end{cautela}

\begin{cautela}
\textbf{Y el nombre completo aparece por fin.} La nómina de autoridades de mesa de la edición
1205 imprime «\textbf{Urquiza Francisco Solano}» presidiendo la mesa del circuito 6. Es la
\textbf{sexta calidad} del mismo hombre ---concesionario, lindero, miembro de la comisión
municipal, subcomisario \textit{ad honorem} de Vaqueros, presidente de mesa en 1924 y otra vez en
1928--- y la primera vez que el archivo da su nombre entero. \textbf{El Boletín no cruza en ningún
lado los tres hechos del año}: que preside una mesa, que renuncia a la comisión y que su
concesión se revoca.

\textbf{Y una advertencia sobre el salto de numeración}, porque el barrido no lo lee como este
libro. El capítulo \ref{cap:politica} sostiene que la cifra de los miles avanza cada cien
decretos y que por eso 7099 y 8000 son consecutivos. El barrido registra el hecho ---en la misma
hoja y con la misma fecha, la serie pasa de \textbf{7097 a 8001}, novecientos números en un día,
con los dos tramos firmados por el mismo gobernador y el mismo ministro--- y \textbf{declina
conciliarlo}. Las dos lecturas se dejan escritas: la de este libro explica el patrón y no está
probada, y el barrido prefiere no explicarlo. El apéndice \ref{cap:pedidos} pide el Registro
Oficial del año, que lo resolvería.
\end{cautela}

\section{1929--1932: el agua se legisla desde el departamento vecino}

Los cuatro años se leyeron enteros y son los tramos decimonoveno a vigésimo segundo: las
ediciones \textbf{1252 a 1303}, \textbf{1304 a 1355}, \textbf{1357 a 1407} y \textbf{1408 a
1461}. \textbf{Ninguno de los cuatro publica un solo plano}: las 1.321 hojas de 1931, por
ejemplo, son texto a dos columnas sin una lámina, un croquis ni una cartela.

\subsection*{La ley que le da a Campo Santo poder sobre el agua de La Caldera}

\begin{ficha}{Ley Nº 11.078 --- sancionada el 27 de septiembre y promulgada el 1º de octubre de 1929 (B.O.\ Nº 1292, h.\ 9)}
«\textit{Art.\ 1º.---Autorízase a la \textbf{Municipalidad del Departamento de Campo Santo} a
efectuar los trabajos que conceptúe necesarios en la Playa del Río Mojotoro, para reconcentrar el
agua, \textbf{desde las juntas del Río de Wierna con el de La Caldera}, vigilar y \textbf{prohibir
se abran nuevas bocas-tomas} no autorizadas por Ley}».
\end{ficha}

\begin{cautela}
\textbf{Es la pieza que le faltaba al capítulo \ref{cap:infraestructura}, y es de 1929.}

La ley no le da a La Caldera ninguna facultad sobre su agua: se la da al municipio de aguas
abajo. \textbf{Quien queda autorizado a intervenir la playa del río, a reconcentrar el caudal y a
prohibir bocas-tomas desde la confluencia del Wierna con el río La Caldera es la Municipalidad de
Campo Santo}, un departamento distinto. Y el punto de arranque que la ley elige ---«las juntas»---
está dentro del territorio de La Caldera.

\textbf{Es el único acto de los cuatro años en que el departamento aparece como origen de un
recurso y no como jurisdicción.} Cinco años antes, la Comisión Municipal de La Caldera había
informado favorablemente la concesión de agua de Urquiza\index{Urquiza, Francisco S.} y la de
Campo Santo había reclamado contra ella (capítulo \ref{cap:politica}); un año después de aquel
reclamo rechazado, \textbf{la Legislatura le dio a la reclamante la facultad que el trámite no le
había dado}.

\textbf{Y dos años más tarde el arreglo se institucionaliza.} En octubre de 1931 un decreto crea
una \textbf{Comisión \textit{ad hoc}} sobre las aguas de los ríos Vaqueros y Wierna, afluentes
del Mojotoro, \textbf{integrada por los Comisionados Municipales de Campo Santo y de La Caldera}
(B.O.\ Nº 1399, h.\ 7 y 8). Es la primera vez que el archivo muestra a los dos municipios
sentados a decidir juntos sobre esa agua ---y los dos estaban, en 1931, gobernados por
comisionados designados y no por cuerpos electos---.
\end{cautela}

\subsection*{Y la ley que le compra al pueblo su Casa Municipal}

\begin{ficha}{Ley Nº 11.088 --- sancionada el 27 de septiembre y promulgada el 2 de octubre de 1929 (B.O.\ Nº 1293, h.\ 7)}
«\textit{Art.\ 1º---Autorízase al P.\ Ejecutivo a invertir hasta la suma de \textbf{quinientos
pesos} m/nacional (\$500.) en la adquisición de un terreno en el \textbf{pueblo de La Caldera,
Capital del mismo Departamento}, con destino a la \textbf{Casa Municipal}}».
\end{ficha}

\begin{cautela}
\textbf{Veintiún años después del presupuesto de 1908, el municipio no tenía dónde funcionar.} La
ley no amplía ni refacciona: \textbf{compra el terreno}, y lo hace con quinientos pesos de Rentas
Generales de la Provincia. \textbf{El Estado provincial le pone al municipio la sede que el
municipio no podía pagarse}, en el mismo año en que otra ley le da al municipio vecino poder sobre
su río.

Es también el único acto del archivo temprano que llama al pueblo «\textbf{Capital del mismo
Departamento}». Y el archivo no dice si el terreno se compró: \textbf{la ley autoriza y el corpus
no vuelve sobre ella}.
\end{cautela}

\subsection*{Y en 1929 alguien catastró el departamento}

\begin{ficha}{Decreto 10.910 --- 5 de agosto de 1929, expediente 3463 D (B.O.\ Nº 1286, h.\ 17)}
El catastrador \textbf{Adeodato Aybar} manifiesta «\textit{habiendo terminado la
\textbf{catastración} de los Departamentos de \textbf{Campo Santo, Caldera y La Viña}}» y pide que
se le liquiden quinientos pesos a cuenta de la remuneración «\textit{que le sea fijada
oportunamente}». El gasto se imputa a la \textbf{Ley 9720 del 6 de octubre de 1928}.
\end{ficha}

\begin{cautela}
\textbf{Es el primer catastro del departamento que este libro encuentra}, y está terminado en
agosto de 1929. \textbf{El decreto lo da por hecho y no publica nada de él}: ni superficie, ni
número de parcelas, ni plano, ni dónde quedó. \textbf{Y la remuneración del catastrador queda sin
fijar}, lo que hace dudar de cuánto avanzó realmente el expediente.

Puesto al lado del acto de marzo de 1931, en que la Dirección General de Obras Públicas
«\textit{ratifica la necesidad de confeccionar el plano catastral de la Provincia [\ldots] en
láminas que comprendiendo una cada departamento}», queda una pregunta con fecha: \textbf{si el
departamento estaba catastrado desde 1929, por qué en 1931 había que ratificar la necesidad de
hacer el plano}. El apéndice \ref{cap:pedidos} reclama el expediente 3463 D y los registros de esa
catastración: \textbf{es el antecedente más antiguo del catastro con el que este libro trabaja hoy}
(capítulo \ref{cap:fincas}).
\end{cautela}

\begin{cautela}
\textbf{Y tres cosas más de estos cuatro años, que valen por lo que muestran del aparato.}

\textbf{En 1929 la receptoría de rentas del departamento la ejerce Clara Iradi} (decreto 10.514):
\textbf{es la primera mujer que el archivo temprano registra con un cargo público en La Caldera},
y el cargo es el de la caja. \textbf{Y en marzo de ese año el comisario de policía es Manuel
Condori} (decreto 10.389), que lleva uno de los seis apellidos que el revalúo de 1944 registrará
entre sus titulares: \textbf{es el primer apellido de origen andino al frente de la comisaría del
departamento}, veinte años después del padrón de comisarios auxiliares de 1909.

\textbf{En septiembre de 1930 el gobierno pasa a un Interventor Nacional}, y desde octubre los
actos del departamento los firma él: la Comisión Municipal, las subcomisarías, los nombramientos.
\textbf{Es la segunda intervención federal que este capítulo registra sobre el departamento en
ocho años}, después de la de 1922, y otra vez el cuerpo municipal se compone por decreto.

\textbf{Y el departamento cambia de colegio electoral por tercera vez.} Fue el Nº 3 de 1918 a
1923, el Nº 5 en 1925, 1926 y 1927, y desde 1929 es el \textbf{Nº 2}, que es el que tendrá en el
padrón de 1940. \textbf{Tres numeraciones en once años}, y ninguna de las tres la decide el
departamento.
\end{cautela}

\section{Veintiocho años de cargos, vistos juntos}

Los actos que este capítulo fue reuniendo año por año forman, puestos en una sola escala, la
lámina \ref{fig:cargos}: \textbf{setenta y cuatro segmentos y cincuenta y cuatro personas} en seis cargos del departamento entre
1908 y 1946, más el estado del cuerpo municipal y los cuatro episodios de tutela.

\begin{figure}[p]
\centering
\includegraphics[width=\textwidth,height=0.88\textheight,keepaspectratio]{fig-cargos-linea.png}
\caption[Quién administra el departamento de La Caldera, 1908--1946]{\textbf{Quién administra el
departamento de La Caldera, 1908--1946.} Un segmento por titular en los seis cargos unipersonales
que el Boletín publica ---comisaría de policía, juzgado de paz, Registro Civil, receptoría de
rentas y las subcomisarías de Mojotoro y de Vaqueros---; arriba, separado, el estado de la
Comisión Municipal, que no es un cargo de una persona sino un cuerpo; y en línea vertical, los
cinco episodios de tutela. \textbf{Los bordes degradados son los que el archivo no fecha}: el
segmento sigue, y lo que termina ahí es la prueba ---cincuenta y ocho de los noventa y cuatro
bordes salen de un acto publicado, el cincuenta y nueve por ciento---. Todo sale de los años leídos
edición por edición (apéndice \ref{cap:fuentes}). Elaboración propia; dibujo de
\texttt{scripts/grafico\_cargos.py} sobre \texttt{datos/cargos/} del repositorio
\texttt{mapas\_caldera} (apéndice \ref{cap:fuentes}).}
\label{fig:cargos}
\end{figure}

\begin{cautela}
\textbf{Cuatro cosas se leen ahí que ningún año suelto deja ver.}

\textbf{La primera: el primer tercio está casi vacío.} De 1908 a 1920 hay cuatro segmentos en
toda la lámina. \textbf{No es que el departamento no tuviera cargos}: es que el Boletín no
publicaba los nombramientos, y sobre esos trece años el archivo sólo entrega un juez de paz en
1909, otro en 1912, un tercero en 1918 y un receptor de rentas en 1916. \textbf{Lo que la lámina
mide en esa mitad no es el gobierno del departamento sino la publicidad de sus actos.}

\textbf{La segunda: 1922 y 1923 son un enjambre.} La comisaría tiene cinco titulares entre
noviembre de 1921 y marzo de 1923; el Registro Civil, cuatro entre marzo y junio de 1922. Los
rótulos de esos segmentos no entran dentro de sus barras y hubo que escalonarlos en cuatro
alturas: \textbf{la ilegibilidad del dibujo en ese tramo es el dato}.

\textbf{La tercera: 1928 es un corte vertical.} Entre marzo y junio se cortan la comisaría, el
juzgado de paz, el Registro Civil, la receptoría y la subcomisaría de Mojotoro, y la Comisión
Municipal pasa por acefalía. \textbf{Ningún segmento cruza ese año entero.}

\textbf{Y una quinta, que el tramo 1937--1946 agrega.} La lámina llega ahora hasta 1946 y
muestra \textbf{el bloque más largo de toda la serie}: Augusto Regis\index{Regis, Augusto}
presidiendo la Comisión Municipal sin interrupción de 1937 a 1943, con su propio segmento
además en la fila de la comisaría en 1942 y 1943. \textbf{Y termina con la tercera
intervención}, que corre de 1944 a agosto de 1946 y cuyo borde izquierdo se degrada porque el
archivo no dice cuándo empezó.

\textbf{Y la cuarta: el caso de Urquiza\index{Urquiza, Francisco S.} ocupa dos filas a la vez}
---la subcomisaría de Vaqueros y, arriba, el cuerpo municipal--- en los mismos meses de 1924 en
que recibe la concesión de agua. \textbf{Y su segmento de la subcomisaría se degrada hacia la
izquierda}, porque el acto que lo nombró no aparece en ninguno de los tramos leídos: la lámina
dibuja esa ignorancia en lugar de taparla.
\end{cautela}

\begin{cautela}
\textbf{Qué no prueba la lámina \ref{fig:cargos}.} No prueba que un cargo estuviera vacante donde
no hay segmento: prueba que el Boletín no publicó quién lo ocupaba. No prueba la fecha exacta de
ningún relevo cuyo borde esté degradado. \textbf{Y mezcla, en la fila del juzgado de paz,
propietarios y suplentes}, porque varios actos no distinguen; el apéndice \ref{cap:personas} sí
los separa cuando el acto lo dice. La tabla que la produce ---\texttt{datos/cargos/}--- trae por
cada segmento el acto y la edición de donde sale, de modo que cualquiera puede ir a discutir
segmento por segmento.
\end{cautela}

\section{1937 y 1938: la primera defensa del río, y la renta escolar}

Los dos años se leyeron edición por edición y son los tramos vigésimo tercero y
vigésimo cuarto: las ediciones \textbf{1669 a 1721} y \textbf{1722 a 1773}. Es
el primer tramo del libro leído con una novedad material: \textbf{los PDF de
esta época ya traen capa de texto del proveedor}, de modo que el noventa y
nueve por ciento de las hojas no necesitó reconocimiento propio (apéndice
\ref{cap:fuentes}).

\begin{ficha}{Decreto de 1938 --- defensa sobre la margen izquierda del río La Caldera (B.O.\ Nº 1724, h.\ 19)}
Vecinos del departamento solicitan «\textit{la construcción de defensas sobre la
margen izquierda del río Caldera}», y el decreto recoge su fundamento:
«\textit{todos los años en la época de lluvias, las aguas del río Caldera
amenazan las propiedades con inminente peligro de perjuicio para los
cultivos}».

De los cómputos métricos, presupuestos y croquis de la obra «\textbf{defensa río
Caldera}», \textbf{proyectada en una longitud de cincuenta metros}, resulta su
costo.
\end{ficha}

\begin{cautela}
\textbf{Es el antecedente más antiguo de la defensa del río que este libro
persigue, y está ochenta años antes del amparo.} El capítulo
\ref{cap:amparo} cuenta que el Plan de Manejo de 2024 cifra la defensa del río
La Caldera en el cincuenta y ocho por ciento de toda su inversión, y que el
aluvión que lo motivó es de diciembre de 2018. \textbf{En 1938 los vecinos ya
pedían lo mismo, con el mismo argumento} ---que todos los años, en la época de
lluvias, el río amenaza las propiedades--- y el Estado provincial lo proyectó.

\textbf{Y la escala es la que hay que mirar.} Aquella obra tenía
\textbf{cincuenta metros}. La del Plan de 2024 se mide en kilómetros y en miles
de millones de pesos. \textbf{Ochenta y seis años entre el primer pedido
documentado y el primer plan cifrado.}

\textbf{Y el barrido de 1940 contesta qué pasó con aquellos cincuenta metros: no
se hicieron.} El decreto del Nº 1833, h.\ 46, pedido «\textit{con carácter de
urgencia}», lo dice entero: «\textit{el año ppdo.\ debieron efectuarse las obras
de defensa en el río de la Caldera, \textbf{las que fueron suspendidas en virtud
de que el mencionado río desvió su curso}}»; y agrega que «\textit{hoy
nuevamente amenazan peligro las aguas de este río \textbf{para la población de la
Calderilla}}». Lo que autoriza es «\textit{la construcción de un \textbf{reparo
provisorio} para defensa en el río de la Caldera en el lugar denominado la
Calderilla}».

\textbf{Tres cosas quedan establecidas, y las tres importan.} La primera: entre
1938 y 1940 hay \textbf{dos intentos de defender el río y ninguna obra
terminada}. La segunda: \textbf{el motivo de la suspensión es el propio río}
---desvió su curso--- y no la falta de fondos ni la desidia, que es el
argumento que un lector supondría. Y la tercera: \textbf{el peligro se corrió de
lugar}, de la margen izquierda en general a \textbf{la población de la
Calderilla} en particular.

\textbf{Eso es exactamente lo que el capítulo \ref{cap:amparo} discute ochenta
años después}, y con el mismo vocabulario: obras que se proyectan, un río que
cambia de curso, y defensas que se ejecutan como reparo provisorio. \textbf{Y el
archivo no registra que el reparo de 1940 se haya hecho tampoco.}

\textbf{Y el croquis, otra vez, se cita y no se publica.} El decreto lo nombra
entre las piezas que fundan el costo, como el de la acequia de Urquiza en 1923 y
como el plano de Hessling en 1911.
\end{cautela}

\subsection*{Y el municipio no pagaba la renta escolar}

\begin{ficha}{Decretos de 1938 sobre el balance municipal de 1937 (B.O.\ Nº 1750 h.\ 12 y Nº 1763 h.\ 6)}
La Comisión Municipal del distrito de La Caldera eleva el balance del ejercicio
1937. El Consejo de Educación informa que esa Comisión «\textit{no ha efectuado
ingreso alguno de contribución al fondo propio de educación}», y el Poder
Ejecutivo \textbf{devuelve el balance con observaciones}.

Un mes más tarde, con nuevo informe de la Comisión, el decreto del expediente
1117 M \textbf{aprueba el ejercicio 1937 y, de acuerdo al monto de los ingresos
habidos, deja establecido el porcentaje} que corresponde a ese fondo.
\end{ficha}

\begin{cautela}
\textbf{Cierra un hilo que este capítulo abrió en 1927 y dejó suelto.} El
decreto 5024 de aquel año aprobó la primera rendición de cuentas del municipio
«\textit{atento el informe del Consejo General de Educación que manifiesta su
conformidad en lo que respecta a la renta escolar que por dicho ejercicio le
corresponde}». Once años después, el mismo Consejo informa lo contrario: que el
municipio \textbf{no aportó nada}.

\textbf{De modo que la renta escolar no era una formalidad: era una
obligación que se controlaba, y que este municipio incumplió.} Y la Provincia
no la ejecuta ni la sanciona: devuelve el balance, pide informe y después fija
el porcentaje sobre lo efectivamente recaudado. \textbf{Sigue sin publicarse
una sola cifra}: ni el balance, ni el monto de los ingresos, ni el porcentaje.
Son el sexto y el séptimo acto económico del municipio publicados sin su
número.

\textbf{Y hay un renglón de 1938 que conviene transcribir aunque no nombre al
departamento en su parte dispositiva} (Nº 1738, h.\ 6): al discutir el
cumplimiento de las comunas de campaña, el acto dice que no lo ha hecho
«\textit{ninguna de las municipalidades de la campaña, a excepción de la de La
Caldera, que ocasionalmente lo ha hecho}». \textbf{Es la única vez en todo el
archivo temprano en que este departamento aparece como la excepción que
cumple}, y aparece en el mismo año en que el Consejo de Educación le observa el
balance por no aportar. Este libro transcribe las dos cosas y no las concilia.
\end{cautela}

\begin{cautela}
\textbf{Y tres piezas más de los dos años, que el apéndice
\ref{cap:cronologia} detalla.}

\textbf{«La Helvecia» era «la Calderilla».} El edicto de deslinde, mensura y
amojonamiento de 1937 (Nº 1691, h.\ 41) individualiza el inmueble
«\textit{denominada «la Helvecia» antes «Calderilla»}», con lindero norte en la
propiedad de \textbf{Augusto Regis}\index{Regis, Augusto} ---presidente de la
Comisión Municipal del departamento en ese mismo año---, al naciente los
herederos de \textbf{Pedro Valencia} y al poniente el \textbf{río de La
Caldera}. \textbf{Los dos nombres están en la lámina \ref{fig:parajesnom} como
topónimos distintos}, y este edicto dice que son el mismo lugar en dos momentos:
la nómina de 1938 los sigue imprimiendo separados, en el mismo circuito.

\textbf{El pueblo tuvo teléfono antes que agua de red.} El decreto de 1937
(Nº 1711, h.\ 5) paga «\textit{las reparaciones de las líneas telefónicas del
pueblo de La Caldera}», con los trabajos a cargo de \textbf{Rafael Delgado},
que ese mismo año es juez de paz propietario del distrito. \textbf{Reparar
supone que existían}, y el archivo no registra cuándo se instalaron.

\textbf{Y la nómina de 1938 resulta ser la de 1940, nombre por nombre.} El
departamento es el \textbf{Colegio Electoral Nº 2}, con los circuitos 5 y 6: la
mesa 1 vota en la Escuela Elemental Mixta con nueve parajes, la mesa 2 en la
Iglesia Parroquial con treinta y uno, y la del circuito 6 en la Escuela
Infantil Provincial de Vaqueros con siete. \textbf{Cuarenta y siete renglones y
cuarenta y seis nombres distintos} ---«Calderilla» se imprime dos veces, en las
mesas 1 y 2---, y \textbf{son exactamente los cuarenta y seis del padrón de
1940}: las cuatro diferencias del cotejo son variantes de grafía del mismo
lugar, «Corral de Barranca» por Barrancas, «Monte de las Garrapatas» por
Monte, «Potrero de Castilla» por Castillo y «Pueblo» por Pueblo (La Caldera).

\textbf{Eso corrige la lectura de la lámina en un punto que importa.} Catorce
nombres ---Helvecia, Getsemaní, Abra de la Sierra, Cañada Ancha, Severino,
Villa Urquiza, Santa Gertrudis, Entre Ríos, Mosquera, Nogalito, Los Matos,
Higueritas, Los Chusos y La Sierra--- figuraban como si aparecieran en 1940 por
primera vez. \textbf{Están en 1938 y son los mismos}, de modo que lo que la
lámina muestra en su última columna no es una ampliación de 1940 sino
\textbf{una lista que ya estaba armada dos años antes y no se tocó}. Y uno de
esos nombres tiene explicación en el mismo año: «La Helvecia» es el nombre
nuevo de «la Calderilla», y las dos figuran en el mismo circuito.

\textbf{Y la segunda concesión de agua del departamento se pide en 1938 y se
resuelve en 1939.} Un particular ---Isaac Fernández--- solicita «\textit{la
concesión de agua}» para regadío de la finca \textbf{«Campo Alegre»}, con título
inscripto «\textit{al folio 374, asiento 475 del libro D de títulos de la
Caldera}», y el trámite corre «\textit{con intervención del señor Fiscal de
Gobierno y de la Municipalidad de la Caldera}» (Nº 1746, h.\ 27 y 28).
\textbf{Catorce años después de la de Urquiza, el mecanismo es idéntico}: la
solicita un propietario, informa el municipio, dictamina el fiscal.

\textbf{Y el acto que lo resuelve, en 1939, trae la pieza que este libro venía
buscando} (Nº 1796, h.\ 23): entre sus considerandos, la oposición al pedido
invoca «\textit{la existencia de una \textbf{ordenanza dictada por el municipio
de la Caldera en fecha octubre 13 de 1914}, que reglamenta la distribución de la
totalidad}» del caudal.

\textbf{Es la ordenanza de 1914 que el capítulo \ref{cap:infraestructura} usa
como prueba de que el municipio decidía sobre el agua de su suelo, y acá aparece
citada veinticinco años después como derecho vigente}, en un expediente
provincial, y esgrimida contra una concesión nueva. \textbf{No es sólo que la
ordenanza existió: es que todavía se invocaba}, y que la Provincia tuvo que
pronunciarse sobre ella. Lo que el archivo no dice es cómo resolvió.

\textbf{Y reaparece el apellido Quipildor, ahora como sucesión y con
superficie.} Un pedimento minero de \textbf{dos mil hectáreas (cuatro unidades)}
se ubica sobre las fincas \textbf{Rosal y Potrero, «propiedad de la sucesión
Quipildor»}, en los departamentos de Rosario de Lerma y Caldera (Nº 1726 h.\ 40
y Nº 1730 h.\ 45). \textbf{Es el mismo apellido que en 1925 aparecía como dueño
de los terrenos donde Larrán pidió su cateo}, y trece años después sigue siendo
propietario en el departamento ---ahora como sucesión, y con dos fincas
nombradas---. Coincidencia de nombre, como siempre.

\subsection*{Y en 1940 el departamento entra en la red vial de la provincia}

El decreto que fija la red vial provincial de 1940 incorpora dos trazas que
pasan por el departamento (Nº 1875, h.\ 21): «\textit{de la Calderilla a río
Saladillo por Gallinato y Campo Santo}», y una ruta que corre «\textit{por el
Jardín, Pampa Grande, Guachipas, La Viña, Cerrillos, Salta, \textbf{Vaqueros,
Calderilla, hasta Tres Cruces en el límite con Jujuy}}».

\begin{cautela}
\textbf{Esa segunda traza es la que hoy se llama Ruta Nacional 9}, y el
decreto la describe atravesando el departamento de sur a norte, por Vaqueros y
la Calderilla, hasta el límite provincial. \textbf{Es la primera vez que el
archivo temprano ubica al departamento dentro de una red y no como punto
aislado}: hasta 1939 los caminos que lo nombran son tramos ---Calderilla al
Desmonte, Calderilla a Campo Santo, Calderilla al río Las Pavas--- y en 1940 es
un pasaje.

\textbf{El capítulo \ref{cap:infraestructura} lee ese pasaje noventa años
después como la forma del problema}: un territorio que la ciudad atraviesa. El
decreto de 1940 lo escribe sin ninguna intención de decirlo.
\end{cautela}

\subsection*{Y Regis es, a la vez, presidente del municipio y diputado provincial}

\begin{cautela}
\textbf{Es la décima superposición de la serie y la más directa de todas.} En
1940, al pedir ciento cincuenta pesos de ayuda para la comisión pro-fiestas
patronales del pueblo, el Boletín identifica al solicitante como
«\textit{el señor presidente de la comisión municipal del distrito de la
Caldera, \textbf{diputado provincial don Augusto Regis}}»\index{Regis, Augusto}
(Nº 1875, h.\ 2).

\textbf{No es un cargo que coincide con otro dentro del departamento ni un
propietario que legisla: es la misma persona presidiendo el municipio y
sentada en la Legislatura que lo regula.} Y no es un año: Regis preside la
Comisión Municipal \textbf{sin interrupción desde 1937}, reelecto por decreto en
1937, 1938, 1939 y 1940 ---cuatro «nuevos períodos legales de funciones»
seguidos---, y es el mismo hombre que el Interventor Nacional nombró a ese
cuerpo en abril de 1922 y que renunció dos veces, en 1922 y en 1923.
\textbf{Dieciocho años en el mismo cuerpo, con dos salidas en el medio}, y al
final presidiéndolo desde una banca.
\end{cautela}

\begin{cautela}
\textbf{Y cuatro cosas más que los dos años dejan, y que el apéndice
\ref{cap:cronologia} detalla.}

\textbf{El cuerpo municipal cambia de nombre.} En 1940 el Boletín ya no habla
de miembros de la Comisión Municipal sino de \textbf{concejales titulares y
suplentes}: por La Caldera, Eugenio Muñoz y José Manuel Alfaro. \textbf{Es la
primera vez en treinta y dos años de archivo} que el cuerpo del departamento se
nombra así.

\textbf{Dos escuelas nacionales más.} La Ley 4874 ---la ley Láinez--- instala
escuela en \textbf{San Alejo} en 1939 y en \textbf{Gallinato} en 1940, después
de la de San Alejo de 1924. \textbf{El Estado nacional pone tres escuelas en el
departamento entre 1924 y 1940}, y la Provincia se limita a consentir cada vez.

\textbf{Un apellido andino preside, por fin, del lado de la mesa.} En la mesa 1
del circuito 5 de 1940 es suplente primero \textbf{Severino Mamani}: es
\textbf{el primer apellido de origen andino que el archivo registra como
autoridad de mesa} del departamento, y Mamani es uno de los seis que el revalúo
de 1944 traerá entre sus titulares.

\textbf{Y un comisario destituido con causa escrita.} En 1939 una información
sumaria comprueba contra \textbf{Abraham Alegre}, comisario de policía de La
Caldera, «\textit{actos graves de inconducta y comisión de hechos}» que el acto
enumera, y lo separa del cargo «\textit{sin perjuicio de las acciones
judiciales a que hubiere lugar}»; lo reemplaza \textbf{Claudio Murúa}. \textbf{De
los más de setenta relevos que este capítulo registra entre 1908 y 1946 ---la
lámina \ref{fig:cargos} los cuenta---, es el único que dice la causa con esas palabras}: los demás usan la fórmula o callan.
\end{cautela}

\textbf{Y las autoridades del bienio se renuevan enteras otra vez.} La
comisaría pasa por \textbf{Ramón Juárez} ---que pide licencia en enero de 1937 y
renuncia en marzo---, \textbf{Fabián Alonso} interino, \textbf{Jorge O.\ Blasco}
y, al cierre de 1938, \textbf{Abraham Alegre}, permutado con el comisario de
San Carlos. La subcomisaría de Vaqueros cambia cuatro veces. Y la elección de
1938 da por el departamento al senador \textbf{Florentino M.\ Serrey}
\index{Serrey, Carlos} por un partido y a \textbf{Francisco Mercado} con
\textbf{Teófilo Reyes} de suplente por el otro, y compone la Comisión Municipal
con \textbf{Eustaquio Murúa}\index{Murúa, Eustaquio} y Eugenio Muñoz como
titulares y \textbf{Genaro Molina} entre los suplentes. \textbf{Murúa y Molina
llevan entonces veintidós y veintinueve años apareciendo en este archivo}, y
Serrey es el apellido que en 1923 era dueño de la finca Vaqueros y que en 1985
todavía figura en la expropiación de la Ley 6334 (capítulo
\ref{cap:expropiacion}).

\textbf{Y los caminos se pagan quincena por quincena.} Los dos años traen
decenas de órdenes de pago de Vialidad por los caminos \textbf{Calderilla--El
Desmonte por Gallinato}, \textbf{Calderilla--Campo Santo} y
\textbf{Calderilla--Güemes por Gallinato}: jornales de cuadrilla, ripio de
Jorge Cornejo, liquidaciones al contratista Adolfo Martínez, capataces con
nombre y apellido. \textbf{Es el primer tramo del archivo en que una obra del
departamento se puede seguir semana a semana}, y queda como serie a reconstruir.
\end{cautela}

\section{1941: un agente muerto en la Calderilla, y tierra que se compra para un camino}

El año se leyó edición por edición y es el vigésimo séptimo tramo: las ediciones
\textbf{1878 a 1929}. Deja, además del cierre de la mina, dos actos que el libro
no tenía.

\begin{ficha}{Expediente de pensión --- B.O.\ Nº 1907, h.\ 5 (1941)}
La viuda de \textbf{don Paulino Fernández}, «\textit{extinto agente de policía de
la comisaría de la Caldera}», pide ser acogida a los beneficios de la ley, y
acredita el vínculo «\textit{en mérito de los testimonios expedidos por la
oficina del registro civil de la Caldera}».

El informe del expediente dice que el causante murió «\textit{[en el] punto
denominado “la Calderilla” de resultas de un tiro de revólver que le
descerrajara \textbf{Abel Murúa}}».
\end{ficha}

\begin{cautela}
\textbf{Es la primera muerte violenta que el archivo temprano registra en el
departamento}, y llega por la vía menos previsible: un trámite de pensión.

\textbf{Lo que el libro afirma es lo que el acto dice, y nada más}: que el
Estado provincial, al resolver una pensión, consignó esa causa de muerte con
ese nombre. \textbf{No es una condena ni consta aquí su resultado judicial}: el
archivo leído no trae la causa penal, y este libro no la busca ni la supone. El
apéndice \ref{cap:pedidos} la reclama, porque es lo único que permitiría decir
algo más.

\textbf{Y el apellido vuelve.} Murúa recorre este capítulo desde 1912 ---jurado
de patentes, comisario del departamento en 1916, 1923, 1930 y 1932, causante del
sucesorio que remata La Calderilla en 1923, titular de La Despensa en 1925,
miembro de la Comisión Municipal en 1938 y comisario otra vez en 1939---.
\textbf{Coincidencia de nombre, como siempre, y en este caso con más razón que
nunca para no inferir nada.}
\end{cautela}

\begin{cautela}
\textbf{Y por primera vez el Estado compra tierra en el departamento para hacer
un camino.} El decreto del Nº 1880, h.\ 8, considera «\textit{imprescindible}»
la adquisición ---o venta--- de «\textit{los terrenos necesarios para la
ejecución de las mejoras del camino Calderilla a río de las Pavas, tramo
Desmonte-Campo Santo}», y encarga la gestión.

\textbf{Hasta 1940 los caminos del departamento aparecen como jornales, ripio y
certificados de obra: nunca como tierra.} Éste es el primer acto en que la traza
exige comprar el suelo por donde pasa, y cae sobre el mismo camino que desde
1937 concentra casi todo el gasto vial del departamento. \textbf{Quién era dueño
de esos terrenos y cuánto se pagó no se publica}, y el apéndice
\ref{cap:pedidos} lo pide: es el antecedente más antiguo de la relación entre
una traza y el dominio del suelo que este libro discute en el capítulo
\ref{cap:expropiacion} para 1985.

\textbf{Y un apellido andino al frente de una cuadrilla.} El mismo año autoriza
la permanencia en obra de la cuadrilla «\textit{a cargo del peón de 1.ª don
\textbf{Manuel Apaza}}». Apaza es otro de los seis apellidos del revalúo de 1944,
y aquí aparece dirigiendo trabajo en el camino.
\end{cautela}

\section{1942 y 1943: la concesión de 1885 aparece entera}

Los dos años se leyeron edición por edición y son los tramos vigésimo octavo y
vigésimo noveno: las ediciones \textbf{1930 a 1982} y \textbf{1983 a 2035}.

\begin{ficha}{Inscripción de derechos de agua --- B.O.\ Nº 1948, h.\ 24 (1942)}
Del registro inmobiliario consta que «\textit{por \textbf{acuerdo de la
Municipalidad de la Caldera de fecha 4 de junio de 1885}, se concedió a don
\textbf{Eustaquio Murúa}}»\index{Murúa, Eustaquio} el uso de \textbf{treinta
litros de agua por segundo del río Wierna}, y que a la finca «\textit{denominada
\textbf{Sausal o Curuzú}, ubicada en el distrito de Vaqueros, departamento de la
Caldera}» le corresponde ese uso por ser su titular sucesor «\textit{en la
posesión}» de los anteriores propietarios.
\end{ficha}

\begin{cautela}
\textbf{El capítulo \ref{cap:planteo} situaba el punto de partida de la voz
municipal sobre el agua en 1885 y no tenía el acto.} Ahora está: \textbf{la
fecha exacta ---4 de junio de 1885---, el beneficiario, el caudal y la finca}.

\textbf{Y el beneficiario es un Murúa}, cincuenta y siete años antes de que el
apellido recorra este capítulo como comisario, jurado de patentes, causante de
sucesorio y titular de La Despensa. \textbf{Es la aparición más antigua de ese
apellido en el departamento} que este libro registra, y llega por un acto del
municipio concediéndole agua.

\textbf{Y hay una segunda, del mismo año, que dice algo distinto} (Nº 1964,
h.\ 22): sobre el \textbf{río de las Nieves, naciente del Wierna}, se reconocen
derechos «\textit{otorgado[s] por la Municipalidad de la Caldera}» y
«\textit{reconocidos como tales \textbf{desde fecha inmemorial}}», para la finca
«Yacones» o «Tránsito». \textbf{«Desde fecha inmemorial» es una fórmula, no una
fecha}, y este libro no la convierte en una: registra que en 1942 la Provincia
inscribió como propios derechos que el municipio había otorgado sin que conste
cuándo.

\textbf{Las dos inscripciones son del mismo año en que el capítulo
\ref{cap:infraestructura} sitúa la pérdida de esa voz.} Y ahora se entiende
mejor por qué: \textbf{1942 es el año en que los derechos que el municipio había
concedido pasan a los registros provinciales}. No se los quitaron: se los
anotaron a nombre de otro registro.
\end{cautela}

\subsection*{Y Regis es, además, comisario de policía}

\begin{cautela}
\textbf{Es la undécima superposición y la que completa el cuadro.} El decreto
del Nº 1952, h.\ 16, de 1942, «\textit{nombra comisario de policía de la
Caldera, con anterioridad al día 1.º del mes en curso, al señor \textbf{Augusto
Regis}}»\index{Regis, Augusto}. En 1943 se le liquidan haberes «\textit{como
comisario de policía de la localidad de la Caldera, durante nueve días del mes
de setiembre}», y después \textbf{renuncia} al cargo (Nº 2021, h.\ 14).

\textbf{En el mismo bienio es reelecto presidente de la Comisión Municipal}, por
decreto, en 1942 y en 1943 ---sexto y séptimo «nuevo período legal de
funciones» consecutivos---. De modo que Augusto Regis es, a la vez y en actos
publicados: \textbf{presidente del cuerpo que gobierna el departamento,
diputado provincial y comisario de policía del mismo departamento}.

\textbf{No hace falta inferir nada.} Este capítulo viene registrando desde 1909
que el que manda la policía integra el cuerpo municipal ---Aparicio\index{Aparicio, Escolástico M.}
en 1910, Rauch en 1922---. \textbf{En 1942 las dos cosas son la misma persona y
además legisla.} Es el punto más alto de una serie de treinta y tres años, y el
archivo lo publica sin comentario.
\end{cautela}

\begin{cautela}
\textbf{Y en 1943 la Provincia regulariza por prescripción el terreno de su
propia comisaría.} Un edicto anuncia «\textit{juicio de posesión treintañal de un
inmueble ubicado en el pueblo de la Caldera, donde se encuentra emplazado el
edificio ocupado por la \textbf{comisaría de policía}}», \textbf{deducido por el
Gobierno de la Provincia de Salta}, con cuarenta metros de frente por doscientos
de fondo, lindando al oeste con «\textit{la calle principal del pueblo}» y con
\textbf{Daniel Linares}\index{Linares, Daniel}; y el auto ordena oficiar
«\textit{a la Municipalidad de la Caldera para que informe si afecta propiedad
municipal}» (Nº 1994, h.\ 48, y Nº 2001, h.\ 22).

\textbf{Treinta años de posesión quiere decir que el Estado ocupaba ese terreno
desde 1913 o antes, sin título.} Y quien tiene que informar si le pertenece es
el municipio, que en ese momento preside el hombre que acababa de ser comisario
en ese mismo edificio. \textbf{El archivo no dice cómo terminó}, y el apéndice
\ref{cap:pedidos} lo pide.
\end{cautela}

\section{1944: el revalúo en su sitio, y otra comuna intervenida}

El año se leyó edición por edición y es el trigésimo tramo. Trae dos cosas que
corrigen la lectura de piezas que el libro ya usaba.

\begin{cautela}
\textbf{El revalúo de 1944 no es un padrón del departamento: es un padrón
provincial en el que el departamento ocupa un tramo.} El barrido lo ubica en el
Boletín \textbf{Nº 2166}, y el cuadro corre por lo menos de la hoja 6 a la 76,
con números de partida que pasan de los cinco mil. \textbf{El renglón
«\textit{departamento de la Caldera}» abre una sección dentro de esa serie}
(h.\ 6).

\textbf{Eso confirma la cautela que el capítulo \ref{cap:fincas} ya traía} ---que
no se sabe si las cuarenta y siete partidas publicadas son todas las del
departamento o sólo las revaluadas--- \textbf{y la vuelve más precisa}: lo que
el Boletín publica es un revalúo general de la Provincia, y el tramo del
departamento es lo que cabe en su sección. \textbf{La pregunta ya no es si el
padrón está completo, sino qué criterio decidió qué partidas entraban}, y el
apéndice \ref{cap:pedidos} la pide así.

\textbf{Y un apellido del cuadro merece registrarse}: entre los titulares del
tramo departamental figura \textbf{«Caldera, Dominga M.\ de»}, con quinientos
pesos de valuación. \textbf{Es el nombre del departamento como apellido de una
titular}, y este libro no infiere nada de eso salvo que conviene no confundirlo
en una búsqueda.
\end{cautela}

\begin{ficha}{Decreto de 1944 sobre el presupuesto municipal (B.O.\ Nº 2053, h.\ 4)}
«\textit{Visto este expediente en el que el señor \textbf{Interventor de la
comuna de la Caldera} eleva a conocimiento y aprobación del Poder Ejecutivo el
proyecto de ordenanza}»: se aprueba «\textit{el cálculo de recursos y
presupuesto de gastos de la \textbf{Municipalidad de la Caldera}, para regir
durante el ejercicio económico 1944 en curso}», y se devuelve el expediente
6463-M.
\end{ficha}

\begin{cautela}
\textbf{Es la tercera intervención del cuerpo municipal que este capítulo
registra}, después de la provincial de 1923 y la federal de 1930, y la primera
en que el interventor aparece ejerciendo la función ordinaria: elevar el
presupuesto. \textbf{El decreto no dice cuándo se dispuso la intervención ni
quién es el interventor}, y ningún acto del año leído lo nombra: \textbf{el
archivo registra el efecto y no el acto}, que es el mismo mecanismo que el
capítulo \ref{cap:opacidad} describe.

\textbf{Y el cuerpo cambió de nombre otra vez.} En 1940 eran concejales, en 1942
y 1943 volvió a ser Comisión Municipal, y en 1944 el mismo decreto lo llama
\textbf{«comuna»} en el vistos y \textbf{«Municipalidad»} en el artículo.
\textbf{Cuatro nombres para el mismo cuerpo en cuatro años}, y este libro los
transcribe sin normalizar.
\end{cautela}

\begin{cautela}
\textbf{Y la comisaría del pueblo se refacciona el año siguiente al juicio por
su terreno.} En 1943 la Provincia promovió posesión treintañal del inmueble
«\textit{donde se encuentra emplazado el edificio ocupado por la comisaría de
policía}»; en 1944 se licitan y ejecutan las obras de refección de ese local,
con presupuesto de \textbf{\$4.363,50}, contratista \textbf{Francisco Crescini} y
certificado final de \textbf{\$5.018,02}, más un adicional (Nº 2048 h.\ 48,
Nº 2053 h.\ 18, Nº 2094 h.\ 6 y Nº 2112 h.\ 7).

\textbf{Primero el título, después la obra}, y en ese orden. Es la primera vez
en treinta y seis años de archivo que el departamento tiene una obra pública
propia seguida paso a paso ---licitación, adjudicación, certificado, adicional
y devolución del depósito en garantía--- \textbf{y no es una escuela, ni un
camino, ni una defensa del río: es la comisaría}.
\end{cautela}

\section{1945: la defensa del río, tercer intento}

El año se leyó edición por edición y es el trigésimo primer tramo. Su hallazgo
cierra un hilo que este capítulo viene siguiendo desde 1938.

\begin{ficha}{Las obras de defensa en el río de la Caldera, 1945}
\textbf{Julio:} la Dirección General de Hidráulica llama a \textbf{licitación
pública} para «\textit{las obras de defensa a efectuarse en el río de la
Caldera}», con presupuesto de \textbf{\$2.515} y sobre el expediente 18385. El
aviso se repite en dieciséis ediciones seguidas.

\textbf{3 de septiembre:} el decreto 8549 autoriza la ejecución.

\textbf{Noviembre:} un decreto resuelve «\textit{ejecutar \textbf{por vía
administrativa} las obras de defensa en el río de la Caldera}» (Nº 2429,
h.\ 5).
\end{ficha}

\begin{cautela}
\textbf{Es el tercer intento en siete años, y el primero que llega a llamarse a
licitación.} En 1938 el proyecto era de cincuenta metros y quedó en croquis; en
1940 la obra «\textit{fue suspendida en virtud de que el mencionado río desvió
su curso}» y se autorizó un reparo provisorio; en 1945 hay presupuesto, llamado
público y expediente.

\textbf{Y termina como los otros dos: sin licitación.} «Por vía administrativa»
quiere decir que el Estado la hace por su cuenta, sin contratista, y el archivo
no publica después ni certificado ni recepción ---a diferencia de la refección
de la comisaría de 1944, que sí dejó licitación, adjudicación, certificado
final, adicional y devolución del depósito---. \textbf{Del edificio de la
comisaría queda el expediente completo; de la defensa del río, sólo la
autorización.}

\textbf{Y la escala cierra el argumento.} 1938: cincuenta metros. 1945:
\textbf{dos mil quinientos quince pesos}, cuando la refacción de la comisaría
del año anterior costó \textbf{cinco mil dieciocho}. \textbf{La defensa del río
valía la mitad que pintar y arreglar la comisaría}, y el capítulo
\ref{cap:amparo} cuenta que ochenta años después esa misma defensa se cifra en
el cincuenta y ocho por ciento de un plan de veintitrés mil millones.
\end{cautela}

\begin{cautela}
\textbf{Y dos actos de escuela que miden población.} El Consejo pide trasladar
la escuela Nº 250 de \textbf{Wierna} «\textit{por encontrarse despoblado dicho
lugar}», a «\textit{San Francisco, partido de Yacones}»; y aprueba instalar una
escuela nacional de la Ley 4874 en \textbf{San Alejo}, «\textit{lugar donde
según censo practicado existe una población de \textbf{51 niños} en edad
escolar}».

\textbf{Es el primer dato de población de un paraje del departamento que el
archivo temprano publica}, y no viene de un censo general sino de un trámite de
escuela: cincuenta y un chicos en San Alejo en 1945. \textbf{Y en el mismo año
Wierna se declara despoblado.} Los dos parajes están en las doce nóminas de la
lámina \ref{fig:parajesnom}; uno tiene escuela y el otro deja de tenerla.
\end{cautela}

\section{1946: el agua se reparte por zonas, y el municipio colabora}

El año se leyó edición por edición y es el trigésimo segundo tramo. Era el último de
esta ventana hasta que se incorporaron 1947 y 1948; \textbf{es el último leído con lectura sobre la imagen}, y en ese sentido cierra la serie verificada.

\begin{ficha}{Distribución de caudales del río Mojotoro (B.O.\ Nº 2686, h.\ 5 y 6)}
Considera «\textit{que el río Mojotoro forma su caudal con el proveniente de los
ríos \textbf{Wierna, la Caldera y Castellanos}}», y divide el sistema en zonas
para el estudio de la distribución: la \textbf{zona II} «\textit{abarca desde la
confluencia de los ríos \textbf{la Caldera y Vaqueros} hasta aguas arriba de la
primera toma del sistema}».

Dispone que \textbf{cuando el caudal de la zona I se mantenga inferior a
\textbf{2.500 litros por segundo}} se prohíban determinados usos, y su
\textbf{artículo 5º} establece que «\textit{las municipalidades de \textbf{Campo
Santo y de la Caldera} prestarán a la Dirección General de Hidráulica la
\textbf{colaboración necesaria}}».
\end{ficha}

\begin{cautela}
\textbf{Es el final del arco que este capítulo abrió en 1885, y lo dice en una
línea.} En 1885 la Municipalidad de La Caldera \textbf{concedía} treinta litros
por segundo del Wierna. En 1914 \textbf{reglamentaba} la distribución de la
totalidad del caudal por ordenanza. En 1929 una ley le dio a Campo Santo la
facultad de \textbf{vigilar y prohibir} bocas-tomas desde las juntas del Wierna.
En 1931 los dos municipios integraron una comisión \textit{ad hoc}. En 1942 los
derechos que el municipio había otorgado pasaron a los registros provinciales.
\textbf{Y en 1946 la Dirección General de Hidráulica fija zonas y umbrales, y a
los dos municipios les asigna «la colaboración necesaria».}

\textbf{Conceder, reglamentar, integrar una comisión, colaborar.} Sesenta y un
años, y el verbo cambia cuatro veces en la misma dirección. \textbf{Ninguno de
los actos deroga al anterior}: la facultad no se quita, se diluye.

\textbf{Y el umbral es la primera cifra hidrológica del departamento que el
archivo publica}: dos mil quinientos litros por segundo en la zona I. El
capítulo \ref{cap:ribera} trabaja con las series de aforo desde 1944; ésta es la
primera vez que un acto fija un número y le cuelga una prohibición.
\end{cautela}

\begin{cautela}
\textbf{Y el municipio publica un balance con cifras, por primera vez desde
1908.} El Boletín Nº 2473 trae el «\textit{Balance del movimiento de la renta}»
de la \textbf{comuna del departamento Caldera}, cuarto trimestre de 1945, con
sus renglones de entradas y salidas, sus saldos mensuales ---\$98,15 al cerrar
octubre, \$98,02 en noviembre--- y la firma de \textbf{F.\ Mercado} al pie de
cada mes.

\textbf{Rompe una serie de treinta y ocho años.} Entre 1922 y 1944 este capítulo
registró \textbf{siete actos económicos del municipio publicados sin una sola
cifra}: presupuestos aprobados por remisión al expediente, rendiciones
aprobadas sin transcribir, balances devueltos con observaciones. \textbf{En 1946
se publica el balance mismo}, y lo que muestra es un municipio que mueve
\textbf{menos de cien pesos por mes}.

\textbf{Y lo publica un interventor.} Francisco Mercado firma ese balance como
interventor de la comuna, y en abril de 1946 renuncia al cargo; el 6 de agosto
un decreto \textbf{da por terminada la intervención} a la Municipalidad de La
Caldera y designa presidente de la Comisión Municipal a Enrique Alejandro
Pfister. \textbf{La tercera intervención del cuerpo duró al menos tres años}
---estaba vigente en 1944 y el archivo no dice cuándo empezó--- y el único
balance con cifras de todo el período lo firma el interventor.
\end{cautela}

\begin{cautela}
\textbf{Y un falso positivo que conviene dejar escrito porque puede repetirse.}
En 1946 se autoriza instalar una escuela de la Ley 4874 «\textit{en la localidad
de \textbf{la Calderilla (departamento de la Capital)}}». \textbf{Hay una
Calderilla en el departamento Capital}, distinta de la de este departamento, y
el barrido la trae porque el término es el mismo. Este libro no la cuenta, y la
declara para que nadie la sume por descuido.
\end{cautela}

\section{1947: el municipio es el departamento entero}

\begin{figure}[p]
\centering
\includegraphics[width=\textwidth,height=0.88\textheight,keepaspectratio]{fig-defensas-1947.png}
\caption[Las defensas contra crecientes del presupuesto provincial de 1947]{\textbf{Las defensas
contra crecientes del presupuesto provincial de 1947: tres de los nueve ríos son de este
departamento, y se llevan el 35,7 por ciento.} El Ítem 6 de la Ley 834 ---«\textit{Obras de
defensas permanentes o eventuales para contener daños por crecientes o inundaciones sobre}»---
reparte \$560.000 entre nueve ríos y un renglón de imprevistos. \textbf{El Mojotoro, el de La
Caldera y el Wierna suman \$200.000}, y son un solo sistema: el acto de 1946 sobre distribución de
caudales establece que el Mojotoro forma el suyo con el Wierna, el de La Caldera y el Castellanos.
\textbf{Presupuestado no es ejecutado}: el barrido de 1947 y 1948 no encontró después ni licitación
ni adjudicación ni certificado de estas obras, y las fojas 12 a 18 del plan hidráulico 1948--1950
---que el decreto 5466-E aprueba sin publicar--- son la pieza que podría contestarlo. Para
comparar, la defensa que sí se ejecutó, por vía administrativa en 1945, costó \$2.515. Fuente:
Ley 834 de Presupuesto para el ejercicio 1947, separata del B.O.\ del 19 de mayo de 1947,
págs.\ 141--142, verificada sobre el facsímil. Elaboración propia; dibujo de
\texttt{scripts/grafico\_defensas.py} sobre \texttt{datos/fiscal/} del repositorio
\texttt{mapas\_caldera} (apéndice \ref{cap:fuentes}).}
\label{fig:defensas}
\end{figure}


El año se leyó edición por edición y es el \textbf{trigésimo tercer tramo}: las
ediciones \textbf{2869 a 3026}, \textbf{158 archivos} y \textbf{2.585 hojas}, barridas en
septiembre de 2026. Con él y con los tres siguientes ---1948, 1949 y 1950---, el barrido que se cerraba en 1946 se
extiende \textbf{hasta 1950}, aunque la ventana en que un cero es un resultado sigue cerrándose en 1946 (apéndice \ref{cap:fuentes}),, y el hueco de sesenta y seis años que este libro declaraba
entre 1947 y 2012 pasa a ser de sesenta y dos.

\begin{cautela}
\textbf{Todo lo que sigue de 1947 y de 1948 sale del barrido automático y no se miró en la
imagen.} El propio informe lo dice en su encabezado. La confianza media del reconocimiento
es de 62 puntos en 1947 y de 66,5 en 1948, con noventa y ciento ochenta y siete hojas
marcadas para mirar enteras, y \textbf{ninguna hoja de ninguno de los dos años alcanza el
umbral de publicable}. De modo que las cifras, los nombres y las fechas de estos dos
apartados \textbf{son lectura de capa de texto y están sujetos a verificación sobre el
facsímil}, y el apéndice \ref{cap:pedidos} la reclama. Las tapas de 1947 agravan la
reserva: el lector automático devuelve números de edición falsos ---«Nº 1805», «Nº 4780»,
«Nº 28911»--- en una parte grande del año, y una tapa imprime «4947».
\end{cautela}

El hallazgo mayor del año no es un acto sobre el departamento sino una \textbf{definición
de qué es el departamento}.

\begin{ficha}{\textbf{Ley Nº 853}, delimitación territorial de los distritos municipales (B.O.\ Nº 2912, 7 de agosto de 1947, h.\ 4 y 5)}
\textbf{Verificada sobre el facsímil.} No es un decreto: es una ley sancionada por el
Senado y la Cámara de Diputados, y su artículo 1º fija «\textit{la delimitación territorial
de los distritos municipales de la Provincia}» entera, departamento por departamento.

El del departamento se enuncia en una sola línea: «\textit{El Municipio de La Caldera
abarca todo el territorio del Departamento de su mismo nombre}». \textbf{La misma fórmula,
palabra por palabra, se usa para Cachi, para Cafayate y para San Antonio de los Cobres}
---éste sobre el departamento de Los Andes---: el departamento de un solo municipio es una
categoría de la ley de 1947, no una singularidad de La Caldera.

Los municipios vecinos se describen \textbf{contra} él, y con los mojones de siempre:
\textbf{San Lorenzo} linda «\textit{por el Norte: Con el Departamento de La Caldera}» y
«\textit{por el Este: Desde el río Vaqueros}»; el \textbf{municipio de Salta},
«\textit{por el Norte: Con los Departamentos de La Caldera y Campo Santo}»; \textbf{El
Bordo}, «\textit{Con el departamento de La Caldera y con la Provincia de Jujuy}», y su
límite corre «\textit{al río Mojotoro en el esquinero Norte-Este de la finca Entre Ríos}»;
y \textbf{Güemes}, al norte, con el departamento y con Jujuy.
\end{ficha}

\begin{cautela}
\textbf{En 1947 hay un solo municipio en el departamento y coincide con él}, y la ley lo
dice con la fórmula que reserva para los departamentos que no se subdividen. Vaqueros
---que en 1948 es un circuito electoral propio y en 1926 ya había pasado al Colegio
Electoral Nº 6 (capítulo \ref{cap:siglo})--- no es municipio: es territorio del municipio de
La Caldera. El segundo municipio aparece recién con su ley de creación de 1970, y el
capítulo \ref{cap:cot} documenta que el límite entre los dos es el río Wierna.

\textbf{Y los mojones con que el acto describe a los vecinos son los mismos que este libro
viene siguiendo desde 1909}: el río Vaqueros, el río Mojotoro y la finca \textbf{Entre
Ríos} ---la «Vaqueros o Entre Ríos» que da nombre al municipio y cuyo remate de 1935 el
apéndice \ref{cap:pedidos} reclama---. Un acto que sólo pretende trazar límites
administrativos vuelve a usar, cuarenta años después, la toponimia de las fincas.
\end{cautela}

\subsection*{El agua de La Helvecia, y quién la pide}

Cinco ediciones seguidas de diciembre de 1947 ---las \textbf{3019 a 3023}--- publican el
mismo edicto: el pedido de reconocimiento de un derecho de concesión de \textbf{cincuenta
litros por segundo}, «\textit{a derivar del río de la Caldera}», para regar la finca
\textbf{«La Helvecia»}, «\textit{ubicada en el departamento de la Caldera, en una extensión
aproximada de 24 hectáreas}». \textbf{Esos cincuenta litros son lo que se pide, no lo que
se otorga}: el decreto que cierra el trámite cuatro meses después reconoce
\textbf{12,6}, porque cincuenta está «\textit{por sobre a lo que tendría derecho las 24
Has.}». \textbf{El edicto publica el pedido y el decreto publica la cuarta parte}, y entre
uno y otro hay cuatro meses en los que sólo el edicto está a la vista.

\textbf{La Helvecia es la finca que antes se llamaba «Calderilla»}, y este libro ya la
tenía: su expediente de mensura de 1937, con perito Jorge de Bancarel y por auto del Juzgado
de 1ª Nominación en lo Civil, tiene por límite occidental el río La Caldera y por lindero
norte a quien veinte días después presidiría el municipio (capítulo \ref{cap:fincas} y
apéndice \ref{cap:pedidos}). \textbf{Diez años después, la misma finca pide el agua del mismo
río que la limita.}

\section{1948: el agrimensor que midió el departamento pide su agua}

El año se leyó edición por edición y es el \textbf{trigésimo cuarto tramo}: \textbf{248
archivos} y \textbf{4.674 hojas}, desde la edición \textbf{3060}. Es el tramo más voluminoso
de toda la serie temprana.

\begin{ficha}{\textbf{Decreto Nº 9041-E} del 5 de abril de 1948, expediente 4887/A/1948 (B.O.\ Nº 3101, h.\ 4)}
\textbf{Verificado sobre el facsímil.} «\textit{Visto este expediente en el que el señor
\textbf{Hermann Pfister}, solicita reconocimiento de derecho al uso de agua pública por una
dotación de \textbf{50 l/s}, a derivarse del Río de La Caldera, para irrigar la finca
denominada \textbf{«La Helvecia»} con una superficie total de \textbf{286 Ha.\ 8.608 m2}, con
\textbf{24 Ha.\ regables}, ubicada en el Departamento de La Caldera}».

\textbf{El municipio, textual}: «\textit{Que la Municipalidad de La Caldera a pedido de
Administración General de Aguas de Salta, en cumplimiento del inciso a) del Art.\ Nº 350,
expresa en el informe respectivo, "que no tiene observación alguna que formular a la
solicitud de reconocimiento que motiva las presentes actuaciones"}».

\textbf{El cálculo}: la dotación pedida «\textit{está por sobre a lo que tendría derecho las
24 Has.}» según los artículos 47 y 48 del Código de Aguas, y a fojas 82 se fija que
«\textit{le correspondería un caudal de \textbf{12,6 l/seg}, sin perjuicio de reajustes
posteriores, cuando se conozca en forma cierta, el caudal del citado río en tiempo de
estiaje}».

\textbf{Artículo 1º}: «\textit{Otórgase un reconocimiento de concesión por usos y costumbres
de agua pública, por un caudal de \textbf{12,6 l/seg}, para irrigar la finca denominada "La
Helvecia" con una superficie bajo riego de 24 Has.\ \ldots{} con carácter de \textbf{temporal y
permanente}, concordante con lo previsto en el Art.\ \textbf{254} del Código de Aguas y a
derivar del Río La Caldera}». El \textbf{artículo 2º} lo otorga con las reservas de los
artículos \textbf{17 y 232}. Firman \textbf{Lucio A.\ Cornejo} y \textbf{Juan W.\ Dates}.

\textbf{Y el edicto tiene su autorización}: por resolución \textbf{Nº 722 del 19 de noviembre
de 1947} del H.\ Consejo de la Administración General de Aguas se procedió a la publicación
de edictos «\textit{en los diarios de esta capital}», «\textit{sin que ello hubiera dado
lugar a oposición de terceros}».
\end{ficha}

\begin{cautela}
\textbf{Dos cosas, y las dos importan.}

\textbf{La primera es quién pide.} \textbf{Hermann Pfister} es el perito que deslindó este
departamento: la \textbf{Loma Bola} en 1916, \textbf{Santa Mónica y Quitilipe} en 1916 ---el
segundo deslinde, cuatro años después del de 1912--- y el \textbf{Potrero Gallinato} en 1917,
las tres filas del apéndice \ref{cap:dominio} para esos años. \textbf{Treinta y dos años
después de medir las fincas del departamento, aparece pidiendo agua para una de ellas.}
Y en el mismo año, por el \textbf{Decreto Nº 8278-G del 14 de febrero de 1948} ---verificado
sobre el facsímil---, «\textit{habiendo terminado su período legal de funciones el señor
Presidente de la H.\ Comisión Municipal del Distrito de La Caldera}», el gobernador
«\textit{nómbrase, en carácter de reelección, a don \textbf{ENRIQUE ALEJANDRO PFISTER},
Presidente de la Honorable Comisión Municipal del Distrito de LA CALDERA, por un período
legal de funciones}», con la facultad del \textbf{artículo 178 de la Constitución} y del
\textbf{artículo 35 de la Ley Nº 68 de Organización y funcionamiento de las
Municipalidades}. Firman \textbf{Lucio A.\ Cornejo} y \textbf{Julio Díaz Villalba}.

\textbf{Y no es un decreto solo: es una tanda.} La misma página publica los decretos
\textbf{8274-G a 8279-G}, todos del 14 de febrero, todos con la misma fórmula y la misma
facultad, sobre \textbf{Angastaco} (Celestino Correa), \textbf{Molinos} (Anacleto
Pastraña), \textbf{Aguaray} (Lauro E.\ Roman), \textbf{El Quebrachal} (Pedro J.\ López) y
\textbf{La Caldera}. \textbf{El presidente del municipio no se elige: lo nombra el
gobernador}, y lo nombra el mismo día que a los de media provincia. La palabra
«\textit{reelección}» del decreto no designa una elección repetida sino una designación
renovada, y el capítulo \ref{cap:politica} muestra que esa distinción recién se cierra en
1983. \textbf{No es un nombre nuevo}: es el mismo a quien el decreto del 6
de agosto de 1946 designó presidente al dar por terminada la tercera intervención del
municipio (apéndice \ref{cap:cronologia}). \textbf{De modo que en 1948 el municipio lo
preside un Pfister y el agua del río la pide un Pfister, y el acto que reconoce el agua
registra al municipio interviniendo en el trámite.} \textbf{Y en diciembre vuelve a aparecer, en otra lista.} El Honorable Tribunal Electoral
publica la «\textit{lista de candidatos a \textbf{constituyentes} propuestos por el PARTIDO
PERONISTA de Salta en los comicios del \textbf{5 de diciembre de 1948}}», dos por
departamento, y por el nuestro van «\textit{\textbf{Alberto Ovejero Paz y Alejandro
Pfister}}» (B.O.\ Nº 3282, h.\ 16, verificado sobre el facsímil).

\begin{cautela}
\textbf{El mismo apellido preside el municipio por decreto y se postula a constituyente por
elección, en el mismo año.} \textbf{Enrique Alejandro Pfister} es designado presidente de la
Comisión Municipal en agosto de 1946 al terminar la tercera intervención, \textbf{reelecto por
decreto del gobernador en febrero de 1948} ---en la tanda de seis municipios del mismo día---
y en diciembre de ese año figura como \textbf{candidato a constituyente} por el departamento.
\textbf{Los dos cargos que puede ocupar una persona en el departamento se resuelven por vías
opuestas}: el municipal, por designación del Ejecutivo provincial; el constituyente, por
voto. Y en 1948 el mismo nombre está en los dos.

\textbf{La convención que esos comicios eligen es la que sanciona la Constitución provincial
de 1949}, y el capítulo \ref{cap:politica} sigue desde ahí la historia de las comisiones
municipales hasta 1983. \textbf{El departamento llega a esa convención con dos candidatos y
sin municipio electo}, que es la situación que la propia convención venía a discutir.

\textbf{Y en la misma lista hay un nombre que este capítulo acaba de citar}: \textbf{Dantón
Cermesoni}, candidato por Chicoana, es uno de los tres que firman en Acuerdo de Ministros la
insistencia que permitió pagar la obra sanitaria de la escuela del pueblo. \textbf{Coincidencia
de nombre, como siempre}, y anotada porque el archivo la ofrece en dos hojas del mismo año.
\end{cautela}

\textbf{Este libro no afirma que sean la misma persona ni que sean parientes}: son
coincidencias de apellido, con la regla que el capítulo \ref{cap:metodo} fija para todas.
Lo que sí queda registrado es que el apellido del agrimensor que midió el departamento en
1916 y 1917 reaparece en 1948 en los dos lugares donde se decide sobre él: pidiendo agua y
presidiendo el municipio. \textbf{Es el mismo patrón que este capítulo documenta para 1909},
cuando el que clasificaba el impuesto mandaba la policía; la diferencia es que ahora el
archivo no permite cerrar la identificación, y por eso el apéndice \ref{cap:pedidos} pide el
expediente.

\textbf{La segunda es qué hace el municipio.} El acto lo registra actuando \textbf{a pedido de
la Administración General de Aguas}, en cumplimiento de algo que la lectura de capa no
alcanza a leer. No concede: informa o verifica cuando el organismo provincial se lo pide. Es
exactamente el registro del artículo 5º del acto de 1946, donde las municipalidades
«\textit{prestarán a la Dirección General de Hidráulica la colaboración}» que se les
requiera. \textbf{La serie que el capítulo \ref{cap:planteo} arma sobre la pérdida de la voz
municipal ---conceder en 1885, proponer en 1918, inscribir en la provincia en 1942--- gana
un cuarto tiempo y es el más bajo: colaborar a pedido, en 1946 y en 1948.}
\end{cautela}

\subsection*{El Angosto también pide agua}

Siete ediciones de 1948 ---de la \textbf{3126} a la \textbf{3132}--- publican el edicto
\textbf{Nº 3720}, verificado sobre el facsímil: por el artículo 350 del Código de Aguas se hace
saber que «\textit{el señor \textbf{Daniel Linares}, solicitando en expediente Nº 4498/47,
reconocimiento de concesión de agua pública a derivarse del río La Caldera, para regar su
propiedad denominada \textbf{«El Angosto»}, ubicada en el Departamento de La Caldera, en una
superficie aproximada de \textbf{29 hectáreas}}». El reconocimiento sería «\textit{por una
dotación de \textbf{15,22 litros por segundo}, con carácter de \textbf{temporal y eventual}}».

\begin{cautela}
\textbf{Los dos actos publican la superficie y el caudal, y con eso el criterio queda a la
vista.} Pfister pide para \textbf{24 hectáreas} y se le reconocen \textbf{12,6 litros por
segundo}; Linares pide para \textbf{29} y se le reconocen \textbf{15,22}. \textbf{Las dos
razones dan lo mismo: 0,525 litros por segundo por hectárea.} No es una coincidencia
aproximada ---29 por 0,525 da 15,225--- sino un coeficiente aplicado.

\textbf{Este libro sostiene que la opacidad de este dispositivo opera en el nivel del
criterio y no del acto, y éste es el contraejemplo que hay que poner.} Acá el criterio no se
enuncia en ningún lado ---ni el edicto ni el decreto dicen «0,525»--- pero \textbf{se deduce
del par}, porque los dos actos publican los dos términos de la división. \textbf{Publicar la
superficie junto al caudal vuelve verificable la regla sin necesidad de enunciarla}, y es
exactamente lo que ochenta años después los instrumentos del capítulo \ref{cap:opacidad} dejan
de hacer.

\textbf{Lo que no es igual es el carácter}: el de Pfister se reconoce «\textit{temporal y
permanente}» y el de Linares «\textit{temporal y eventual}». \textbf{Misma agua, mismo río,
mismo coeficiente, y distinta seguridad jurídica}, sin que ningún acto diga por qué. El
apéndice \ref{cap:pedidos} lo pregunta.

\textbf{Y un tercer edicto de la misma hoja marca el límite del coeficiente}: Ricardo
Echenique pide para «\textit{La Ciénega}», 120 hectáreas en el departamento Capital, y obtiene
\textbf{14,66 litros por segundo} ---0,12 por hectárea---. \textbf{El coeficiente no es
provincial: es del río}, o del sistema. Pero deriva de la \textbf{acequia Isasmendi} del
sistema de riego de La Silleta y no de un río, de modo que este libro no los compara: registra
que dentro del río La Caldera la razón es la misma en los dos casos.
\end{cautela}

\textbf{El Angosto es la finca que Antonia P.\ de Peralta hizo deslindar en 1921}, y cuyos
cuatro linderos repiten tres de los que el remate del Angosto de Arrieta declaraba en 1909
(capítulos \ref{cap:fincas} y \ref{cap:siglo}). \textbf{Dos de las fincas del archivo
temprano piden agua del río en el mismo año}, y las dos por la vía del reconocimiento de un
uso anterior y no por concesión nueva: es el mecanismo del artículo 254 del Código de Aguas,
que convierte la costumbre en derecho inscripto. \textbf{Lo que en 1914 repartía una
ordenanza municipal, en 1948 lo reconoce un decreto provincial finca por finca.}

\subsection*{Y cuatro actos más, que miden la escala}

\begin{figure}[p]
\centering
\includegraphics[width=\textwidth,height=0.88\textheight,keepaspectratio]{fig-coparticipacion-1947.png}
\caption[El criterio de la coparticipación municipal, 1947]{\textbf{El criterio de la
coparticipación municipal, 1947: el porcentaje no es una asignación, es el espejo de lo que cada
municipio ya cobraba.} Las dos columnas son la misma columna. El artículo 6º de la Ley 834 manda
liquidar la participación en los impuestos a los Réditos, Ventas, Beneficios Extraordinarios y
Contribución Territorial «\textit{en proporción a los recursos efectivamente percibidos por cada
una de ellas durante el año 1946}», de modo que \textbf{el que recauda poco recibe poco, y lo
recibe poco porque recauda poco}. La Caldera recaudó \$3.068 ---el 0,24 por ciento de lo que
recaudó la capital--- y le corresponde el \textbf{0,127 por ciento}: \$266,41 de los \$209.775,06
repartidos. La línea punteada marca dónde caería un reparto parejo entre las cuarenta y dos
municipalidades. \textbf{Muestra: los órdenes 1 a 32, que es hasta donde llega la hoja publicada;
los diez que faltan están todos por debajo de La Caldera.} Escala logarítmica en los dos paneles.
Fuente: Decreto 5276-E del 30 de julio de 1947, expediente 1873/C/1947, B.O.\ Nº 2909, h.\ 5,
verificado sobre el facsímil. Elaboración propia; dibujo de
\texttt{scripts/grafico\_coparticipacion.py} sobre \texttt{datos/fiscal/} del repositorio
\texttt{mapas\_caldera} (apéndice \ref{cap:fuentes}).}
\label{fig:coparticipacion}
\end{figure}

\begin{ficha}{1947--1948, lo que el Estado pone y lo que cobra}
\textbf{La escuela, y lo que costó pagarla.} El \textbf{decreto 12.581 del 19 de noviembre
de 1948} aprueba y adjudica «\textit{a la firma \textbf{Conrado Marcuzzi} la licitación de la
obra "Refecciones y ampliaciones en las instalaciones sanitarias de la \textbf{Escuela
Provincial de La Caldera}", por un importe total de \textbf{\$35.026,99}}». Ocho días
después, «\textit{atento a las observaciones formuladas por \textbf{Contaduría General}}», el
\textbf{decreto 12.744-E} ---dictado \textbf{en Acuerdo de Ministros} y firmado por tres---
\textbf{insiste} en el cumplimiento del anterior (B.O.\ Nº 3285, h.\ 6, verificado sobre el
facsímil). Y en diciembre de 1947 un acto autoriza al \textbf{Banco de Préstamos y Asistencia Social} a acordar «\textit{un subsidio extraordinario en la suma de \textbf{CIEN PESOS} a cada una de las \textbf{cooperadoras} de las escuelas nacionales y provinciales}» de toda la provincia. \textbf{El departamento entra con tres} (B.O.\ Nº 3022, h.\ 6, verificado sobre el facsímil): la \textbf{nacional Nº 51, Calderilla}, y las \textbf{provinciales Nº 21, La Caldera}, en el pueblo, y \textbf{Nº 22, Vaqueros}. Firman \textbf{Cornejo} y \textbf{José T.\ Solá Torino}.

\textbf{Tres escuelas, y la de Calderilla es nacional.} La distinción importa porque el
capítulo \ref{cap:presencia} mide qué pone la Provincia y qué pone la Nación: en 1947 el
departamento tiene \textbf{dos escuelas provinciales} ---el pueblo y Vaqueros--- y
\textbf{una nacional} en Calderilla. \textbf{Y el subsidio no va a las escuelas sino a sus
cooperadoras}, que es otra cosa: cien pesos a la asociación de padres, por una vez, y a través
de un banco. \textbf{Ese mismo año la Provincia presupuesta \$690.000 para un puente en el
mismo departamento.}

\textbf{Y una trampa que conviene dejar anotada}: la misma lista registra la escuela nacional
\textbf{Nº 212, «Entre Ríos»}, y no está en este departamento sino en \textbf{La Viña}. El
topónimo que este libro sigue desde 1909 ---la finca «Vaqueros o Entre Ríos»--- tiene
homónimo provincial, y cualquier búsqueda por nombre lo va a traer.

\textbf{Lo que el departamento recibe, con su denominador.} Un decreto de septiembre de 1947
reparte la \textbf{participación de la Ley nacional 12.956} entre \textbf{cuarenta y dos
municipalidades} de la provincia, por un total de \textbf{\$209.775,06}, con cargo a la cuenta
especial «\textit{Municipalidades de la Provincia}». \textbf{La Caldera figura en el orden 30
con \$266,41} (B.O.\ Nº 2932, h.\ 5, verificado sobre el facsímil), y otras ediciones del
bienio publican su coeficiente: \textbf{0,127 por ciento}, el mismo en 1947 y en 1948.

\textbf{Y el decreto que fija los porcentajes dice de dónde salen.} El \textbf{decreto 5276-E
del 30 de julio de 1947} (expte.\ 1873/C/1947), verificado sobre el facsímil, encadena tres
normas: el artículo 4º de la \textbf{Ley nacional 12.956} obliga a las provincias a distribuir
trimestralmente entre sus municipalidades «\textit{no menos del \textbf{10 \%} de la
participación que reciban de la Nación}»; el artículo 50 de la \textbf{Ley provincial 833} da
a los municipios «\textit{una participación del producto del Impuesto a la Contribución
territorial igual al \textbf{20 \%}}»; y el artículo 6º de la \textbf{Ley de Presupuesto 834} lo dispone.
\textbf{Ese artículo se leyó sobre la separata del Boletín, y no sólo en la cita del
decreto}: «\textit{Los proporcionales que sobre la recaudación del impuesto a los Réditos,
Ventas, y Beneficios Extraordinarios y Contribución Territorial le corresponde a las
Municipalidades de Salta y de la Campaña, serán liquidados \textbf{en proporción a los
recursos efectivamente percibidos por cada una de ellas durante el año 1946}. En lo que
respecta a la participación que le corresponde en el impuesto a los \textbf{Espectáculos
Públicos} éste será distribuido en proporción a la recaudación que del mismo efectúe cada
Municipalidad}».

\textbf{Y publica la tabla.} La recaudación propia de cada municipio en 1946, y el porcentaje
que de ella resulta: Salta capital, \textbf{\$1.298.124} y el \textbf{54,116 \%}; Orán,
\$201.982 y el 8,446; \ldots{} y \textbf{La Caldera, en el orden 30, \$3.068 y el
\textbf{0,127 \%}} (lámina \ref{fig:coparticipacion}). Por debajo sólo quedan Cachi, con \$2.406, y Urundel, con \$2.312.

\begin{cautela}
\textbf{El 0,127 por ciento no es una asignación: es un espejo.} La ley no reparte según
población ni según superficie ni según necesidad, sino \textbf{en proporción a lo que cada
municipio ya recaudaba por su cuenta}. \textbf{El que recauda poco recibe poco de lo que la
Nación manda}, y lo recibe poco \textit{porque} recauda poco.

\textbf{Este libro venía suponiendo que el coeficiente podía responder a población, a
superficie o a recaudación, y que el archivo no lo decía. Lo dice, y es la tercera.} La
regla está escrita en el artículo 6º de una ley de presupuesto, publicada, y con su tabla al
lado: cualquiera puede rehacer la cuenta. \textbf{Es un criterio explícito, y es regresivo por
construcción}: la coparticipación no corrige la asimetría entre municipios, la reproduce.

\textbf{Y la ley trae una segunda oración que el decreto no cita, y que aprieta más la
regla.} El impuesto a los \textbf{espectáculos públicos} se reparte «\textit{en proporción a
la recaudación que del mismo efectúe cada Municipalidad}». \textbf{Un departamento sin
cine, sin teatro y sin cancha con entrada recauda cero por ese impuesto, y por lo tanto recibe
cero.} No es que reciba poco: queda fuera del reparto. \textbf{Y eso vale para las dos
oraciones del artículo}, porque las dos reparten según lo que cada municipio ya cobra. La
diferencia es que en la primera el departamento entra con el 0,127 por ciento y en la segunda
no entra.

\textbf{Y fija el suelo del departamento en una cifra.} En 1946 el municipio de La Caldera
recaudó \textbf{\$3.068}, que es el \textbf{0,24 por ciento} de lo que recaudó la capital.
El capítulo \ref{cap:presencia} venía midiendo la presencia del Estado en el departamento por
lo que el Estado ponía; \textbf{ésta es la primera cifra de lo que el departamento ponía de sí
mismo}.

\textbf{Y hay un dato de ausencia en el mismo acto.} Su artículo 3º manda retener la parte
proporcional de los convenios que cuatro municipalidades ---\textbf{Salta, Metán, Rosario de
la Frontera y Campo Santo}--- celebraron con la Dirección Provincial de Sanidad. \textbf{La
Caldera no tiene convenio de sanidad}, y por eso no se le retiene nada: cobra menos y también
debe menos, que es la forma en que un departamento sin servicios aparece en una planilla.
\end{cautela}

\textbf{Y el plan hidráulico que no se publica.} El \textbf{decreto 5466-E del 14 de agosto
de 1947} (expte.\ 17847/1947) aprueba «\textit{provisoriamente el \textbf{plan hidráulico
orgánico de financiación y construcción de obras en el período 1948--1950} preparado por la
Administración General de Aguas de Salta y que representa en dicho período una inversión de
\textbf{\$35.100.000}}», por la resolución Nº 326 del acta Nº 16 del H.\ Consejo, del 18 de
julio de 1947, «\textit{y que corre de \textbf{fs.\ 12 a 18} del expediente precitado}». Se
dicta por el artículo 92 inciso i) del \textbf{Código de Aguas entonces vigente, la Ley
Nº 775}, y el considerando aclara que el plan «\textit{no se superpone con el plan de obras
previsto en el \textbf{Plan Quinquenal} del Superior Gobierno de la Nación}».

\begin{cautela}
\textbf{El acto publica el monto y no publica el plan.} Treinta y cinco millones de pesos para
obras hidráulicas en tres años, y qué obras, dónde y en qué orden queda en las fojas 12 a 18
de un expediente que el Boletín no reproduce. \textbf{Es, en 1947, la misma forma que el
capítulo \ref{cap:opacidad} documenta para los instrumentos de 2013 en adelante}: se publica
el marco y se reserva el contenido operativo.

\textbf{Y para este departamento la pregunta es precisa.} El capítulo viene siguiendo desde
1938 los intentos de defensa del río de la Caldera ---croquis en 1938, obra suspendida en
1940, licitación y ejecución por vía administrativa en 1945 por \$2.515--- y no sabe qué pasó
después. \textbf{Si esa defensa está en las fojas 12 a 18 de este plan, el hilo sigue; si no
está, se corta acá y el libro puede decirlo.} El apéndice \ref{cap:pedidos} pide esas siete
fojas.
\end{cautela}

\textbf{Y el presupuesto de 1947 sí pone plata en el río.} La \textbf{Ley 834}, que fija el
presupuesto provincial del año en \textbf{\$13.560.293,76}, trae en su \textbf{Ítem 6}
---«\textit{Obras de defensas permanentes o eventuales para contener daños por crecientes o
inundaciones sobre}»--- una lista de nueve ríos por un total de \textbf{\$560.000}. \textbf{Tres
de los nueve son de este departamento}: el \textbf{Mojotoro}, \textbf{\$100.000}; \textbf{La
Caldera}, \textbf{\$50.000}; y el \textbf{Wierna}, \textbf{\$50.000}, los tres en primera
etapa. Los otros seis son el Rosario (Toro), Chuñapampa, el Cachi, el Calchaquí, el Juramento y
el de Iruya, más \$60.000 de imprevistos (lámina \ref{fig:defensas}) (separata del B.O.\ del 19 de mayo de 1947,
págs.\ 141--142).

\begin{cautela}
\textbf{Doscientos mil pesos de quinientos sesenta mil: el 35,7 por ciento de todo lo que la
Provincia presupuesta en 1947 contra las crecientes va a los tres ríos de este departamento.}
Y los tres son uno solo: el acto de 1946 que este capítulo transcribe establece que el Mojotoro
forma su caudal con el Wierna, el de La Caldera y el Castellanos.

\textbf{Esto corrige lo que este libro venía dando por sentado.} El capítulo sigue los intentos
de defensa del río desde el croquis de 1938 y la obra suspendida de 1940 hasta la ejecución por
vía administrativa de 1945, \textbf{por \$2.515}, y la lectura implícita era que la Provincia no
ponía plata. \textbf{La ponía}: en 1947 hay una partida en la ley de presupuesto, votada por las
dos cámaras, de \textbf{veinte veces} lo gastado dos años antes, y el departamento concentra más
de un tercio del ítem provincial.

\textbf{Lo que no se sabe es si se ejecutó}, y ahí vuelve el patrón de este bienio: el mercado
presupuestado en julio y desafectado en diciembre, el camino objetado, el puente de \$690.000
sin licitación publicada. \textbf{Presupuestado no es ejecutado.} Y la pieza que podría
contestarlo son las fojas 12 a 18 del plan hidráulico 1948--1950, que el decreto 5466-E aprueba
sin publicar: \textbf{la partida de 1947 y el plan de 1948 son dos tramos del mismo expediente,
y de los dos se conoce el monto y no el destino.}
\end{cautela}

\textbf{Y la estación sanitaria vuelve, con otro origen.} El \textbf{Inciso VI} del mismo
presupuesto ---«\textit{Plan de obras con fondos de la \textbf{Ley Nº 770}}»--- destina
\textbf{\$370.000} a «\textit{estudios, proyectos y construcción de nuevos edificios para
\textbf{estaciones sanitarias}, sus instalaciones y amueblamiento, en las siguientes
localidades: El Tala; \textbf{La Caldera}; Pichanal; Joaquín V.\ González, Iruya, Santa
Victoria}». \textbf{Seis localidades en una sola partida}, sin decir cuánto le toca a cada una.
Es la \textbf{misma Ley 770} que financia el plan vial donde está el puente sobre el río, de
modo que \textbf{el puente y la estación sanitaria salen del mismo empréstito}, y el capítulo
\ref{cap:presencia} puede ahora nombrar la norma que paga lo que mide.

\textbf{El mercado que se presupuestó y se desafectó en el mismo año.} El decreto
\textbf{Nº 4945 del 1º de julio de 1947} destinó \textbf{\$600.000} a la construcción de
mercados en municipios, y entre ellos un \textbf{Mercado Tipo I para La Caldera}. En
diciembre, el decreto que adjudica los mercados de los demás departamentos
\textbf{«desaféctase provisoriamente»} el de La Caldera, junto con los de La Candelaria, La
Poma y San Carlos y el Tipo II de San Antonio de los Cobres, «\textit{para que con dicha
economía \ldots{} se haga viable la construcción de un \textbf{Mercado Frigorífico} en
\textbf{Embarcación}}» (B.O.\ Nº 3016, h.\ 9, verificado sobre el facsímil).

\begin{cautela}
\textbf{El motivo que el acto da es el que conviene leer despacio}: «\textit{a mérito de
considerarse ineficiosa la construcción de dichas obras en los pueblos antedichos,
\textbf{expresión que se había manifestado por los propios habitantes}}».

\textbf{El archivo no publica quién lo manifestó, cuándo ni por qué vía.} No hay nota, ni
petición, ni acta de la comisión municipal, ni expediente citado: la voluntad de los
habitantes aparece en el considerando de un decreto que les quita una obra, invocada por el
mismo organismo que la quita. \textbf{Es la única vez en el archivo temprano en que un acto
provincial dice hablar por los habitantes del departamento}, y lo hace para justificar una
desafectación. El apéndice \ref{cap:pedidos} pide el expediente.

\textbf{Y hay un antecedente exacto, tres años antes.} En mayo de 1944 el decreto 3016-H
aprueba la \textbf{eliminación del camino de Vaqueros a Los Yacones} del plan de obras de la
Ley nacional 12.776 (apéndice \ref{cap:cronologia}). \textbf{Dos obras presupuestadas para el
departamento y sacadas del plan en tres años}, mientras en 1947 los mercados sí se construyen
en Cerrillos, Rosario de Lerma, Campo Quijano, La Merced, San Lorenzo y Chicoana. \textbf{Lo
que el departamento recibe no es menos presupuesto: es más desafectación.}
\end{cautela}

\textbf{Y el contratista reaparece.} El mismo acto nombra a \textbf{Conrado Marcuzzi} entre
los oferentes que dieron cumplimiento al pliego. Es el mismo a quien, once meses después, se
le adjudica la obra sanitaria de la escuela del pueblo por \$35.026,99 ---la que necesitó una
insistencia en Acuerdo de Ministros para pagarse---.

\textbf{El camino, y la única objeción del Ejecutivo.} El \textbf{decreto 5916-E del 22 de
septiembre de 1947} aprueba el Acta Nº 188 del H.\ Consejo de Administración de Vialidad
\textbf{«a excepción de la resolución Nº 4967»}, y ordena que la Administración
«\textit{deberá informar en mérito a qué consideraciones ha sido incluido, \textbf{fuera del
plan}, el camino denominado «La Caldera a Santa Rufina», en la red de caminos
provinciales}» (B.O.\ Nº 2949, h.\ 12, verificado sobre el facsímil).

\textbf{El puente.} El \textbf{plan de obras de la Administración de Vialidad} de 1947
consigna «\textit{Puente S/Río La Caldera. Obras de arte Luz 150 mts.}» por
\textbf{\$690.000}, sobre un plan provincial total de \textbf{\$12.867.275,27}
(B.O.\ Nº 2983, h.\ 7, verificado sobre el facsímil).

\textbf{El Registro Civil, y quién fabrica sus muebles.} El \textbf{decreto 5537-G del 20 de
agosto de 1947} (expte.\ 6770/947) adjudica «\textit{a la \textbf{CÁRCEL PENITENCIARIA} la
confección de una mesa escritorio y tres sillas \ldots{} con destino a la Oficina de Registro
Civil de "La Caldera", al precio total de \textbf{ciento treinta y dos pesos con cuarenta}}»
(B.O.\ Nº 2923, h.\ 6, verificado sobre el facsímil).

\textbf{La salud, y un traslado que el libro no tenía.} El \textbf{decreto 9089-A del 8 de
abril de 1948}, verificado sobre el facsímil, dice: «\textit{Encontrándose vacante el cargo de
Enfermero del pueblo de \textbf{Rivadavia} \ldots{} nómbrase Auxiliar 7º (Enfermero de
Rivadavia) \ldots{} al actual Ayudante 5º (Enfermero de La Caldera), don \textbf{CLAUDIO
MURÚA}}»; y en su artículo 2º, «\textit{nómbrase Ayudante 5º (Enfermero de La Caldera), al
señor \textbf{TEODORO GÓMEZ}, en la vacante dejada por don Claudio Murúa}».

\textbf{La mina.} La \textbf{Resolución de Minas Nº 591}, del 2 de abril de 1948
---verificada sobre el facsímil, hoja 5---, adjudica «\textit{a favor del señor
\textbf{Enrique A.\ Zuccarino} y Sra.\ \textbf{María Josefina Méndez Alcorta de Zuccarino}
\ldots{} la presente mina de \textbf{plomo y plata} denominada "CALDERA No.\ 1, 2 y 3",
ubicada en el departamento de La Caldera, Exp.\ No.\ \textbf{44-L año 1929}}», por la última
parte del artículo 7º de la \textbf{Ley Nacional 10.273}. Firma \textbf{Luis Víctor Outes};
ante el escribano de minas \textbf{Oscar M.\ Aráoz Alemán}.
\end{ficha}

\begin{cautela}
\textbf{De toda el acta, lo único que el gobernador no aprueba es el camino del
departamento.} No lo rechaza: lo suspende y pide explicaciones, que es una figura distinta y
más rara. \textbf{Un camino entró a la red provincial por fuera del plan de obras}, y el
Poder Ejecutivo quiere saber con qué criterio. Es la primera vez en el archivo temprano que
un acto provincial \textbf{pregunta por el criterio} de una decisión sobre el departamento en
lugar de tomarla o publicarla, y este libro no tiene la respuesta: el apéndice
\ref{cap:pedidos} pide el informe que Vialidad debía elevar. \textbf{Santa Rufina}, además,
no es un destino cualquiera: la nómina electoral de 1948 la registra como paraje del circuito
del pueblo, y el apellido asociado a esa finca aparece en los pedimentos mineros de 1925.

\textbf{La escala del bienio, en una línea.} \$132,40 el mobiliario del Registro Civil,
\$35.026,99 los sanitarios de la escuela, \textbf{\$690.000 el puente}. \textbf{El puente
solo vale veinte escuelas y cinco mil doscientos Registros Civiles}, y es \textbf{el 5,4 por
ciento de todo el plan vial de la provincia}. El capítulo \ref{cap:presencia} mide el mismo
reparto para 1923--1945 y encuentra que lo único caro eran los caminos: 1947 y 1948 no lo
desmienten, lo confirman con tres cifras del mismo bienio, que \textbf{sí son comparables
entre sí} porque no median años de inflación ---las de 1934 no lo son, y este libro no las
corrige---.

\textbf{Y la escuela deja ver algo más.} La obra se adjudica el 19 de noviembre y ocho días
después el gobernador tiene que \textbf{insistir en Acuerdo de Ministros} contra las
observaciones de Contaduría General. \textbf{La partida más grande que el archivo temprano
destina a un edificio del departamento que no sea un camino necesitó saltar el control
contable para pagarse}, y el acto no dice cuál era la observación. El apéndice
\ref{cap:pedidos} la pide.

\textbf{Y la mina cierra una serie de diecinueve años, con tres datos que el libro no
tenía.} El expediente \textbf{44-L} es el que el padrón minero provincial cerrado el 30 de
marzo de 1933 registra bajo el número de orden 93, inscripto el 30 de abril de 1930 a nombre
de \textbf{Juan Larrán}: veintisiete hectáreas, tres pertenencias (capítulo
\ref{cap:siglo}).

\textbf{Primero: la mina es de plomo y plata}, y es la única sustancia metalífera que el
archivo temprano asocia a este departamento ---todo lo demás, antes y después, son áridos---.
\textbf{Segundo: las «tres pertenencias» del padrón de 1933 son las tres minas} que el acto
de 1948 nombra, «Caldera No.\ 1, 2 y 3». \textbf{Y tercero, el que importa: no se adjudica a
Larrán.} Entre la inscripción de 1930 y la adjudicación de 1948 el expediente cambió de
manos, y el acto no dice cuándo ni cómo: conserva el número de 1929 y nombra a otros dos.
\textbf{Diecinueve años entre inscribir y adjudicar}, y en el medio un traspaso que el
Boletín no publicó. El apéndice \ref{cap:pedidos} pide el expediente completo, que es lo
único que puede contar ese medio.

\begin{cautela}
\textbf{Murúa no deja el cargo: asciende y se va.} Pasa de \textbf{Ayudante 5º} en La Caldera
a \textbf{Auxiliar 7º} en Rivadavia, que es una categoría más alta y otro departamento.
\textbf{Es el primer caso documentado en el archivo temprano de un empleado del departamento
que sale de él por promoción}: todos los movimientos anteriores que este capítulo registra son
relevos dentro del departamento, ceses o intervenciones. \textbf{Y el que llega no viene de
ningún lado conocido}: Teodoro Gómez entra directamente en la vacante.
\end{cautela}

\textbf{Y \textit{Murúa} otra vez.} El apellido que este capítulo sigue desde 1912 ---jurado
de patentes, comisario en 1916, 1923, 1930 y 1932, causante del sucesorio que remata La
Calderilla en 1923, titular de La Despensa en 1925, miembro de la Comisión Municipal en 1938,
comisario en 1939 y nombrado en el expediente de pensión de 1941--- aparece en 1948 dejando
vacante el cargo de enfermero del departamento. \textbf{Coincidencia de nombre, como siempre.}

\textbf{Y en 1947 y 1948 el apellido aparece dos veces más, y las dos en los tribunales.} El
sumario de la edición 3019 de 1947 publica una \textbf{notificación de sentencia}
«\textit{en ejecución seguido por Banco Provincial de Salta vs.\ \textbf{Silvano U.\ Murúa} y
otros}» (aviso Nº 3347, h.\ 13 y 14); y el de la edición 3064 de 1948 abre el
\textbf{juicio sucesorio} «\textit{de don \textbf{Silvano, I.\ O.\ Ignacio Murúa}}» (aviso
Nº 3442, h.\ 12). \textbf{Los dos avisos se verificaron sobre el facsímil} ---son de las
hojas de sumario que el repositorio sí recortó---. Un ejecutado por el banco provincial en
diciembre de 1947 y un sucesorio abierto al año siguiente: \textbf{el archivo no dice que
sean la misma persona, y este libro tampoco}, pero el apellido que administró el
departamento desde 1912 termina el bienio en ejecución y en sucesión.
\end{cautela}

\subsection*{Una trampa de homonimia dentro del propio archivo}

\textbf{Las dos puntas están verificadas sobre el facsímil}, y dicen esto.

La \textbf{Resolución Nº 423-E} del 7 de noviembre de 1947 (expte.\ 18646/47; B.O.\ Nº 2985,
h.\ 7) autoriza una nota de crédito de \textbf{\$541,20} «\textit{a favor del Receptor de
Rentas de la localidad de La Caldera, \textbf{Departamento de la Capital}, don GERARDO
GUERRERO}», por estampillas «\textit{aplicadas \textbf{por error} en las boletas de
Contribución Territorial}». La adscripción a Capital \textbf{aparece dos veces en el mismo
acto}, en el visto y en la parte resolutiva, de modo que no es un desliz de una línea.

El \textbf{Decreto Nº 8706-E} del 10 de marzo de 1948 (expte.\ 4567/F/1948; B.O.\ Nº 3083,
h.\ 5) «\textit{desígnase Receptor de Rentas y Expendedor de Guías, Transferencia de Cueros
Marcas y Señales, y Multas Policiales del \textbf{Departamento de La Caldera}, al señor
GERARDO GUERRERO}». Y lo hace «\textit{atento a lo solicitado por el \textbf{Banco Provincial
de Salta}}».

\begin{cautela}
\textbf{El acto de 1948 no confirma al de 1947: lo designa de nuevo.} Cuatro meses después de
cobrar como receptor de una localidad del departamento Capital, el mismo hombre es
\textbf{designado} receptor del departamento de La Caldera. Caben las dos lecturas ---que el
Boletín se equivocó dos veces en 1947, o que Guerrero cambió de destino--- y este libro no
elige: registra las dos y pide el legajo.

\textbf{Y sobre quién pide la designación, este libro se corrige a sí mismo.} Que la pida el
\textbf{Banco Provincial de Salta} parecía una anomalía ---no la piden el municipio, ni la
Dirección General de Rentas, ni la policía, que son los tres organismos cuyos ingresos el
cargo recauda---. \textbf{No lo es, y el propio archivo lo explica cuatro meses antes.} El
\textbf{decreto 5919-E del 22 de septiembre de 1947} (expte.\ 2376/B/47) designa a los
expendedores de guías de \textbf{Orán} y de \textbf{Campo Santo} con la misma fórmula
---«\textit{atento lo solicitado por el Banco Provincial de Salta \ldots{} y lo dispuesto por
el art.\ 3º de la \textbf{Ley Nº 744}}»--- y dice para qué: cada designado «\textit{deberá
prestar \textbf{fianza a satisfacción del Banco Provincial de Salta}}».

\textbf{El banco no elige al recaudador: es ante quien el recaudador afianza.} La solicitud
es el trámite que el artículo 3º de la Ley 744 le asigna, y no una preferencia. Lo que en un
solo acto parecía una decisión bancaria sobre un departamento, visto en serie es un
procedimiento provincial que se aplica igual en tres departamentos distintos.
\textbf{Registrarlo así es el punto}: un acto aislado se lee mal, y la serie corrige la
lectura.
\end{cautela}

\subsection*{La decimotercera nómina, y lo que le hace a la lámina}

El \textbf{Decreto 8176 del 5 de febrero de 1948} fija ubicación y circuitos de los colegios
electorales 1 a 6 ---«\textit{La Capital, La Caldera, Campo Santo, Metán, Anta y
Rivadavia}»---, y el del departamento se publica entero, repetido en seis ediciones
seguidas (B.O.\ Nº 3060 a 3065; el sumario lo lista bajo el aviso Nº 3459, hoja 3).
\textbf{Verificado sobre el facsímil} de la hoja 19 de la edición 3062.

\textbf{La Caldera es el Colegio Electoral Nº 2}, y tiene dos circuitos.

\textbf{El circuito Nº 5} es el pueblo y el interior, con tres mesas: dos en la escuela
elemental mixta y una en la iglesia parroquial. Las mesas 1 y 3 votan
\textbf{Calderilla, Chalchanio, Helvecia, Getsemaní, Peñones, Potrero de Castilla, Pueblo,
San Alejo y Severino}. La mesa 2, \textbf{Abra de la Sierra, Acheral, El Angosto, Angosto
de Arrieta, Angostura, El Bordo, Calderilla, Campo Alegre, Cañada Ancha, Corral de
Barranca, Churcal, Los Chusos, Despensa, Despensilla, Higueritas, Potrero de Gallinato,
Las Lagunas, Los Matos, Mesada, Monte de las Garrapatas, Mosquera, Nogalito, Potrero de
Valencia, Los Porongos, Quesería, Santa Gertrudis, Santa Rufina, Los Sauces, La Sierra,
Wierna y Yacones}.

\textbf{El circuito Nº 6 es Vaqueros}, con dos mesas ---la Escuela Infantil Provincial y el
Hogar Agrícola «San Cayetano»--- y siete nombres: \textbf{Entre Ríos, Estación
F.\ C.\ C.\ N.\ A., Mojotoro, Lesser, El Túnel, Vaqueros y Villa Urquiza}.

\begin{cautela}
\textbf{Es la decimotercera nómina del mismo territorio, y la primera posterior a 1940.} La
lámina \ref{fig:parajesnom} reúne \textbf{doce} nóminas oficiales entre 1909 y 1940 y da
cincuenta y dos nombres distintos. \textbf{Esta no está en la lámina}, y hasta que se
incorpore ---lo que exige rehacer la tabla con el script que la produce, sobre el facsímil y
no sobre la capa de texto--- las dos afirmaciones que la lámina sostiene quedan acotadas a
1909--1940 y no al archivo temprano entero.

\textbf{Lo que ya se puede decir es que la lista de 1948 trae nombres que las doce no
traían}: \textbf{Helvecia} ---que es la Calderilla rebautizada, y aparece como paraje el mismo
año en que pide agua---, \textbf{Getsemaní} ---que el deslinde de 1916 nombraba como propiedad
de Daniel Linares, el mismo que en 1948 pide el agua de El Angosto, y no como paraje---, \textbf{Severino}, \textbf{Abra de la Sierra},
\textbf{Cañada Ancha}, \textbf{Corral de Barranca}, \textbf{Churcal}, \textbf{Los Chusos} y
\textbf{El Bordo}. \textbf{Y trae los que no se mueven}: Calderilla, Campo Alegre, San Alejo,
Potrero de Castilla, Peñones, Wierna y Yacones, que estaban en las doce.

\textbf{Y el Tribunal Electoral publica quiénes presiden las mesas.} Para los dos circuitos
del departamento designa \textbf{quince personas} ---presidente y dos suplentes por cada una de
las cinco mesas--- (B.O.\ Nº 3069, h.\ 14, verificado sobre el facsímil). \textbf{Este libro no
las transcribe}: son designaciones de 1948 y, por la regla de los cien años que el capítulo
\ref{cap:metodo} fija, esas personas se presumen vivas. Se nombra una sola, porque el
argumento la necesita.

\begin{cautela}
\textbf{En la mesa 2 del circuito de Vaqueros, el suplente segundo es un \textit{Urquiza}.}
Veinticuatro años antes, en 1924, \textbf{Francisco S.\ Urquiza} era presidente de mesa
electoral en Vaqueros, y ése era uno de los cinco roles simultáneos que este capítulo le
documenta. \textbf{Y la nómina de parajes del mismo decreto de 1948 registra \textit{Villa
Urquiza} como uno de los siete nombres del circuito.}

\textbf{Coincidencia de apellido, con la regla de siempre}: el libro no afirma parentesco ni
identidad. \textbf{Lo que sí queda registrado es que el apellido está en los tres lugares a la
vez} ---en el topónimo, en la mesa electoral y en la serie de 1924--- y que el lugar es el
mismo. Es el patrón que el capítulo viene siguiendo desde 1909, y en 1948 sigue en pie.
\end{cautela}

\textbf{Y el circuito Nº 6 contesta una pregunta que este libro traía abierta desde 1909}:
dónde cae la finca \textbf{Entre Ríos}. Cae en Vaqueros, y encabeza su nómina. La
«\textit{Vaqueros o Entre Ríos}» de los remates y los deslindes es, en 1948, el primer
paraje de su propio circuito electoral. Y con ella aparecen dos nombres que el archivo
temprano no traía: \textbf{El Túnel} ---la obra del ferrocarril, convertida en topónimo--- y
\textbf{Villa Urquiza}, que es la primera vez que el archivo nombra una \textit{villa} en el
departamento.

\textbf{El circuito electoral parte al departamento en dos
veintidós años antes de que la ley cree el segundo municipio}, y ocho meses después de que el
acto de límites de 1947 declarara que el municipio de La Caldera abarca el departamento
entero. \textbf{Las dos cosas son del mismo Estado y del mismo bienio.}
\end{cautela}


\section{1949: la obra que sí se hizo, y el decreto que se la dio a otro}

El año se leyó edición por edición y es el \textbf{trigésimo quinto tramo}: \textbf{281
archivos} y \textbf{4.869 hojas}, el más voluminoso de toda la serie temprana. Rige la misma
reserva que para 1947 y 1948 ---barrido automático, \textbf{sin lectura sobre la imagen} salvo
las hojas que el apéndice \ref{cap:fuentes} enumera---, y las tapas repiten la patología de
1947: el lector devuelve «Nº 4780» tres veces, «Nº 1805», y fechas imposibles como
«1349-01-07» y «3949-01-20». \textbf{El número de edición de la tapa no es confiable en 1949};
se toma del nombre del archivo.

\begin{ficha}{\textbf{Decreto Nº 13793-E} del 31 de enero de 1949, orden de pago 501, expediente 404/E/1949 (B.O.\ Nº 3337, h.\ 11)}
\textbf{Verificado sobre el facsímil.} «\textit{Apruébase el certificado No.\ 1 \ldots{} por
trabajos ejecutados en la obra: \textbf{Estación Sanitaria en La Caldera}, por la Empresa
\textbf{E.C.O.R.M.\ S.R.L.}, el que asciende a la suma de \textbf{\$27.024,10} \ldots{}
adjudicada al contratista Empresa E.C.O.R.M.\ S.R.L.\ \textbf{por decreto No.\ 10668 del 29 de
julio de 1948}}».

El artículo 3º manda retener el \textbf{10 por ciento} en concepto de garantía de obra, y el
4º imputa el gasto al \textbf{Anexo I, Inciso III, Principal 1/c} del presupuesto de 1948.
Ediciones posteriores del mismo año publican el acta «\textit{para la entrega de la obra}».
\end{ficha}

\begin{cautela}
\textbf{Este libro se corrige.} La sección anterior sostiene que el bienio 1947--1948 es de un
Estado que presupuesta y retira, y enumera cuatro obras de las que tres salieron del plan.
\textbf{La cuarta no salió: se hizo.} La estación sanitaria estaba en el Inciso VI de la Ley
834 de mayo de 1947, dentro de los \$370.000 repartidos entre seis localidades; se
\textbf{adjudicó el 29 de julio de 1948}; su primer certificado se pagó el 31 de enero de 1949;
y se entregó ese año. \textbf{La cadena completa está publicada, acto por acto, y se puede
seguir.}

\textbf{De modo que el patrón existe pero no es una ley, y conviene decir de qué depende.} Lo
que se cayó fue el mercado, el camino y el puente ---obra pública de fomento y de
infraestructura---. Lo que se hizo fue la escuela y la estación sanitaria. \textbf{En este
departamento, en este bienio, lo que el Estado termina es lo que atiende cuerpos.}
\end{cautela}

\subsection*{Y en la misma página está el mercado que no se hizo}

El decreto que aprueba el certificado de la estación sanitaria comparte hoja con otro que
aprueba el certificado parcial Nº 1 de la obra «\textit{\textbf{Mercado Tipo II en
Pichanal}}», por \textbf{\$9.783,82}, «\textit{trabajos adjudicados al señor Arquitecto
\textbf{Luis Moreno Díaz} por \textbf{decreto Nº 7321 del 13 de diciembre de 1947}}».

\begin{cautela}
\textbf{Ése es el mismo decreto que desafectó el mercado de La Caldera.} El acto del 13 de
diciembre de 1947 que quitó el Mercado Tipo I de este departamento ---«\textit{a mérito de
considerarse ineficiosa la construcción de dichas obras en los pueblos antedichos, expresión
que se había manifestado por los propios habitantes}»--- es el que \textbf{adjudicó el de
Pichanal}. Y en febrero de 1949 el de Pichanal se está certificando y pagando.

\textbf{No son dos decisiones: es una.} La economía que justificó la desafectación de cinco
mercados ---La Caldera, La Candelaria, La Poma, San Carlos y San Antonio de los Cobres--- es
la que financió los que sí se construyeron. \textbf{El mismo número de decreto está en las dos
puntas}, y el archivo lo publica en la misma página trece meses después. Lo que el libro no
tiene es la constancia de esa expresión de los habitantes, que el apéndice
\ref{cap:pedidos} reclama.
\end{cautela}

\subsection*{El río La Caldera pasa a ser el primero de la provincia}

El presupuesto de 1949 repite la estructura de la Ley 834, y su \textbf{Ítem 6} ---las obras de
defensa contra crecientes--- reparte \textbf{\$460.000} entre diecisiete ríos y un renglón de
imprevistos, cien mil menos que en 1947 (B.O.\ Nº 3341, h.\ 12, verificado sobre el facsímil).
El reparto del departamento cambia de forma:

\begin{center}
\small
\begin{tabular}{lrr}
\toprule
 & \textbf{1947} & \textbf{1949} \\
\midrule
La Caldera              & 50.000  & \textbf{80.000} \\
Mojotoro                & 100.000 & 50.000 \\
Wierna                  & 50.000  & 10.000 \\
Vaqueros                & ---     & 10.000 \\
\midrule
\textbf{Del departamento} & \textbf{200.000} & \textbf{150.000} \\
Total del ítem          & 560.000 & 460.000 \\
Participación           & 35,7 \% & 32,6 \% \\
\bottomrule
\end{tabular}
\end{center}

\begin{cautela}
\textbf{Con ochenta mil pesos, el río de la Caldera es la partida más alta de las dieciocho del
ítem}, por encima de Iruya, del Colorado y de San Antonio de los Cobres. \textbf{Y Vaqueros
aparece como río propio por primera vez}: en 1947 el presupuesto lo contaba dentro del
Mojotoro, y en 1949 tiene su renglón. \textbf{Es la primera vez que un instrumento provincial
separa a Vaqueros del sistema}, veintiún años antes de que la ley lo separe como municipio.

\textbf{La participación del departamento baja y su río sube}, que no es lo mismo: el ítem
entero se achicó, el Mojotoro y el Wierna perdieron la mitad y cuatro quintos, y lo que quedó
se concentró en el río del pueblo. \textbf{Este libro no sabe si esos ochenta mil se
ejecutaron}, y por las mismas razones de siempre: el archivo publica el presupuesto y no el
certificado.

\textbf{Y hay un segundo ítem con el departamento adentro.} El \textbf{Ítem 5} destina
\textbf{\$120.000} a «\textit{expropiación, estudios, proyectos, conservación y refección de
servicios de extensión y reducido (con captación superficial), en las siguientes localidades:
El Tala, Tolombón, Molinos, \textbf{La Caldera}, La Poma, Iruya y otras localidades de la
Provincia (obras de fomento)}».
\end{cautela}

\subsection*{La mina se adjudica por segunda vez}

\begin{ficha}{\textbf{Resolución de Minas Nº 672}, del 10 de febrero de 1949 (B.O.\ Nº 3346, h.\ 9)}
\textbf{Verificada sobre el facsímil.} «\textit{Y VISTOS: El escrito que antecede y de acuerdo
con lo dispuesto por la Ley Nacional 10.273, art.\ 7º, \textbf{in fine}, adjudícase a favor del
señor \textbf{Enrique A.\ Zuccarino} y señora \textbf{María Josefina Méndez Alcorta de
Zuccarino}, la mina de plomo y plata denominada "Caldera 1, 2 y 3", ubicada en el Departamento
del mismo nombre, \textbf{lugar denominado San José}, Exp.\ No.\ 44---L, año 1929}».
\end{ficha}

\begin{cautela}
\textbf{Es la misma adjudicación de diez meses antes, palabra por palabra}, incluida la
invocación del artículo 7º \textit{in fine} de la Ley 10.273: la Resolución Nº 591 del 2 de
abril de 1948 ya la había dictado, con los mismos adjudicatarios y el mismo expediente.
\textbf{Dos adjudicaciones idénticas del mismo derecho, con diez meses de diferencia}, y el
Boletín no publica nada en el medio. O la primera no se registró, o no quedó firme, o hubo un
trámite que no se publicó; el expediente 44-L es lo único que puede decirlo, y el apéndice
\ref{cap:pedidos} lo pide.

\textbf{Y la resolución de 1949 trae un dato que la de 1948 no tenía: el lugar.} La mina está
en «\textbf{San José}», que \textbf{no figura entre los parajes de la nómina electoral de
1948} ---donde sí están San Alejo, Santa Rufina y Santa Gertrudis---. Es el único dato de
ubicación que el archivo da de esta mina en veinte años.
\end{cautela}


\section{1950: se cierra la obra y se fracciona Entre Ríos}

El año se leyó edición por edición y es el \textbf{trigésimo sexto tramo}: \textbf{277
archivos} y \textbf{4.156 hojas}, con la misma reserva que los tres anteriores.

\subsection*{La estación sanitaria se recibe}

En la edición \textbf{3684} se aprueba otro certificado por trabajos ejecutados, y en la
\textbf{3864} el trámite se cierra: se devuelven «\textit{los certificados correspondientes a
la obra "estación sanitaria tipo "A" en La Caldera", \textbf{con imputación a la cuenta
especial "Depósitos en garantía"}}», a favor de la empresa constructora E.C.O.R.M.\ S.R.L.

\begin{cautela}
\textbf{Ésa es la devolución de la retención del diez por ciento, y por lo tanto la recepción
definitiva.} La cadena queda completa y publicada, acto por acto, a lo largo de cuatro años:
presupuestada en el Inciso VI de la Ley 834 en \textbf{mayo de 1947}; adjudicada por decreto
10668 el \textbf{29 de julio de 1948}; primer certificado pagado el \textbf{31 de enero de
1949}; entregada ese año; garantía devuelta en \textbf{1950}.

\textbf{Es la única obra del departamento que el archivo temprano permite seguir de punta a
punta}, y conviene medir lo que eso significa. El capítulo \ref{cap:presencia} venía
sosteniendo que de los edificios públicos del departamento sólo se puede seguir la historia
completa de uno, y era aquel que el Estado ocupaba sin título desde hacía treinta años.
\textbf{Ahora son dos, y el segundo es el que se hizo bien.}
\end{cautela}

\subsection*{Y aparece la palabra «fracción»}

Tres ediciones seguidas ---\textbf{3692, 3693 y 3694}--- publican el mismo edicto de la
Administración General de Aguas: alguien solicita reconocimiento de concesión sobre la
propiedad «\textit{denominada \textbf{"Fracción Entre Ríos"}, ubicada en \textbf{Vaqueros},
departamento de La Caldera}», y «\textit{el reconocimiento a otorgarse sería para una dotación
de agua proveniente}» del sistema. El aviso invoca además el \textbf{decreto Nº 3649 del 11 de
julio de 1944}.

\begin{cautela}
\textbf{«Fracción Entre Ríos» no es «Entre Ríos».} La finca que este libro sigue desde los
remates de 1909, que da nombre al municipio de Vaqueros y que en 1948 encabeza la nómina de su
circuito electoral, \textbf{en 1950 ya está dividida}, y una de sus partes pide agua con nombre
propio. \textbf{Es la primera vez que el archivo temprano nombra una fracción de una finca de
este departamento}: hasta acá las fincas aparecían enteras, deslindadas y mensuradas, pero
nunca partidas.

\textbf{Y el agua llega antes que el plano.} El capítulo \ref{cap:loteo} documenta para los
loteos recientes que el orden habitual es el inverso ---primero el fraccionamiento aprobado,
después los servicios---, y acá el reconocimiento de agua sobre una fracción se publica sin
que el Boletín haya publicado antes su plano. \textbf{Este libro no afirma que no exista}: el
plano puede estar en Dirección General de Inmuebles sin pasar por el Boletín, que es
exactamente lo que el capítulo \ref{cap:opacidad} describe. El apéndice \ref{cap:pedidos}
pide el expediente y el plano de esa fracción.
\end{cautela}

\section{Qué hace el bienio 1947--1948, y por qué no es más de lo mismo}

Los dos años que este capítulo acaba de incorporar no aportan un hallazgo suelto: aportan
\textbf{una forma de operar que los treinta y dos tramos anteriores no mostraban junta}. Vale
la pena enunciarla, con la reserva de que \textbf{los dos años entraron al libro por lectura
de capa de texto} y que, de las hojas que sostienen lo que sigue, este libro verificó sobre el
facsímil las que el apéndice \ref{cap:fuentes} enumera y no las demás.

\begin{cautela}
\textbf{No es un bienio de ausencia de Estado. Es uno de Estado que presupuesta y retira.}

\textbf{El mercado.} El decreto 4945 del 1º de julio de 1947 destina \$600.000 a mercados
municipales y prevé uno \textbf{Tipo I para La Caldera}. En diciembre, el mismo Estado lo
\textbf{desafecta} y pasa el dinero al Mercado Frigorífico de Embarcación, «\textit{a mérito
de considerarse ineficiosa la construcción \ldots{} expresión que se había manifestado por los
propios habitantes}». \textbf{Los mercados sí se construyen} en Cerrillos, Rosario de Lerma,
Campo Quijano, La Merced, San Lorenzo y Chicoana.

\textbf{El camino.} El decreto 5916-E aprueba el Acta 188 del Consejo de Vialidad
\textbf{«a excepción de la resolución Nº 4967»}, la única que exceptúa, y exige que se informe
por qué el camino «\textit{La Caldera a Santa Rufina}» entró \textbf{fuera del plan} a la red
provincial.

\textbf{El puente.} El plan vial de 1947 lo presupuesta en \textbf{\$690.000} ---el 5,4 por
ciento de todo el plan provincial--- y el archivo no publica después ni licitación ni
adjudicación ni certificado.

\textbf{La escuela.} La única obra del bienio que llega a ejecutarse, \$35.026,99 en
sanitarios, \textbf{necesita una insistencia en Acuerdo de Ministros} contra las observaciones
de Contaduría General para poder pagarse.

\textbf{Y hay antecedente.} En mayo de 1944 el decreto 3016-H aprueba \textbf{eliminar el
camino de Vaqueros a Los Yacones} del plan de la Ley nacional 12.776. \textbf{Cuatro obras
presupuestadas para este departamento en cinco años, y tres salen del plan.}

\textbf{La cuarta es la excepción y hay que decirla acá, no sólo después.} La estación
sanitaria del Inciso VI se adjudicó en julio de 1948, se certificó en enero de 1949 y se
entregó ese año, con toda su cadena publicada (sección siguiente). \textbf{De modo que la
frase correcta no es que el Estado presupuesta y retira, sino que retira lo que construye y
termina lo que atiende}: se cayeron el mercado, el camino y el puente; se hicieron la escuela
y la estación sanitaria.
\end{cautela}

\begin{cautela}
\textbf{Y el criterio, cuando está escrito, es regresivo.}

El artículo 6º de la Ley de Presupuesto 834 manda liquidar la coparticipación municipal
«\textit{en proporción a los recursos efectivamente percibidos por cada una de ellas durante
el año 1946}», y el decreto 5276-E publica la tabla entera: \textbf{lo que cada municipio
recaudó y el porcentaje que le toca}. La Caldera recaudó \textbf{\$3.068} y le corresponde el
\textbf{0,127 por ciento}; la capital recaudó \$1.298.124 y le corresponde el 54,116.

\textbf{Este libro sostiene que la opacidad de este dispositivo opera en el nivel del criterio
y no del acto, y éste es el contraejemplo que hay que poner.} Acá el criterio está escrito, es
público y su tabla está al lado: cualquiera puede rehacer la cuenta. \textbf{Que se pueda leer
no lo vuelve menos duro. Lo vuelve discutible}, que es exactamente lo que el capítulo
\ref{cap:opacidad} reclama para los instrumentos de 2013 en adelante y no obtiene.

\textbf{Y hay un segundo criterio, éste no enunciado pero deducible.} Los dos reconocimientos
de agua del bienio publican la superficie y el caudal: 24 hectáreas y 12,6 litros por segundo
para La Helvecia, 29 y 15,22 para El Angosto. \textbf{Las dos razones dan 0,525 litros por
segundo por hectárea.} Nadie escribió «0,525» en ningún lado, pero \textbf{publicar los dos
términos de la división vuelve verificable la regla sin necesidad de enunciarla}.

\textbf{De modo que el archivo de 1947--1948 ofrece las tres posiciones posibles}: un criterio
escrito y publicado (la coparticipación), uno no escrito pero reconstruible (el agua), y uno
retenido (las fojas 12 a 18 del plan hidráulico de \$35.100.000, que el decreto 5466-E aprueba
sin publicar). \textbf{La tercera posición es la que el resto del libro documenta para los
ochenta años siguientes}, y encontrar las otras dos conviviendo con ella en el mismo bienio
obliga a decir que la opacidad de criterio \textbf{no es una condición del Estado provincial
sino una elección que se toma acto por acto}.
\end{cautela}

\begin{cautela}
\textbf{Y los nombres de 1916 vuelven en 1948.}

\textbf{Hermann Pfister}, el perito que deslindó Loma Bola, Santa Mónica, Quitilipe y el
Potrero Gallinato en 1916 y 1917, pide en 1948 el agua del río para La Helvecia.
\textbf{Daniel Linares}, a quien el deslinde de 1916 nombra como propietario de Getsemaní,
pide el agua para El Angosto. \textbf{Los dos peticionantes de agua del bienio salen de los
deslindes de treinta años antes.} Y un \textbf{Urquiza} vuelve a ser autoridad de mesa
electoral en Vaqueros, veinticuatro años después de que Francisco S.\ Urquiza acumulara allí
cinco roles, en el circuito donde el mismo decreto registra el paraje \textbf{Villa Urquiza}.

\textbf{Coincidencia de apellido, con la regla de siempre}, y por eso este libro no afirma
identidad ni parentesco. \textbf{Lo que afirma es más simple y más verificable}: que el
conjunto de apellidos que aparece midiendo y poseyendo la tierra en 1916 es el mismo que
aparece pidiendo su agua y presidiendo sus mesas en 1948. \textbf{El que mide, el que posee y
el que administra no son tres conjuntos distintos en este departamento}, y eso el capítulo lo
venía diciendo para 1909 y 1924. \textbf{1948 muestra que todavía vale cuando el archivo se
cierra.}
\end{cautela}


\section{Y la tierra cambiaba de manos}

En julio de 1922, el mismo martillero que tres años después remataría la fundición anuncia
otra venta: «por disposición del señor Juez de Primera Instancia Dr. Mendioroz y como
correspondiente a la ejecución seguida por el Dr. Marcos Alsina contra don José Nogales,
venderé con base de \$800 la finca denominada Potrero del Gallinato, ubicada en el
Departamento de La Caldera y de propiedad del ejecutado».

\textbf{Y hay un acto anterior, que explica quién era el ejecutado.} En enero de 1917, Camilo
Cáceres pide el deslinde, la mensura y la \textbf{división de condominio} de la finca «\textit{Potrero
Gallinato}» ---así, sin «de»---, «\textit{que fue antes de Inocencio Aramayo}», con el perito Herman
Pfister\index{Pfister, Hermann} y por auto del juez de feria Augusto F.\ Torino. Entre sus linderos, al
este, \textbf{José Nogales} (B.O.\ Nº 636, h.\ 4, y Nº 637, h.\ 4). \textbf{Es el mismo José Nogales
que cinco años después aparece como el ejecutado cuya finca se remata}: en 1917 lindaba con el Potrero
Gallinato y en 1922 era su dueño. Entre las dos fechas, la finca estuvo en condominio y se dividió.

Con ese remate, la finca Potrero del Gallinato aparece en \textbf{seis} momentos de su historia, y
los seis son actos del Estado sobre el dominio. En 1909 es un partido policial con
comisario auxiliar propio. En 1917 se deslinda y se divide un condominio, y José Nogales es lindero.
En 1922 se anuncia su remate judicial por una deuda ---para el 28 de julio, seis días después del final del tramo leído, de modo que \textbf{el archivo relevado no dice si se vendió ni a quién}---.
En 2015 la Provincia declara operada a su favor la prescripción adquisitiva de una fracción
de tres mil doscientos metros cuadrados, para regularizar una escuela que estaba allí desde 1968. Y
en el mismo 2015 se aprueba sobre ella el club de campo Chacras del Gallinato, de mil
quinientas sesenta y nueve hectáreas.

Y hay un sexto acto. En octubre de 1930, el mismo martillero anuncia otro remate por el
juez de Comercio, en los autos «Embargo Preventivo Melchor Rocha vs.\ José González Yufra»:
una casa en la calle 25 de Mayo de la capital y, en el mismo acto, «una \textbf{estancia}
denominada Potrero de Gallinato en el departamento de la Caldera».

Ocho años después del primer remate, la finca vuelve a la subasta —ahora con el nombre de
estancia— y otra vez junto a un inmueble urbano de la ciudad de Salta, en la ejecución de un
mismo deudor.

Partido en 1909, deslinde y división de condominio en 1917, remate en 1922, remate en 1930,
prescripción en 2015, club de campo en 2015: \textbf{seis} actos del Estado sobre el dominio de la
misma tierra en ciento seis años, bajo cinco regímenes distintos. Este libro no afirma que los seis actos recaigan sobre
idéntica superficie —una finca de 1922 pudo subdividirse muchas veces—, pero el remate es
un eslabón identificado, con juez, ejecutante, ejecutado y martillero. Es por donde se
empieza a tirar del hilo.


